\documentclass[12pt,a4paper,openany]{book}
\usepackage[utf8]{inputenc}
\usepackage[T1]{fontenc}
\usepackage{lmodern}
\usepackage[margin=2.5cm]{geometry}
\usepackage{amsmath,amssymb,amsthm,mathrsfs,mathtools}
\usepackage{tikz}
\usetikzlibrary{arrows.meta,calc,positioning,patterns,decorations.markings,decorations.pathmorphing}
\usepackage{pgfplots}
\pgfplotsset{compat=1.18}
\usepackage[most]{tcolorbox}
\tcbuselibrary{theorems,skins,breakable}
\usepackage[colorlinks=true,linkcolor=blue!60!black,citecolor=green!50!black]{hyperref}
\usepackage{makeidx}
\makeindex
\usepackage{enumitem}
\usepackage{fancyhdr}
\usepackage{caption}
\usepackage{booktabs}
\usepackage{array}
\usepackage{longtable}
\usepackage{microtype}
\usepackage{xcolor}

\pagestyle{fancy}
\fancyhf{}
\fancyhead[LE]{\slshape\leftmark}
\fancyhead[RO]{\slshape\rightmark}
\fancyfoot[C]{\thepage}

% Theorem environments — \newtheorem (amsthm), NOT \newtcbtheorem
\theoremstyle{definition}
\newtheorem{definition}{Definition}[chapter]
\newtheorem{example}[definition]{Example}
\newtheorem{exercise}{Exercise}[chapter]

\theoremstyle{plain}
\newtheorem{theorem}[definition]{Theorem}
\newtheorem{proposition}[definition]{Proposition}
\newtheorem{lemma}[definition]{Lemma}
\newtheorem{corollary}[definition]{Corollary}

\theoremstyle{remark}
\newtheorem{remark}[definition]{Remark}
\newtheorem{notation}[definition]{Notation}

% tcolorbox wrapping
\tcolorboxenvironment{definition}{enhanced,breakable,colback=blue!5,colframe=blue!70!black,fonttitle=\bfseries,before skip=10pt,after skip=10pt,left=8pt,right=8pt}
\tcolorboxenvironment{theorem}{enhanced,breakable,colback=red!5,colframe=red!70!black,fonttitle=\bfseries,before skip=10pt,after skip=10pt,left=8pt,right=8pt}
\tcolorboxenvironment{proposition}{enhanced,breakable,colback=orange!5,colframe=orange!70!black,fonttitle=\bfseries,before skip=10pt,after skip=10pt,left=8pt,right=8pt}
\tcolorboxenvironment{lemma}{enhanced,breakable,colback=yellow!5,colframe=yellow!50!black,fonttitle=\bfseries,before skip=10pt,after skip=10pt,left=8pt,right=8pt}
\tcolorboxenvironment{corollary}{enhanced,breakable,colback=green!5,colframe=green!50!black,fonttitle=\bfseries,before skip=10pt,after skip=10pt,left=8pt,right=8pt}
\tcolorboxenvironment{remark}{enhanced,breakable,colback=gray!5,colframe=gray!50!black,before skip=8pt,after skip=8pt,left=6pt,right=6pt}
\tcolorboxenvironment{example}{enhanced,breakable,colback=purple!5,colframe=purple!50!black,before skip=8pt,after skip=8pt,left=6pt,right=6pt}

% Custom commands
\newcommand{\N}{\mathbb{N}}
\newcommand{\Z}{\mathbb{Z}}
\newcommand{\Q}{\mathbb{Q}}
\newcommand{\R}{\mathbb{R}}
\newcommand{\C}{\mathbb{C}}
\newcommand{\abs}[1]{\left|#1\right|}
\newcommand{\set}[1]{\left\{#1\right\}}
\newcommand{\powerset}{\mathcal{P}}
\DeclareMathOperator{\card}{card}

% ============================================================
\begin{document}

\begin{titlepage}
\begin{tikzpicture}[remember picture, overlay]
  \fill[left color=blue!15!white, right color=blue!5!white]
    (current page.south west) rectangle (current page.north east);
  \fill[color=orange!70!yellow]
    ([yshift=-2cm]current page.north west) rectangle ([yshift=-2.3cm]current page.north east);
  \fill[color=blue!80!black, rounded corners=8pt]
    ([xshift=3cm, yshift=-4cm]current page.north west) rectangle
    ([xshift=-3cm, yshift=-10cm]current page.north east);
  \node[anchor=center, text=white, font=\Huge\bfseries]
    at ([yshift=-6cm]current page.north) {Discrete Mathematics};
  \node[anchor=center, text=orange!80!yellow, font=\Large]
    at ([yshift=-7.5cm]current page.north) {Lecture Notes};
  \node[anchor=center, text=white!80, font=\large]
    at ([yshift=-8.8cm]current page.north) {Licence L2 --- 2025--2026};
  \node[anchor=center, text=blue!80!black, font=\Large\itshape]
    at ([yshift=-12cm]current page.north) {Ya\'e Ulrich Gaba};
  \draw[orange!70!yellow, line width=2pt]
    ([xshift=5cm, yshift=-13.5cm]current page.north west) --
    ([xshift=-5cm, yshift=-13.5cm]current page.north east);
  \node[anchor=center, text width=12cm, align=center, text=blue!60!black, font=\itshape]
    at ([yshift=-16cm]current page.north) {``God made the integers; all else is the work of man.''\\[4pt]--- Leopold Kronecker};
  \node[anchor=south, text=gray, font=\small]
    at ([yshift=1.5cm]current page.south) {\today};
\end{tikzpicture}
\end{titlepage}
% ============================================================
% PREFACE
% ============================================================
\chapter*{Preface}
\addcontentsline{toc}{chapter}{Preface}
\label{ch:preface}

Discrete mathematics provides the theoretical backbone of computer science, combinatorics,
cryptography, and many other fields of modern mathematics. Unlike continuous mathematics,
which studies objects that vary smoothly, discrete mathematics is concerned with structures
that are fundamentally separate and distinct: integers, graphs, logical statements, and
algorithms.

These notes offer a rigorous yet accessible treatment of the core topics in discrete
mathematics. Each chapter develops the theory from first principles, states and proves the
main results carefully, and provides numerous worked examples and exercises.

\medskip
\noindent\textbf{Prerequisites.} A basic familiarity with high-school algebra and an
appetite for logical reasoning are the only prerequisites. No prior exposure to formal
proof-writing is assumed; Chapter~2 is devoted entirely to proof techniques.

\medskip
\noindent\textbf{Structure.} The course is organized as follows:
\begin{itemize}[nosep]
  \item Chapters 1--2 lay the logical and proof-theoretic foundations.
  \item Chapters 3--5 develop set theory, relations, functions, and combinatorics.
  \item Chapters 6--8 treat number theory, graph theory, and recurrences.
  \item Chapters 9--10 address Boolean algebra, automata, and further topics.
\end{itemize}

\noindent Throughout, the reader will find:
\begin{itemize}[nosep]
  \item Boxed definitions and theorems for easy reference.
  \item Detailed proofs with commentary on the strategy employed.
  \item TikZ diagrams illustrating key ideas.
  \item End-of-chapter exercises ranging from routine to challenging.
\end{itemize}

\medskip
\noindent\textbf{How to read these notes.} Active engagement is essential. Attempt each
example before reading the solution, and try the exercises before consulting any hints.
Mathematical maturity is built through struggle, not spectation.

% ============================================================
% NOTATION
% ============================================================
\chapter*{Notation and Conventions}
\addcontentsline{toc}{chapter}{Notation and Conventions}
\label{ch:notation}

We collect here the principal symbols and conventions used throughout these notes.

\bigskip
\begin{center}
\begin{tabular}{>{\raggedright}p{3.2cm} p{9cm}}
\toprule
\textbf{Symbol} & \textbf{Meaning} \\
\midrule
$\N$ & The set of natural numbers $\set{0,1,2,\dots}$ \\
$\Z$ & The set of integers \\
$\Q$ & The set of rational numbers \\
$\R$ & The set of real numbers \\
$\C$ & The set of complex numbers \\
$\powerset(A)$ & The power set of $A$ \\
$\abs{A}$ or $\card(A)$ & The cardinality of the set $A$ \\
$\emptyset$ & The empty set \\
$\subset$, $\subseteq$ & Strict and non-strict set inclusion \\
$\cap$, $\cup$, $\setminus$ & Intersection, union, set difference \\
$A \times B$ & Cartesian product of $A$ and $B$ \\
$\lnot$, $\land$, $\lor$ & Logical NOT, AND, OR \\
$\Rightarrow$, $\Leftrightarrow$ & Implication and biconditional \\
$\forall$, $\exists$ & Universal and existential quantifiers \\
$n!$ & Factorial of $n$ \\
$\binom{n}{k}$ & Binomial coefficient ``$n$ choose $k$'' \\
$\lfloor x \rfloor$, $\lceil x \rceil$ & Floor and ceiling of $x$ \\
$a \mid b$ & $a$ divides $b$ \\
$a \equiv b \pmod{n}$ & $a$ is congruent to $b$ modulo $n$ \\
$\square$ or \qedsymbol & End of proof \\
\bottomrule
\end{tabular}
\end{center}

\bigskip
\noindent\textbf{Conventions.}
\begin{enumerate}[nosep]
  \item We include $0$ in the natural numbers: $\N = \set{0,1,2,\dots}$.
  \item The notation $\set{x \in S : P(x)}$ denotes the subset of $S$ satisfying
        property~$P$.
  \item All logarithms without an explicit base are taken to be base~$2$
        (as is common in discrete mathematics and computer science).
  \item The symbol $:=$ means ``is defined to be.''
\end{enumerate}

% ============================================================
\tableofcontents

% ============================================================
% CHAPTER 1: PROPOSITIONAL AND PREDICATE LOGIC
% ============================================================
\chapter{Propositional and Predicate Logic}
\label{ch:logic}
\index{logic}

Logic is the study of correct reasoning. In this chapter we develop the two main
layers of formal logic used throughout mathematics: \emph{propositional logic}, which
deals with statements and their truth values, and \emph{predicate logic}, which extends
the framework to handle variables, properties, and quantification.

% ------------------------------------------------------------
\section{Propositions and Logical Connectives}
\label{sec:propositions}
\index{proposition}

\begin{definition}[Proposition]
\label{def:proposition}
A \textbf{proposition}\index{proposition} is a declarative sentence that is either
\emph{true}~(T) or \emph{false}~(F), but not both.
\end{definition}

\begin{example}
\label{ex:proposition-examples}
The following are propositions:
\begin{enumerate}[nosep]
  \item ``$2 + 3 = 5$'' \quad (true).
  \item ``Every even number greater than $2$ is the sum of two primes''
        \quad (Goldbach's conjecture; either true or false, though currently unproven).
  \item ``$\sqrt{2}$ is irrational'' \quad (true).
\end{enumerate}
The following are \emph{not} propositions:
\begin{enumerate}[nosep]
  \item ``What time is it?'' \quad (a question).
  \item ``$x + 1 = 3$'' \quad (truth value depends on $x$; this is a predicate, not a proposition).
  \item ``Do your homework!'' \quad (a command).
\end{enumerate}
\end{example}

We denote propositions by lowercase letters $p, q, r, \dots$ and build compound
propositions using \emph{logical connectives}\index{connective}.

\begin{definition}[Logical connectives]
\label{def:connectives}
Let $p$ and $q$ be propositions. We define:
\begin{enumerate}[nosep]
  \item \textbf{Negation}\index{negation}: $\lnot p$ is true when $p$ is false,
        and false when $p$ is true.
  \item \textbf{Conjunction}\index{conjunction} (AND): $p \land q$ is true when both
        $p$ and $q$ are true.
  \item \textbf{Disjunction}\index{disjunction} (OR): $p \lor q$ is true when at
        least one of $p$, $q$ is true.
  \item \textbf{Implication}\index{implication}: $p \Rightarrow q$ is false only when
        $p$ is true and $q$ is false.
  \item \textbf{Biconditional}\index{biconditional}: $p \Leftrightarrow q$ is true
        when $p$ and $q$ have the same truth value.
  \item \textbf{Exclusive or}\index{exclusive or}: $p \oplus q$ is true when exactly
        one of $p$, $q$ is true.
\end{enumerate}
\end{definition}

\begin{remark}
\label{rem:implication-vacuous}
The implication $p \Rightarrow q$ is \emph{vacuously true}\index{vacuous truth} whenever
$p$ is false. This convention may seem strange at first but is essential for mathematical
reasoning. For instance, the statement ``If $n$ is a prime less than $2$, then $n$ is
even'' is true, because there is no prime less than~$2$.
\end{remark}

% ------------------------------------------------------------
\section{Truth Tables}
\label{sec:truth-tables}
\index{truth table}

A \emph{truth table} lists the truth value of a compound proposition for every possible
combination of truth values of its components.

\begin{example}[Truth table for basic connectives]
\label{ex:truth-table-basic}
\[
\begin{array}{cc|cccccc}
\toprule
p & q & \lnot p & p \land q & p \lor q & p \Rightarrow q & p \Leftrightarrow q & p \oplus q \\
\midrule
T & T & F & T & T & T & T & F \\
T & F & F & F & T & F & F & T \\
F & T & T & F & T & T & F & T \\
F & F & T & F & F & T & T & F \\
\bottomrule
\end{array}
\]
\end{example}

\begin{example}[Verifying a compound proposition]
\label{ex:truth-table-compound}
Consider the proposition $(p \land q) \Rightarrow p$. We construct its truth table:
\[
\begin{array}{cc|c|c}
\toprule
p & q & p \land q & (p \land q) \Rightarrow p \\
\midrule
T & T & T & T \\
T & F & F & T \\
F & T & F & T \\
F & F & F & T \\
\bottomrule
\end{array}
\]
Since the final column is always~T, the proposition is a tautology.
\end{example}

% ------------------------------------------------------------
\section{Tautologies, Contradictions, and Logical Equivalence}
\label{sec:tautologies}
\index{tautology}\index{contradiction}\index{logical equivalence}

\begin{definition}[Tautology, contradiction, contingency]
\label{def:tautology}
A compound proposition is called:
\begin{enumerate}[nosep]
  \item a \textbf{tautology}\index{tautology} if it is true for every assignment of
        truth values to its variables;
  \item a \textbf{contradiction}\index{contradiction} if it is false for every
        assignment;
  \item a \textbf{contingency} if it is neither a tautology nor a contradiction.
\end{enumerate}
\end{definition}

\begin{definition}[Logical equivalence]
\label{def:logical-equivalence}
Two propositions $\alpha$ and $\beta$ are \textbf{logically equivalent}, written
$\alpha \equiv \beta$, if $\alpha \Leftrightarrow \beta$ is a tautology.
\end{definition}

\begin{theorem}[Key logical equivalences]
\label{thm:logical-equivalences}
Let $p$, $q$, $r$ be propositions. Then:
\begin{enumerate}[nosep]
  \item \textbf{Double negation}\index{double negation}: $\lnot(\lnot p) \equiv p$.
  \item \textbf{Commutativity}\index{commutativity}:
        $p \land q \equiv q \land p$;\quad $p \lor q \equiv q \lor p$.
  \item \textbf{Associativity}\index{associativity}:
        $(p \land q) \land r \equiv p \land (q \land r)$;\quad
        $(p \lor q) \lor r \equiv p \lor (q \lor r)$.
  \item \textbf{Distributivity}\index{distributivity}:
        $p \land (q \lor r) \equiv (p \land q) \lor (p \land r)$;\quad
        $p \lor (q \land r) \equiv (p \lor q) \land (p \lor r)$.
  \item \textbf{Identity laws}: $p \land T \equiv p$;\quad $p \lor F \equiv p$.
  \item \textbf{Domination laws}: $p \lor T \equiv T$;\quad $p \land F \equiv F$.
  \item \textbf{Idempotent laws}: $p \land p \equiv p$;\quad $p \lor p \equiv p$.
  \item \textbf{Absorption laws}\index{absorption}: $p \lor (p \land q) \equiv p$;\quad
        $p \land (p \lor q) \equiv p$.
  \item \textbf{Implication equivalence}\index{implication!equivalence}:
        $p \Rightarrow q \equiv \lnot p \lor q$.
  \item \textbf{Contrapositive}\index{contrapositive}:
        $p \Rightarrow q \equiv \lnot q \Rightarrow \lnot p$.
\end{enumerate}
\end{theorem}

\begin{proof}
Each equivalence can be verified by constructing the appropriate truth table and
checking that both sides yield the same column. We illustrate (9).

\medskip
\noindent\emph{Proof of} (9): Consider the truth table:
\[
\begin{array}{cc|c|c|c}
\toprule
p & q & p \Rightarrow q & \lnot p & \lnot p \lor q \\
\midrule
T & T & T & F & T \\
T & F & F & F & F \\
F & T & T & T & T \\
F & F & T & T & T \\
\bottomrule
\end{array}
\]
Columns~3 and~5 are identical, so $p \Rightarrow q \equiv \lnot p \lor q$.
The remaining equivalences are verified analogously and left as exercises
(Exercise~\ref{exer:verify-equivalences}).
\end{proof}

% ------------------------------------------------------------
\section{De Morgan's Laws}
\label{sec:demorgan}
\index{De Morgan's laws}

De Morgan's laws are among the most frequently used equivalences in mathematics and
computer science. They describe how negation distributes over conjunction and disjunction.

\begin{theorem}[De Morgan's laws for propositions]
\label{thm:demorgan-prop}
For any propositions $p$ and $q$:
\begin{enumerate}[nosep]
  \item $\lnot(p \land q) \equiv \lnot p \lor \lnot q$.
  \item $\lnot(p \lor q) \equiv \lnot p \land \lnot q$.
\end{enumerate}
\end{theorem}

\begin{proof}
We prove both by truth table.

\medskip
\noindent\emph{Part~(1):}
\[
\begin{array}{cc|c|c|cc|c}
\toprule
p & q & p \land q & \lnot(p \land q) & \lnot p & \lnot q & \lnot p \lor \lnot q \\
\midrule
T & T & T & F & F & F & F \\
T & F & F & T & F & T & T \\
F & T & F & T & T & F & T \\
F & F & F & T & T & T & T \\
\bottomrule
\end{array}
\]
Columns~4 and~7 coincide, establishing the equivalence.

\medskip
\noindent\emph{Part~(2):} An analogous table (left to the reader) confirms
$\lnot(p \lor q) \equiv \lnot p \land \lnot q$.
\end{proof}

\begin{remark}
\label{rem:demorgan-generalized}
De Morgan's laws generalize to any finite number of propositions:
\[
  \lnot(p_1 \land p_2 \land \cdots \land p_n)
  \;\equiv\;
  \lnot p_1 \lor \lnot p_2 \lor \cdots \lor \lnot p_n,
\]
and similarly for $\lor$. This follows by induction on $n$
(see Exercise~\ref{exer:demorgan-generalized}).
\end{remark}

% ------------------------------------------------------------
\section{Predicate Logic and Quantifiers}
\label{sec:predicates}
\index{predicate logic}\index{quantifier}

Propositional logic cannot express statements like ``every integer greater than~$1$
has a prime factor.'' For this, we need \emph{predicate logic}.

\begin{definition}[Predicate]
\label{def:predicate}
A \textbf{predicate}\index{predicate} (or \textbf{propositional function}) is an
expression $P(x)$ involving a variable $x$ that becomes a proposition when $x$ is
assigned a specific value from a \textbf{domain of discourse}\index{domain of discourse}~$D$.
\end{definition}

\begin{example}
\label{ex:predicate}
Let $P(x)$: ``$x^2 - 1 = 0$'' with domain $D = \Z$.
Then $P(1)$ is true, $P(-1)$ is true, $P(2)$ is false, and $P(0)$ is false.
\end{example}

\begin{definition}[Universal and existential quantifiers]
\label{def:quantifiers}
Let $P(x)$ be a predicate with domain~$D$.
\begin{enumerate}[nosep]
  \item The \textbf{universal quantification}\index{universal quantifier}
        $\forall x \in D,\; P(x)$ asserts that $P(x)$ is true for
        \emph{every} $x$ in~$D$.
  \item The \textbf{existential quantification}\index{existential quantifier}
        $\exists x \in D,\; P(x)$ asserts that $P(x)$ is true for
        \emph{at least one} $x$ in~$D$.
\end{enumerate}
\end{definition}

\begin{example}
\label{ex:quantifier-statements}
Let $D = \Z$.
\begin{enumerate}[nosep]
  \item $\forall x \in \Z,\; x + 0 = x$ \quad (true).
  \item $\exists x \in \Z,\; x^2 = 2$ \quad (false, since $\sqrt{2} \notin \Z$).
  \item $\forall n \in \N,\; n(n+1) \text{ is even}$ \quad (true).
\end{enumerate}
\end{example}

% ------------------------------------------------------------
\section{Negation of Quantified Statements}
\label{sec:negation-quantifiers}
\index{negation!of quantifiers}

One of the most important skills in mathematics is correctly negating quantified
statements.

\begin{theorem}[De Morgan's laws for quantifiers]
\label{thm:demorgan-quantifiers}
\begin{enumerate}[nosep]
  \item $\lnot\bigl(\forall x,\; P(x)\bigr) \;\equiv\; \exists x,\; \lnot P(x)$.
  \item $\lnot\bigl(\exists x,\; P(x)\bigr) \;\equiv\; \forall x,\; \lnot P(x)$.
\end{enumerate}
\end{theorem}

\begin{proof}
(1) The statement $\forall x,\; P(x)$ is false exactly when there is some element
$x_0$ in the domain for which $P(x_0)$ is false, i.e., $\lnot P(x_0)$ is true.
This is precisely the assertion $\exists x,\; \lnot P(x)$.

\medskip\noindent
(2) Similarly, $\exists x,\; P(x)$ is false exactly when \emph{no} element
satisfies~$P$, i.e., every element satisfies $\lnot P$. This gives
$\forall x,\; \lnot P(x)$.
\end{proof}

\begin{example}[Negating a nested quantifier statement]
\label{ex:negate-nested}
Negate the statement: ``For every $\varepsilon > 0$, there exists $\delta > 0$ such
that $\abs{f(x) - L} < \varepsilon$ whenever $\abs{x - a} < \delta$.''

\medskip\noindent
\textbf{Solution.} The statement has the form
$\forall \varepsilon > 0,\; \exists \delta > 0,\; \forall x,\;
(\abs{x-a} < \delta \Rightarrow \abs{f(x)-L} < \varepsilon)$.

Negating step by step:
\begin{align*}
&\lnot\bigl[\forall \varepsilon > 0,\; \exists \delta > 0,\; \forall x,\;
(\abs{x-a} < \delta \Rightarrow \abs{f(x)-L} < \varepsilon)\bigr] \\
&\equiv \exists \varepsilon > 0,\; \lnot\bigl[\exists \delta > 0,\; \forall x,\;
(\abs{x-a} < \delta \Rightarrow \abs{f(x)-L} < \varepsilon)\bigr] \\
&\equiv \exists \varepsilon > 0,\; \forall \delta > 0,\;
\lnot\bigl[\forall x,\; (\abs{x-a} < \delta \Rightarrow \abs{f(x)-L} < \varepsilon)\bigr] \\
&\equiv \exists \varepsilon > 0,\; \forall \delta > 0,\; \exists x,\;
\lnot(\abs{x-a} < \delta \Rightarrow \abs{f(x)-L} < \varepsilon) \\
&\equiv \exists \varepsilon > 0,\; \forall \delta > 0,\; \exists x,\;
\bigl(\abs{x-a} < \delta \;\land\; \abs{f(x)-L} \geq \varepsilon\bigr).
\end{align*}
In words: ``There exists $\varepsilon > 0$ such that for every $\delta > 0$ there
is an $x$ with $\abs{x-a} < \delta$ and $\abs{f(x)-L} \geq \varepsilon$.''
\end{example}

% ------------------------------------------------------------
\section{Logic Diagrams}
\label{sec:logic-diagrams}
\index{logic diagram}

We can visualize the relationship between connectives and their truth conditions
using TikZ diagrams.

\begin{figure}[ht]
\centering
\begin{tikzpicture}[
  gate/.style={draw, minimum width=1.8cm, minimum height=1cm, font=\small\bfseries},
  input/.style={circle, draw, inner sep=2pt, fill=blue!10},
  output/.style={circle, draw, inner sep=2pt, fill=red!10},
  >=Stealth,
  node distance=1.2cm and 1.5cm
]
% AND gate
\node[input] (p1) {$p$};
\node[input, below=0.5cm of p1] (q1) {$q$};
\node[gate, right=of $(p1)!0.5!(q1)$] (and) {AND};
\node[output, right=of and] (o1) {$p \land q$};
\draw[->] (p1) -- (and.west |- p1);
\draw[->] (q1) -- (and.west |- q1);
\draw[->] (and) -- (o1);

% OR gate
\node[input, below=2cm of p1] (p2) {$p$};
\node[input, below=0.5cm of p2] (q2) {$q$};
\node[gate, right=of $(p2)!0.5!(q2)$] (or) {OR};
\node[output, right=of or] (o2) {$p \lor q$};
\draw[->] (p2) -- (or.west |- p2);
\draw[->] (q2) -- (or.west |- q2);
\draw[->] (or) -- (o2);

% NOT gate
\node[input, below=2cm of p2] (p3) {$p$};
\node[gate, right=1.5cm of p3] (not) {NOT};
\node[output, right=of not] (o3) {$\lnot p$};
\draw[->] (p3) -- (not);
\draw[->] (not) -- (o3);

% IMPLIES diagram
\node[input, below=2cm of p3] (p4) {$p$};
\node[input, below=0.5cm of p4] (q4) {$q$};
\node[gate, right=of $(p4)!0.5!(q4)$] (imp) {$\Rightarrow$};
\node[output, right=of imp] (o4) {$p \Rightarrow q$};
\draw[->] (p4) -- (imp.west |- p4);
\draw[->] (q4) -- (imp.west |- q4);
\draw[->] (imp) -- (o4);
\end{tikzpicture}
\caption{Schematic representation of basic logic gates.}
\label{fig:logic-gates}
\end{figure}

\begin{figure}[ht]
\centering
\begin{tikzpicture}[
  box/.style={draw, rounded corners, minimum width=2.2cm, minimum height=0.8cm,
              font=\small, fill=#1},
  arr/.style={->, >=Stealth, thick},
  node distance=0.8cm and 2.8cm
]
% Implication chain
\node[box=blue!15] (pq) {$p \Rightarrow q$};
\node[box=green!15, right=of pq] (contra) {$\lnot q \Rightarrow \lnot p$};
\node[box=orange!15, below=of pq] (npq) {$\lnot p \lor q$};
\node[box=red!15, below=of contra] (nqnp) {$\lnot(p \land \lnot q)$};

\draw[arr, blue!60!black] (pq) -- node[above, font=\footnotesize]{contrapositive} (contra);
\draw[arr, green!50!black] (pq) -- node[left, font=\footnotesize]{equiv.} (npq);
\draw[arr, orange!60!black] (contra) -- node[right, font=\footnotesize]{equiv.} (nqnp);
\draw[arr, red!60!black] (npq) -- node[below, font=\footnotesize]{De Morgan} (nqnp);
\end{tikzpicture}
\caption{Equivalent forms of the implication $p \Rightarrow q$.}
\label{fig:implication-forms}
\end{figure}

% ------------------------------------------------------------
\section{Exercises}
\label{sec:exercises-ch1}

\begin{exercise}
\label{exer:classify-propositions}
Classify each of the following as a proposition or not. If it is a proposition,
determine its truth value.
\begin{enumerate}[label=(\alph*),nosep]
  \item $3 + 7 = 10$.
  \item $x^2 \geq 0$.
  \item The moon is made of cheese.
  \item $\sqrt{16} = 4$.
  \item Close the door!
  \item Every prime number is odd.
\end{enumerate}
\end{exercise}

\begin{exercise}
\label{exer:truth-tables}
Construct truth tables for each of the following compound propositions:
\begin{enumerate}[label=(\alph*),nosep]
  \item $(p \Rightarrow q) \land (q \Rightarrow p)$.
  \item $p \Rightarrow (q \Rightarrow p)$.
  \item $(p \lor q) \land (\lnot p \lor r) \Rightarrow (q \lor r)$.
\end{enumerate}
\end{exercise}

\begin{exercise}
\label{exer:verify-equivalences}
Verify each of the equivalences in Theorem~\ref{thm:logical-equivalences}
by constructing truth tables.
\end{exercise}

\begin{exercise}
\label{exer:demorgan-generalized}
Use mathematical induction to prove the generalized De Morgan's law:
\[
  \lnot(p_1 \land p_2 \land \cdots \land p_n)
  \;\equiv\;
  \lnot p_1 \lor \lnot p_2 \lor \cdots \lor \lnot p_n
\]
for all $n \geq 2$.
\end{exercise}

\begin{exercise}
\label{exer:negate-statements}
Negate each of the following statements, writing the result in positive form
(i.e., push the negation inward as far as possible):
\begin{enumerate}[label=(\alph*),nosep]
  \item $\forall x \in \R,\; x^2 + 1 > 0$.
  \item $\exists n \in \N,\; \forall m \in \N,\; n \leq m$.
  \item $\forall \varepsilon > 0,\; \exists N \in \N,\; \forall n \geq N,\;
        \abs{a_n - L} < \varepsilon$.
  \item $\forall x \in \R,\; (x > 0 \Rightarrow \exists y \in \R,\; y^2 = x)$.
\end{enumerate}
\end{exercise}

\begin{exercise}
\label{exer:logical-arguments}
Determine whether each of the following arguments is valid. Justify your answer
using truth tables or known equivalences.
\begin{enumerate}[label=(\alph*),nosep]
  \item $p \Rightarrow q$, $q \Rightarrow r$, therefore $p \Rightarrow r$
        (hypothetical syllogism).
  \item $p \Rightarrow q$, $\lnot q$, therefore $\lnot p$
        (modus tollens).
  \item $p \lor q$, $\lnot p$, therefore $q$
        (disjunctive syllogism).
  \item $p \Rightarrow q$, $q$, therefore $p$
        (affirming the consequent).
\end{enumerate}
\end{exercise}

\begin{exercise}
\label{exer:knights-knaves}
On an island, every inhabitant is either a \emph{knight} (always tells the truth)
or a \emph{knave} (always lies). You meet two inhabitants, $A$ and $B$.
\begin{enumerate}[label=(\alph*),nosep]
  \item $A$ says: ``Both of us are knaves.'' What are $A$ and $B$?
  \item $A$ says: ``I am a knave or $B$ is a knight.'' What are $A$ and $B$?
\end{enumerate}
\emph{Hint:} Translate each scenario into propositional logic and analyze.
\end{exercise}

% ------------------------------------------------------------
\section{Chapter Summary}
\label{sec:summary-ch1}

\begin{itemize}[nosep]
  \item A \textbf{proposition} is a declarative sentence with a definite truth value.
  \item The main \textbf{logical connectives} are $\lnot$, $\land$, $\lor$,
        $\Rightarrow$, $\Leftrightarrow$, and $\oplus$.
  \item A \textbf{tautology} is always true; a \textbf{contradiction} is always false.
  \item Two propositions are \textbf{logically equivalent} if they have the same truth
        table.
  \item \textbf{De Morgan's laws}: $\lnot(p \land q) \equiv \lnot p \lor \lnot q$
        and $\lnot(p \lor q) \equiv \lnot p \land \lnot q$.
  \item A \textbf{predicate} $P(x)$ becomes a proposition when $x$ is instantiated.
  \item The \textbf{universal quantifier} $\forall$ and \textbf{existential quantifier}
        $\exists$ bind variables in predicates.
  \item Negation swaps quantifiers:
        $\lnot(\forall x,\; P(x)) \equiv \exists x,\; \lnot P(x)$.
\end{itemize}

% ============================================================
% CHAPTER 2: PROOF TECHNIQUES
% ============================================================
\chapter{Proof Techniques}
\label{ch:proofs}
\index{proof techniques}

A \emph{proof} is a rigorous logical argument establishing the truth of a mathematical
statement beyond any doubt. In this chapter we develop the main proof strategies:
direct proof, proof by contradiction, proof by contrapositive, mathematical induction,
and several additional techniques. Mastery of these methods is essential for all that
follows.

% ------------------------------------------------------------
\section{Direct Proof}
\label{sec:direct-proof}
\index{direct proof}

The most straightforward proof technique: to prove $p \Rightarrow q$, assume $p$ is
true and deduce, through a chain of logical steps, that $q$ must also be true.

\begin{definition}[Even and odd integers]
\label{def:even-odd}
An integer $n$ is \textbf{even}\index{even integer} if $n = 2k$ for some integer~$k$,
and \textbf{odd}\index{odd integer} if $n = 2k + 1$ for some integer~$k$.
\end{definition}

\begin{theorem}
\label{thm:sum-even}
The sum of two even integers is even.
\end{theorem}

\begin{proof}
Let $m$ and $n$ be even integers. By definition, there exist integers $a$ and $b$
such that $m = 2a$ and $n = 2b$. Then
\[
  m + n = 2a + 2b = 2(a + b).
\]
Since $a + b \in \Z$, the integer $m + n$ has the form $2k$ with $k = a + b$,
so $m + n$ is even.
\end{proof}

\begin{theorem}
\label{thm:product-even-odd}
The product of an even integer and an odd integer is even.
\end{theorem}

\begin{proof}
Let $m$ be even and $n$ be odd. Write $m = 2a$ and $n = 2b + 1$ for integers $a, b$.
Then
\[
  mn = 2a(2b + 1) = 2\bigl(a(2b+1)\bigr).
\]
Since $a(2b+1) \in \Z$, the product $mn$ is even.
\end{proof}

\begin{example}[Direct proof: divisibility]
\label{ex:direct-divisibility}
\textbf{Claim.} If $a \mid b$ and $b \mid c$, then $a \mid c$.\index{divisibility!transitivity}

\begin{proof}
Suppose $a \mid b$ and $b \mid c$. Then there exist integers $k$ and $\ell$ such that
$b = ak$ and $c = b\ell$. Therefore,
\[
  c = b\ell = (ak)\ell = a(k\ell).
\]
Since $k\ell \in \Z$, we conclude $a \mid c$.
\end{proof}
\end{example}

\begin{example}[Direct proof: sum of consecutive integers]
\label{ex:direct-consecutive}
\textbf{Claim.} For every positive integer $n$, the sum $1 + 2 + \cdots + n = \frac{n(n+1)}{2}$.

\begin{proof}
Let $S = 1 + 2 + \cdots + n$. Writing the sum in reverse order:
\[
  S = n + (n-1) + \cdots + 1.
\]
Adding term by term:
\[
  2S = (1+n) + (2+n-1) + \cdots + (n+1) = n(n+1).
\]
Hence $S = \frac{n(n+1)}{2}$.
\end{proof}
\end{example}

% ------------------------------------------------------------
\section{Proof by Contradiction}
\label{sec:contradiction-proof}
\index{proof by contradiction}

To prove a statement~$P$, we assume $\lnot P$ and derive a logical contradiction
(a statement of the form $Q \land \lnot Q$). Since a contradiction is always false,
the assumption $\lnot P$ must be false, so $P$ is true.

\begin{theorem}
\label{thm:sqrt2-irrational}
$\sqrt{2}$ is irrational.\index{irrationality of $\sqrt{2}$}
\end{theorem}

\begin{proof}
Suppose, for the sake of contradiction, that $\sqrt{2}$ is rational. Then we can
write $\sqrt{2} = \frac{a}{b}$ where $a, b \in \Z$, $b \neq 0$, and
$\gcd(a, b) = 1$ (the fraction is in lowest terms).

Squaring both sides gives $2 = \frac{a^2}{b^2}$, so $a^2 = 2b^2$. This means $a^2$
is even, so $a$ must be even (if $a$ were odd, $a^2$ would be odd). Write $a = 2c$
for some integer~$c$. Then
\[
  (2c)^2 = 2b^2 \implies 4c^2 = 2b^2 \implies b^2 = 2c^2.
\]
So $b^2$ is even, hence $b$ is even. But then both $a$ and $b$ are even, contradicting
$\gcd(a,b) = 1$.

Therefore, $\sqrt{2}$ is irrational.
\end{proof}

\begin{theorem}
\label{thm:infinitely-many-primes}
There are infinitely many prime numbers.\index{primes!infinitude}
\end{theorem}

\begin{proof}
Suppose, for contradiction, that there are only finitely many primes, say
$p_1, p_2, \dots, p_n$. Consider the integer
\[
  N = p_1 p_2 \cdots p_n + 1.
\]
Since $N > 1$, it has a prime factor~$p$. This prime $p$ must be one of
$p_1, \dots, p_n$ (since these are all the primes). But $N \equiv 1 \pmod{p_i}$
for each~$i$, so $p_i \nmid N$ for any~$i$. This contradicts the fact that $p$
divides~$N$.

Therefore, there are infinitely many primes.
\end{proof}

\begin{example}[Contradiction: no smallest positive rational]
\label{ex:no-smallest-rational}
\textbf{Claim.} There is no smallest positive rational number.

\begin{proof}
Suppose, for contradiction, that $r$ is the smallest positive rational number.
Consider $\frac{r}{2}$. Since $r > 0$, we have $\frac{r}{2} > 0$, and since $r$
is rational, $\frac{r}{2}$ is also rational. But $\frac{r}{2} < r$, contradicting
the assumption that $r$ is the smallest positive rational. Therefore, no such $r$
exists.
\end{proof}
\end{example}

% ------------------------------------------------------------
\section{Proof by Contrapositive}
\label{sec:contrapositive-proof}
\index{proof by contrapositive}

To prove $p \Rightarrow q$, we can instead prove the logically equivalent statement
$\lnot q \Rightarrow \lnot p$ (the contrapositive; see
Theorem~\ref{thm:logical-equivalences}(10)). This is especially useful when the
direct approach seems difficult but the negation of the conclusion provides a
convenient starting point.

\begin{theorem}
\label{thm:n2-even-then-n-even}
If $n^2$ is even, then $n$ is even.
\end{theorem}

\begin{proof}
We prove the contrapositive: if $n$ is odd, then $n^2$ is odd.

Suppose $n$ is odd. Then $n = 2k + 1$ for some integer~$k$. Hence,
\[
  n^2 = (2k+1)^2 = 4k^2 + 4k + 1 = 2(2k^2 + 2k) + 1.
\]
Since $2k^2 + 2k \in \Z$, the integer $n^2$ is odd. By the contrapositive, if $n^2$
is even, then $n$ is even.
\end{proof}

\begin{theorem}
\label{thm:contrapositive-3nminus1}
For every integer $n$, if $3n - 1$ is even, then $n$ is odd.
\end{theorem}

\begin{proof}
We prove the contrapositive: if $n$ is even, then $3n - 1$ is odd.

Suppose $n$ is even, so $n = 2k$ for some $k \in \Z$. Then
\[
  3n - 1 = 3(2k) - 1 = 6k - 1 = 2(3k - 1) + 1.
\]
Since $3k - 1 \in \Z$, the integer $3n - 1$ is odd, as required.
\end{proof}

\begin{example}[Contrapositive: rational square root]
\label{ex:contrapositive-sqrt}
\textbf{Claim.} Let $n$ be a positive integer. If $\sqrt{n}$ is rational, then
$\sqrt{n}$ is an integer.

\begin{proof}
We prove the contrapositive: if $\sqrt{n}$ is not an integer, then $\sqrt{n}$ is
irrational.

Suppose $\sqrt{n}$ is not an integer, but assume for contradiction that $\sqrt{n}$
is rational. Write $\sqrt{n} = \frac{a}{b}$ with $\gcd(a,b) = 1$ and $b \geq 2$
(since $\sqrt{n}$ is not an integer). Then $n = \frac{a^2}{b^2}$, so $a^2 = nb^2$.
Let $p$ be a prime factor of $b$. Then $p^2 \mid b^2 \mid a^2$, so $p \mid a$,
contradicting $\gcd(a,b) = 1$.

Hence, if $\sqrt{n}$ is not an integer, it is irrational. Equivalently, if $\sqrt{n}$
is rational, then $\sqrt{n}$ is an integer.
\end{proof}
\end{example}

% ------------------------------------------------------------
\section{Proof by Mathematical Induction}
\label{sec:induction}
\index{induction!mathematical}

Mathematical induction is a fundamental proof technique for statements involving
natural numbers. It exploits the well-ordered structure of~$\N$.

\begin{theorem}[Principle of Mathematical Induction (PMI)]
\label{thm:pmi}
Let $P(n)$ be a predicate defined for integers $n \geq n_0$. If:
\begin{enumerate}[nosep]
  \item \textbf{Base case}\index{base case}: $P(n_0)$ is true, and
  \item \textbf{Inductive step}\index{inductive step}: for every integer $k \geq n_0$,
        $P(k) \Rightarrow P(k+1)$,
\end{enumerate}
then $P(n)$ is true for all integers $n \geq n_0$.
\end{theorem}

\begin{remark}
\label{rem:induction-analogy}
Induction is often compared to a row of dominoes: the base case knocks over the first
domino, and the inductive step ensures that each falling domino knocks over the next.
\end{remark}

\begin{figure}[ht]
\centering
\begin{tikzpicture}[domino/.style={draw, fill=blue!15, minimum width=0.4cm,
  minimum height=1.2cm, rounded corners=1pt}]
  \foreach \i in {0,...,6} {
    \node[domino, rotate={-\i*8}] at (\i*1.1, {-\i*\i*0.03}) {};
    \node[below, font=\tiny] at (\i*1.1, -0.8) {$P(n_0\!+\!\i)$};
  }
  \node at (8.5, 0) {$\cdots$};
  \draw[->, thick, red, decorate, decoration={snake, amplitude=2pt, segment length=6pt}]
    (-0.6, 0.9) -- (0.0, 0.5);
  \node[red, font=\footnotesize, above] at (-0.3, 0.9) {Base case};
\end{tikzpicture}
\caption{The domino analogy for mathematical induction.}
\label{fig:induction-dominoes}
\end{figure}

\begin{theorem}
\label{thm:sum-squares}
For every positive integer $n$,
\[
  \sum_{k=1}^{n} k^2 = \frac{n(n+1)(2n+1)}{6}.
\]
\end{theorem}

\begin{proof}
We proceed by induction on $n$.

\medskip\noindent
\textbf{Base case} ($n = 1$): The left side is $1^2 = 1$ and the right side is
$\frac{1 \cdot 2 \cdot 3}{6} = 1$. They agree.

\medskip\noindent
\textbf{Inductive step.} Suppose the formula holds for some $k \geq 1$, i.e.,
\[
  \sum_{i=1}^{k} i^2 = \frac{k(k+1)(2k+1)}{6}.
  \tag{Inductive Hypothesis}
\]
We must show it holds for $k + 1$:
\begin{align*}
  \sum_{i=1}^{k+1} i^2
  &= \left(\sum_{i=1}^{k} i^2\right) + (k+1)^2 \\
  &= \frac{k(k+1)(2k+1)}{6} + (k+1)^2 \\
  &= \frac{k(k+1)(2k+1) + 6(k+1)^2}{6} \\
  &= \frac{(k+1)\bigl[k(2k+1) + 6(k+1)\bigr]}{6} \\
  &= \frac{(k+1)(2k^2 + 7k + 6)}{6} \\
  &= \frac{(k+1)(k+2)(2k+3)}{6} \\
  &= \frac{(k+1)\bigl((k+1)+1\bigr)\bigl(2(k+1)+1\bigr)}{6}.
\end{align*}
This is the formula with $n = k + 1$. By induction, the result holds for all
$n \geq 1$.
\end{proof}

\begin{theorem}
\label{thm:bernoulli}
\textbf{(Bernoulli's inequality)}\index{Bernoulli's inequality}
For every real number $x \geq -1$ and every positive integer $n$,
\[
  (1 + x)^n \geq 1 + nx.
\]
\end{theorem}

\begin{proof}
We use induction on $n$.

\medskip\noindent
\textbf{Base case} ($n = 1$): $(1+x)^1 = 1 + x = 1 + 1 \cdot x$. The inequality
holds with equality.

\medskip\noindent
\textbf{Inductive step.} Assume $(1+x)^k \geq 1 + kx$ for some $k \geq 1$.
Since $1 + x \geq 0$ (because $x \geq -1$), we may multiply both sides by $1 + x$:
\begin{align*}
  (1+x)^{k+1} &= (1+x)^k \cdot (1+x) \\
               &\geq (1 + kx)(1 + x) \\
               &= 1 + kx + x + kx^2 \\
               &= 1 + (k+1)x + kx^2 \\
               &\geq 1 + (k+1)x,
\end{align*}
where the last inequality uses $kx^2 \geq 0$. By induction, the result holds for
all $n \geq 1$.
\end{proof}

\begin{example}[Induction: divisibility]
\label{ex:induction-divisibility}
\textbf{Claim.} For every non-negative integer $n$, $3 \mid (n^3 - n)$.\index{divisibility}

\begin{proof}
\textbf{Base case} ($n = 0$): $0^3 - 0 = 0$, and $3 \mid 0$.

\medskip\noindent
\textbf{Inductive step.} Suppose $3 \mid (k^3 - k)$ for some $k \geq 0$. Then
\begin{align*}
  (k+1)^3 - (k+1)
  &= k^3 + 3k^2 + 3k + 1 - k - 1 \\
  &= (k^3 - k) + 3k^2 + 3k \\
  &= (k^3 - k) + 3(k^2 + k).
\end{align*}
By the inductive hypothesis, $3 \mid (k^3 - k)$, and clearly $3 \mid 3(k^2 + k)$.
Therefore $3 \mid \bigl[(k^3 - k) + 3(k^2 + k)\bigr] = (k+1)^3 - (k+1)$.

By induction, $3 \mid (n^3 - n)$ for all $n \geq 0$.
\end{proof}
\end{example}

% ------------------------------------------------------------
\section{Strong Induction}
\label{sec:strong-induction}
\index{induction!strong}

In \emph{strong induction} (also called \emph{complete induction}), the inductive
hypothesis assumes $P(n_0), P(n_0 + 1), \dots, P(k)$ are all true (not just $P(k)$)
in order to prove $P(k+1)$.

\begin{theorem}[Principle of Strong Induction]
\label{thm:strong-induction}
Let $P(n)$ be a predicate defined for integers $n \geq n_0$. If:
\begin{enumerate}[nosep]
  \item \textbf{Base case}: $P(n_0)$ is true, and
  \item \textbf{Inductive step}: for every $k \geq n_0$,
        $\bigl[P(n_0) \land P(n_0+1) \land \cdots \land P(k)\bigr] \Rightarrow P(k+1)$,
\end{enumerate}
then $P(n)$ is true for all $n \geq n_0$.
\end{theorem}

\begin{remark}
\label{rem:strong-vs-weak}
Strong induction is equivalent in power to ordinary (weak) induction; any proof by
strong induction can be converted to one using weak induction (and vice versa).
However, strong induction is often more natural when the proof of $P(k+1)$ requires
information about values \emph{earlier} than $P(k)$.
\end{remark}

\begin{theorem}[Fundamental Theorem of Arithmetic --- existence part]
\label{thm:fta-existence}
Every integer $n \geq 2$ can be written as a product of primes.\index{fundamental theorem of arithmetic}
\end{theorem}

\begin{proof}
We use strong induction on $n$.

\medskip\noindent
\textbf{Base case} ($n = 2$): The integer $2$ is itself prime, so it is trivially
a product of primes (a product with one factor).

\medskip\noindent
\textbf{Inductive step.} Let $k \geq 2$ and assume that every integer
$m$ with $2 \leq m \leq k$ can be written as a product of primes.
Consider $k + 1$.

\emph{Case~1:} $k + 1$ is prime. Then $k+1$ is a product of primes.

\emph{Case~2:} $k + 1$ is composite. Then $k + 1 = ab$ for some integers $a, b$
with $2 \leq a, b \leq k$. By the strong inductive hypothesis, both $a$ and $b$
can be written as products of primes. Hence $k + 1 = ab$ is also a product of primes.

By strong induction, every integer $n \geq 2$ is a product of primes.
\end{proof}

\begin{example}[Strong induction: postage stamps]
\label{ex:stamps}
\textbf{Claim.} Every integer amount of postage $n \geq 12$ cents can be formed
using only $4$-cent and $5$-cent stamps.

\begin{proof}
\textbf{Base cases:}
\begin{itemize}[nosep]
  \item $n = 12$: $12 = 3 \times 4$.
  \item $n = 13$: $13 = 2 \times 4 + 1 \times 5$.
  \item $n = 14$: $14 = 1 \times 4 + 2 \times 5$.
  \item $n = 15$: $15 = 3 \times 5$.
\end{itemize}

\textbf{Inductive step.} Let $k \geq 15$ and assume the claim holds for all integers
$m$ with $12 \leq m \leq k$. We prove it for $k + 1$.

Since $k + 1 \geq 16$, we have $k + 1 - 4 = k - 3 \geq 12$. By the strong inductive
hypothesis, $k - 3$ can be formed using $4$-cent and $5$-cent stamps. Adding one
more $4$-cent stamp gives $k + 1$ cents.

By strong induction, the claim holds for all $n \geq 12$.
\end{proof}
\end{example}

% ------------------------------------------------------------
\section{The Pigeonhole Principle}
\label{sec:pigeonhole}
\index{pigeonhole principle}

The pigeonhole principle is a deceptively simple counting argument with surprisingly
powerful applications.

\begin{theorem}[Pigeonhole Principle]
\label{thm:pigeonhole}
If $n + 1$ or more objects are placed into $n$ containers, then at least one container
holds two or more objects.\index{pigeonhole principle}
\end{theorem}

\begin{proof}
We prove the contrapositive. Suppose every container holds at most one object.
Then the total number of objects is at most $n \cdot 1 = n$. Hence, if the number
of objects exceeds~$n$, at least one container must hold at least two objects.
\end{proof}

\begin{theorem}[Generalized Pigeonhole Principle]
\label{thm:pigeonhole-generalized}
If $N$ objects are placed into $k$ containers, then at least one container holds
at least $\lceil N/k \rceil$ objects.\index{pigeonhole principle!generalized}
\end{theorem}

\begin{proof}
Suppose, for contradiction, that every container holds at most
$\lceil N/k \rceil - 1$ objects. Then the total number of objects is at most
\[
  k \left(\left\lceil \frac{N}{k} \right\rceil - 1\right)
  < k \cdot \frac{N}{k}
  = N,
\]
where we used $\lceil N/k \rceil - 1 < N/k$. This contradicts the fact that there
are $N$ objects in total.
\end{proof}

\begin{example}[Pigeonhole: birthdays]
\label{ex:pigeonhole-birthdays}
In any group of $367$ people, at least two must share the same birthday
(since there are at most $366$ possible birthdays, including February~29).
\end{example}

\begin{example}[Pigeonhole: socks in the dark]
\label{ex:pigeonhole-socks}
A drawer contains $10$ black socks and $10$ white socks. How many socks must you
draw (in the dark) to guarantee a matching pair?

\textbf{Solution.} The ``containers'' are the $2$ colors. By the pigeonhole principle,
drawing $3$ socks guarantees at least $\lceil 3/2 \rceil = 2$ socks of the same color.
\end{example}

\begin{example}[Pigeonhole: subset sums]
\label{ex:pigeonhole-subsets}
\textbf{Claim.} Given any set of $n + 1$ integers from $\set{1, 2, \dots, 2n}$,
there must be two elements in the set such that one divides the other.

\begin{proof}
Write each element $a$ in the form $a = 2^{s} \cdot m$ where $m$ is odd. The odd
part $m$ must be one of the $n$ odd numbers $1, 3, 5, \dots, 2n - 1$ (the
``containers''). Since there are $n + 1$ elements and only $n$ possible odd parts,
by the pigeonhole principle two elements share the same odd part, say $a = 2^s m$
and $b = 2^t m$ with $s < t$. Then $a \mid b$.
\end{proof}
\end{example}

\begin{example}[Pigeonhole: lattice points]
\label{ex:pigeonhole-lattice}
\textbf{Claim.} Among any five points chosen from the integer lattice $\Z^2$,
some pair has a midpoint that is also a lattice point.

\begin{proof}
Each lattice point $(x, y)$ can be classified by the parities of $x$ and $y$:
$(even, even)$, $(even, odd)$, $(odd, even)$, or $(odd, odd)$. These give $4$
classes. Among $5$ points, by the pigeonhole principle, at least two belong to the
same class, say $(x_1, y_1)$ and $(x_2, y_2)$ with $x_1 \equiv x_2 \pmod{2}$ and
$y_1 \equiv y_2 \pmod{2}$. Their midpoint is
\[
  \left(\frac{x_1 + x_2}{2},\; \frac{y_1 + y_2}{2}\right),
\]
which has integer coordinates because $x_1 + x_2$ and $y_1 + y_2$ are both even.
\end{proof}
\end{example}

% ------------------------------------------------------------
\section{The Well-Ordering Principle}
\label{sec:well-ordering}
\index{well-ordering principle}

\begin{theorem}[Well-Ordering Principle]
\label{thm:well-ordering}
Every non-empty subset of $\N$ has a least element.\index{well-ordering principle}
\end{theorem}

\begin{remark}
\label{rem:well-ordering-equivalence}
The Well-Ordering Principle is logically equivalent to the Principle of Mathematical
Induction. We take the Well-Ordering Principle as an axiom (it follows from the Peano
axioms for~$\N$) and note that proofs using it have a distinctive flavour: one considers
the set of counterexamples, argues it is non-empty if the statement fails, then derives
a contradiction from the existence of a minimal counterexample.
\end{remark}

\begin{example}[Well-ordering proof: division algorithm]
\label{ex:well-ordering-division}
\textbf{Claim (Division Algorithm).}\index{division algorithm}
For every integer $a$ and positive integer $d$, there exist unique integers $q$
(quotient) and $r$ (remainder) such that $a = dq + r$ and $0 \leq r < d$.

\begin{proof}[Proof of existence]
Consider the set
\[
  S = \set{a - dk : k \in \Z,\; a - dk \geq 0}.
\]
This set is non-empty: if $a \geq 0$, take $k = 0$; if $a < 0$, take $k = a$
(then $a - da = a(1 - d) \geq 0$ since $d \geq 1$ and $a < 0$).

Since $S \subseteq \N$ and $S \neq \emptyset$, the Well-Ordering Principle
guarantees that $S$ has a least element $r = a - dq$ for some $q$.

We claim $r < d$. Suppose for contradiction that $r \geq d$. Then
\[
  a - d(q + 1) = r - d \geq 0,
\]
so $r - d \in S$. But $r - d < r$, contradicting the minimality of~$r$.
Therefore $0 \leq r < d$, as required.
\end{proof}
\end{example}

\begin{example}[Well-ordering proof: every positive integer $\geq 2$ has a prime factor]
\label{ex:well-ordering-prime-factor}
\textbf{Claim.} Every integer $n \geq 2$ has a prime factor.

\begin{proof}
Suppose the claim is false. Let $S = \set{n \in \Z : n \geq 2 \text{ and $n$ has no prime factor}}$.
By assumption, $S$ is non-empty. By the Well-Ordering Principle, $S$ has a least
element~$m$.

Since $m \in S$, the integer $m$ has no prime factor. In particular, $m$ itself is
not prime (since every prime is its own prime factor). So $m$ is composite:
$m = ab$ with $2 \leq a, b < m$. Since $a < m$ and $m$ is the least element of~$S$,
the integer $a$ has a prime factor~$p$. But $p \mid a$ and $a \mid m$ imply
$p \mid m$, contradicting $m \in S$.

Therefore $S = \emptyset$, and every integer $n \geq 2$ has a prime factor.
\end{proof}
\end{example}

% ------------------------------------------------------------
\section{Additional Examples and Techniques}
\label{sec:additional-proof-examples}

We gather several more worked examples that combine the techniques introduced above.

\begin{example}[Proof by cases]
\label{ex:proof-by-cases}
\index{proof by cases}
\textbf{Claim.} For every integer $n$, the integer $n^2 + n$ is even.

\begin{proof}
We consider two cases.

\emph{Case~1: $n$ is even.} Write $n = 2k$. Then $n^2 + n = 4k^2 + 2k = 2(2k^2 + k)$,
which is even.

\emph{Case~2: $n$ is odd.} Write $n = 2k + 1$. Then
$n^2 + n = (2k+1)^2 + (2k+1) = 4k^2 + 4k + 1 + 2k + 1 = 4k^2 + 6k + 2 = 2(2k^2 + 3k + 1)$,
which is even.

In both cases, $n^2 + n$ is even.
\end{proof}
\end{example}

\begin{example}[Existence proof]
\label{ex:existence-proof}
\index{existence proof}
\textbf{Claim.} There exist irrational numbers $a$ and $b$ such that $a^b$ is
rational.

\begin{proof}
Consider $\sqrt{2}^{\sqrt{2}}$. This number is either rational or irrational.

\emph{Case~1:} $\sqrt{2}^{\sqrt{2}}$ is rational. Then take $a = b = \sqrt{2}$
(both irrational), and $a^b$ is rational. Done.

\emph{Case~2:} $\sqrt{2}^{\sqrt{2}}$ is irrational. Then take
$a = \sqrt{2}^{\sqrt{2}}$ (irrational by assumption) and $b = \sqrt{2}$ (irrational).
We get
\[
  a^b = \left(\sqrt{2}^{\sqrt{2}}\right)^{\sqrt{2}}
      = \sqrt{2}^{\sqrt{2} \cdot \sqrt{2}}
      = \sqrt{2}^{2}
      = 2,
\]
which is rational.

In both cases, we have found irrational $a, b$ with $a^b$ rational.
\end{proof}
\end{example}

\begin{example}[Uniqueness proof]
\label{ex:uniqueness-proof}
\index{uniqueness proof}
\textbf{Claim.} The additive identity in $\Z$ is unique: if $e \in \Z$ satisfies
$n + e = n$ for all $n \in \Z$, then $e = 0$.

\begin{proof}
Suppose $e$ and $e'$ are both additive identities. Then $e = e + e' = e'$, where the
first equality uses that $e'$ is an identity and the second uses that $e$ is an identity.
\end{proof}
\end{example}

\begin{example}[Constructive proof: rational between any two reals]
\label{ex:density-rationals}
\textbf{Claim.} Between any two distinct real numbers there is a rational number.\index{density of rationals}

\begin{proof}
Let $a < b$ be real numbers. The Archimedean property\index{Archimedean property}
guarantees the existence of a positive integer $n$ such that $n(b - a) > 1$, i.e.,
$nb - na > 1$. Therefore, there exists an integer $m$ with $na < m < nb$
(since the interval $(na, nb)$ has length greater than~$1$). Dividing by~$n$:
\[
  a < \frac{m}{n} < b.
\]
The number $\frac{m}{n}$ is rational and lies strictly between $a$ and~$b$.
\end{proof}
\end{example}

% ------------------------------------------------------------
\section{Exercises}
\label{sec:exercises-ch2}

\begin{exercise}
\label{exer:direct-proof-exercises}
Prove each of the following by direct proof:
\begin{enumerate}[label=(\alph*),nosep]
  \item The sum of two odd integers is even.
  \item If $n$ is odd, then $n^2$ is odd.
  \item If $a \mid b$ and $a \mid c$, then $a \mid (b + c)$.
  \item For all $n \in \N$, $n^3 - n$ is divisible by~$6$.
\end{enumerate}
\end{exercise}

\begin{exercise}
\label{exer:contradiction-exercises}
Prove each of the following by contradiction:
\begin{enumerate}[label=(\alph*),nosep]
  \item $\sqrt{3}$ is irrational.
  \item If $n^2$ is odd, then $n$ is odd.
  \item There is no largest integer.
  \item The sum of a rational number and an irrational number is irrational.
\end{enumerate}
\end{exercise}

\begin{exercise}
\label{exer:contrapositive-exercises}
Prove each of the following by contrapositive:
\begin{enumerate}[label=(\alph*),nosep]
  \item If $n^2$ is divisible by~$3$, then $n$ is divisible by~$3$.
  \item If $x + y \geq 2$, then $x \geq 1$ or $y \geq 1$.
  \item If $n^3 + 5$ is odd, then $n$ is even.
\end{enumerate}
\end{exercise}

\begin{exercise}
\label{exer:induction-exercises}
Prove each of the following by mathematical induction:
\begin{enumerate}[label=(\alph*),nosep]
  \item $\displaystyle\sum_{k=1}^{n} k^3 = \left(\frac{n(n+1)}{2}\right)^{\!2}$ for all $n \geq 1$.
  \item $2^n > n^2$ for all $n \geq 5$.
  \item $n! > 2^n$ for all $n \geq 4$.
  \item $\displaystyle\sum_{k=0}^{n} 2^k = 2^{n+1} - 1$ for all $n \geq 0$.
\end{enumerate}
\end{exercise}

\begin{exercise}
\label{exer:strong-induction-exercises}
Prove each of the following by strong induction:
\begin{enumerate}[label=(\alph*),nosep]
  \item Every integer $n \geq 2$ is either prime or can be expressed as a product of
        two or more primes.
  \item Every positive integer can be represented as a sum of distinct powers of~$2$
        (binary representation).
  \item The Fibonacci sequence, defined by $F_1 = 1$, $F_2 = 1$, and
        $F_n = F_{n-1} + F_{n-2}$ for $n \geq 3$, satisfies $F_n < 2^n$ for all
        $n \geq 1$.
\end{enumerate}
\end{exercise}

\begin{exercise}
\label{exer:pigeonhole-exercises}
Solve each of the following using the pigeonhole principle:
\begin{enumerate}[label=(\alph*),nosep]
  \item Show that among any $n+1$ integers chosen from $\set{1, 2, \dots, 2n}$,
        there must be a pair of consecutive integers.
  \item Show that in any group of $6$ people, there are either $3$ mutual friends or
        $3$ mutual strangers. \emph{(Hint: consider the edges of $K_6$ colored with
        two colors.)}
  \item Prove that for every positive integer $n$, there is a multiple of $n$ whose
        decimal representation consists entirely of $0$s and $1$s.
        \emph{(Hint: consider the numbers $1, 11, 111, \dots$)}
  \item Show that among any $52$ integers, there exist two whose difference is
        divisible by~$51$.
\end{enumerate}
\end{exercise}

\begin{exercise}
\label{exer:well-ordering-exercises}
Use the Well-Ordering Principle to prove:
\begin{enumerate}[label=(\alph*),nosep]
  \item There is no integer between $0$ and~$1$.
  \item For every pair of positive integers $a$ and $b$ with $a \mid b$ and
        $b \mid a$, we have $a = b$.
\end{enumerate}
\end{exercise}

\begin{exercise}
\label{exer:mixed-proofs}
Prove or disprove each of the following (choose any appropriate method):
\begin{enumerate}[label=(\alph*),nosep]
  \item For every integer $n$, $n^2 + n + 41$ is prime.
  \item If $a^2 \mid b^2$, then $a \mid b$ (where $a, b$ are positive integers).
  \item For all positive reals $x$ and $y$, $\frac{x + y}{2} \geq \sqrt{xy}$
        (the AM--GM inequality\index{AM--GM inequality}).
  \item Every integer of the form $4k + 3$ has a prime factor of the form $4j + 3$.
\end{enumerate}
\end{exercise}

% ------------------------------------------------------------
\section{Chapter Summary}
\label{sec:summary-ch2}

\begin{itemize}[nosep]
  \item A \textbf{direct proof} of $p \Rightarrow q$ assumes $p$ and derives $q$.
  \item A \textbf{proof by contradiction} assumes $\lnot P$ and derives a
        contradiction, thereby establishing~$P$.
  \item A \textbf{proof by contrapositive} of $p \Rightarrow q$ proves the equivalent
        statement $\lnot q \Rightarrow \lnot p$.
  \item \textbf{Mathematical induction} (weak) proves $P(n)$ for all $n \geq n_0$ by
        establishing a base case and an implication $P(k) \Rightarrow P(k+1)$.
  \item \textbf{Strong induction} allows the inductive hypothesis to assume
        $P(n_0), \dots, P(k)$ when proving $P(k+1)$.
  \item The \textbf{Pigeonhole Principle}: if $n+1$ objects go into $n$ boxes, some
        box contains at least two objects. The generalized version gives a lower bound
        of $\lceil N/k \rceil$.
  \item The \textbf{Well-Ordering Principle}: every non-empty subset of $\N$ has a
        least element. It is equivalent to induction.
  \item Other proof strategies include \emph{proof by cases}, \emph{existence proofs},
        \emph{uniqueness proofs}, and \emph{constructive proofs}.
\end{itemize}

% === END OF CHUNK preamble+ch1-2 ===
% === START OF CHUNK ch3-4 ===

% ############################################################
% CHAPTER 3
% ############################################################
\chapter{Sets, Relations, and Functions}
\label{ch:sets-relations-functions}
\index{set}\index{relation}\index{function}

\begin{quote}
\itshape
``A set is a Many that allows itself to be thought of as a One.''\\
--- Georg Cantor
\end{quote}

\medskip
\noindent
Sets are the foundational language of modern mathematics. In this chapter we
develop the core notions of set theory, explore relations (with special
attention to equivalence relations and partial orders), and study functions
as a particular type of relation. We conclude with a discussion of
cardinality and Cantor's landmark diagonal argument.

% ============================================================
\section{Sets and subsets}
\label{sec:sets-subsets}
\index{set!definition}

\begin{definition}[Set]\label{def:set}
A \emph{set} is an unordered collection of distinct objects, called the
\emph{elements} (or \emph{members}) of the set. If $x$ is an element of a
set $A$ we write $x \in A$; otherwise $x \notin A$.
\end{definition}

\index{set!roster notation}\index{set!set-builder notation}
A set may be specified by \emph{roster notation},
$A = \set{1, 2, 3}$,
or by \emph{set-builder notation},
$A = \set{x \in \Z : x^2 < 10}$.

\begin{definition}[Subset]\label{def:subset}
\index{subset}
Let $A$ and $B$ be sets. We say $A$ is a \emph{subset} of $B$, written
$A \subseteq B$, if every element of $A$ belongs to $B$:
\[
  A \subseteq B \iff \bigl(\forall x,\; x \in A \implies x \in B\bigr).
\]
If $A \subseteq B$ and $A \neq B$, we write $A \subsetneq B$ and call $A$
a \emph{proper subset} of $B$.
\end{definition}

\begin{definition}[Set equality]\label{def:set-equality}
\index{set!equality}
Two sets $A$ and $B$ are \emph{equal}, written $A = B$, if and only if
$A \subseteq B$ and $B \subseteq A$.
\end{definition}

\begin{definition}[Empty set]\label{def:empty-set}
\index{empty set}
The \emph{empty set}, denoted $\varnothing$, is the unique set with no
elements. For every set $A$ we have $\varnothing \subseteq A$.
\end{definition}

% ============================================================
\section{Power set and Cartesian product}
\label{sec:powerset-cartesian}

\begin{definition}[Power set]\label{def:powerset}
\index{power set}
The \emph{power set} of a set $A$, denoted $\powerset(A)$, is the set of all
subsets of $A$:
\[
  \powerset(A) = \set{S : S \subseteq A}.
\]
\end{definition}

\begin{example}\label{ex:powerset}
If $A = \set{1,2,3}$, then
\[
  \powerset(A) = \bigl\{\varnothing,\,\set{1},\,\set{2},\,\set{3},\,
    \set{1,2},\,\set{1,3},\,\set{2,3},\,\set{1,2,3}\bigr\}.
\]
In particular $\card{\powerset(A)} = 2^3 = 8$.
\end{example}

\begin{proposition}\label{prop:powerset-card}
\index{power set!cardinality}
If $A$ is a finite set with $\card{A} = n$, then $\card{\powerset(A)} = 2^n$.
\end{proposition}

\begin{proof}
Each subset of $A$ is determined by, for each of the $n$ elements,
choosing whether to include it or not. There are $2$ choices per element
and the choices are independent, giving $2^n$ subsets in total.
\end{proof}

\begin{definition}[Cartesian product]\label{def:cartesian}
\index{Cartesian product}
The \emph{Cartesian product} of sets $A$ and $B$ is
\[
  A \times B = \set{(a, b) : a \in A,\; b \in B}.
\]
More generally, $A_1 \times A_2 \times \cdots \times A_n =
\set{(a_1, \ldots, a_n) : a_i \in A_i \text{ for each } i}$.
\end{definition}

% ============================================================
\section{Set operations}
\label{sec:set-operations}
\index{set!operations}

\begin{definition}[Union, intersection, difference, complement]
\label{def:set-ops}
Let $A$ and $B$ be subsets of a universal set $U$.
\begin{enumerate}[label=(\roman*)]
  \item \textbf{Union:}\index{union} $A \cup B = \set{x : x \in A \text{ or } x \in B}$.
  \item \textbf{Intersection:}\index{intersection} $A \cap B = \set{x : x \in A \text{ and } x \in B}$.
  \item \textbf{Difference:}\index{set difference} $A \setminus B = \set{x : x \in A \text{ and } x \notin B}$.
  \item \textbf{Complement:}\index{complement} $A^c = U \setminus A = \set{x \in U : x \notin A}$.
  \item \textbf{Symmetric difference:}\index{symmetric difference}
        $A \mathbin{\triangle} B = (A \setminus B) \cup (B \setminus A)$.
\end{enumerate}
\end{definition}

\begin{figure}[ht]
\centering
\begin{tikzpicture}[scale=0.8]
  % Union
  \begin{scope}[shift={(-5,0)}]
    \draw[thick] (-2.2,-1.8) rectangle (2.2,1.8);
    \fill[blue!20] (0.6,0) circle (1.1);
    \fill[blue!20] (-0.6,0) circle (1.1);
    \draw[thick] (-0.6,0) circle (1.1) node[left=8pt] {$A$};
    \draw[thick] (0.6,0) circle (1.1) node[right=8pt] {$B$};
    \node at (0,-2.2) {$A \cup B$};
  \end{scope}
  % Intersection
  \begin{scope}
    \begin{scope}
      \clip (-0.6,0) circle (1.1);
      \fill[green!30] (0.6,0) circle (1.1);
    \end{scope}
    \draw[thick] (-2.2,-1.8) rectangle (2.2,1.8);
    \draw[thick] (-0.6,0) circle (1.1) node[left=8pt] {$A$};
    \draw[thick] (0.6,0) circle (1.1) node[right=8pt] {$B$};
    \node at (0,-2.2) {$A \cap B$};
  \end{scope}
  % Difference A \ B
  \begin{scope}[shift={(5,0)}]
    \begin{scope}
      \clip (-0.6,0) circle (1.1);
      \fill[red!25] (-0.6,0) circle (1.1);
      \fill[white] (0.6,0) circle (1.1);
    \end{scope}
    \draw[thick] (-2.2,-1.8) rectangle (2.2,1.8);
    \draw[thick] (-0.6,0) circle (1.1) node[left=8pt] {$A$};
    \draw[thick] (0.6,0) circle (1.1) node[right=8pt] {$B$};
    \node at (0,-2.2) {$A \setminus B$};
  \end{scope}
\end{tikzpicture}
\caption{Venn diagrams for union, intersection, and set difference.}
\label{fig:venn-diagrams}
\end{figure}

% ---- De Morgan's Laws ----
\begin{theorem}[De Morgan's Laws]\label{thm:de-morgan}
\index{De Morgan's laws}
Let $A$ and $B$ be subsets of a universal set $U$. Then:
\begin{enumerate}[label=(\roman*)]
  \item $(A \cup B)^c = A^c \cap B^c$,
  \item $(A \cap B)^c = A^c \cup B^c$.
\end{enumerate}
\end{theorem}

\begin{proof}
We prove (i); the proof of (ii) is analogous.

$(\subseteq)$\; Let $x \in (A \cup B)^c$. Then $x \notin A \cup B$, which
means $x \notin A$ and $x \notin B$. Hence $x \in A^c$ and $x \in B^c$,
so $x \in A^c \cap B^c$.

$(\supseteq)$\; Let $x \in A^c \cap B^c$. Then $x \notin A$ and
$x \notin B$. Therefore $x \notin A \cup B$, i.e.\ $x \in (A \cup B)^c$.

Since each side is a subset of the other, equality holds.
\end{proof}

\begin{remark}\label{rem:de-morgan-general}
De Morgan's laws generalise to arbitrary families of sets:
$\bigl(\bigcup_{i \in I} A_i\bigr)^c = \bigcap_{i \in I} A_i^c$ and
$\bigl(\bigcap_{i \in I} A_i\bigr)^c = \bigcup_{i \in I} A_i^c$.
\end{remark}

% ============================================================
\section{Relations}
\label{sec:relations}
\index{relation}

\begin{definition}[Binary relation]\label{def:relation}
A \emph{binary relation} $R$ from a set $A$ to a set $B$ is a subset
$R \subseteq A \times B$. We often write $a \mathrel{R} b$ instead of
$(a,b) \in R$. When $B = A$, we say $R$ is a relation \emph{on} $A$.
\end{definition}

\begin{definition}[Properties of relations]\label{def:relation-props}
\index{relation!reflexive}\index{relation!symmetric}\index{relation!transitive}\index{relation!antisymmetric}
Let $R$ be a relation on a set $A$. We say $R$ is:
\begin{enumerate}[label=(\roman*)]
  \item \textbf{reflexive} if $a \mathrel{R} a$ for all $a \in A$;
  \item \textbf{symmetric} if $a \mathrel{R} b \implies b \mathrel{R} a$;
  \item \textbf{antisymmetric} if $a \mathrel{R} b$ and $b \mathrel{R} a$
        imply $a = b$;
  \item \textbf{transitive} if $a \mathrel{R} b$ and $b \mathrel{R} c$
        imply $a \mathrel{R} c$.
\end{enumerate}
\end{definition}

% ---- Equivalence relations ----
\subsection{Equivalence relations and partitions}
\label{subsec:equiv-rel}
\index{equivalence relation}

\begin{definition}[Equivalence relation]\label{def:equiv-rel}
A relation $\sim$ on a set $A$ is an \emph{equivalence relation} if it is
reflexive, symmetric, and transitive.
\end{definition}

\begin{definition}[Equivalence class]\label{def:equiv-class}
\index{equivalence class}
Let $\sim$ be an equivalence relation on $A$ and let $a \in A$. The
\emph{equivalence class} of $a$ is
\[
  [a] = \set{x \in A : x \sim a}.
\]
The set of all equivalence classes is called the \emph{quotient set} and
denoted $A / {\sim}$.
\end{definition}

\begin{definition}[Partition]\label{def:partition}
\index{partition}
A \emph{partition} of a set $A$ is a collection $\mathcal{P}$ of non-empty,
pairwise disjoint subsets of $A$ whose union is $A$:
\begin{enumerate}[label=(\roman*)]
  \item $P \neq \varnothing$ for every $P \in \mathcal{P}$;
  \item $P \cap Q = \varnothing$ for all $P, Q \in \mathcal{P}$ with $P \neq Q$;
  \item $\bigcup_{P \in \mathcal{P}} P = A$.
\end{enumerate}
\end{definition}

\begin{theorem}[Partition theorem]\label{thm:partition}
\index{partition theorem}
Let $A$ be a non-empty set.
\begin{enumerate}[label=(\alph*)]
  \item If $\sim$ is an equivalence relation on $A$, then the collection of
        equivalence classes $A / {\sim}$ is a partition of $A$.
  \item Conversely, if $\mathcal{P}$ is a partition of $A$, then the relation
        $\sim_{\mathcal{P}}$ defined by $a \sim_{\mathcal{P}} b$ if and only
        if $a$ and $b$ belong to the same block of $\mathcal{P}$ is an
        equivalence relation on $A$, and $A / {\sim_{\mathcal{P}}} = \mathcal{P}$.
\end{enumerate}
\end{theorem}

\begin{proof}
\textbf{(a)}\; We verify the three conditions of a partition.

\emph{(i)}\; For every $a \in A$, reflexivity gives $a \sim a$, so
$a \in [a]$ and each class is non-empty.

\emph{(iii)}\; Since every $a \in A$ belongs to $[a]$, we have
$A = \bigcup_{a \in A} [a]$.

\emph{(ii)}\; Suppose $[a] \cap [b] \neq \varnothing$; pick
$c \in [a] \cap [b]$. Then $c \sim a$ and $c \sim b$. By symmetry,
$a \sim c$, and by transitivity $a \sim b$. For any $x \in [a]$ we have
$x \sim a \sim b$, so $x \in [b]$; hence $[a] \subseteq [b]$. By a
symmetric argument $[b] \subseteq [a]$, giving $[a] = [b]$.

\medskip\noindent
\textbf{(b)}\; Define $a \sim_{\mathcal{P}} b$ iff there exists
$P \in \mathcal{P}$ with $a, b \in P$.

\emph{Reflexive:}\; Since $\mathcal{P}$ covers $A$, every $a$ belongs to
some block, so $a \sim_{\mathcal{P}} a$.

\emph{Symmetric:}\; If $a, b \in P$, then $b, a \in P$.

\emph{Transitive:}\; If $a, b \in P$ and $b, c \in Q$, then $b \in P \cap Q$.
Since blocks are pairwise disjoint, $P = Q$, so $a, c \in P$.

Finally, the class $[a]$ equals the unique block containing $a$, so
$A / {\sim_{\mathcal{P}}} = \mathcal{P}$.
\end{proof}

\begin{example}\label{ex:congruence-equiv}
\index{congruence}
Fix $n \geq 1$. The relation $\equiv \pmod{n}$ on $\Z$ is an equivalence
relation. The equivalence classes are the residue classes
$[0], [1], \ldots, [n-1]$, which partition $\Z$.
\end{example}

% ---- Partial orders ----
\subsection{Partial orders and Hasse diagrams}
\label{subsec:partial-orders}
\index{partial order}

\begin{definition}[Partial order]\label{def:partial-order}
A relation $\preceq$ on a set $A$ is a \emph{partial order} if it is
reflexive, antisymmetric, and transitive. The pair $(A, \preceq)$ is called a
\emph{partially ordered set} (or \emph{poset}).
\end{definition}

\begin{example}\label{ex:divisibility-poset}
\index{divisibility!partial order}
On $\N^{*}$, the divisibility relation $a \mid b$ defines a partial order.
\end{example}

\begin{definition}[Hasse diagram]\label{def:hasse}
\index{Hasse diagram}
A \emph{Hasse diagram} for a finite poset $(A, \preceq)$ is a graph in which:
\begin{enumerate}[label=(\roman*)]
  \item each element of $A$ is a vertex;
  \item if $a \prec b$ and there is no $c$ with $a \prec c \prec b$
        (i.e.\ $b$ \emph{covers} $a$), an edge is drawn from $a$ upward to $b$;
  \item edges implied by transitivity are omitted.
\end{enumerate}
\end{definition}

\begin{example}\label{ex:hasse-divides-12}
The Hasse diagram of the divisors of $12$, ordered by divisibility:
\end{example}

\begin{figure}[ht]
\centering
\begin{tikzpicture}[
  every node/.style={draw, circle, minimum size=7mm, inner sep=1pt, font=\small},
  edge/.style={thick},
  scale=1.1
]
  \node (1)  at (0,0)   {$1$};
  \node (2)  at (-1.5,1) {$2$};
  \node (3)  at (1.5,1) {$3$};
  \node (4)  at (-2.2,2) {$4$};
  \node (6)  at (0,2)   {$6$};
  \node (12) at (0,3)   {$12$};

  \draw[edge] (1)  -- (2);
  \draw[edge] (1)  -- (3);
  \draw[edge] (2)  -- (4);
  \draw[edge] (2)  -- (6);
  \draw[edge] (3)  -- (6);
  \draw[edge] (4)  -- (12);
  \draw[edge] (6)  -- (12);
\end{tikzpicture}
\caption{Hasse diagram for $(\set{1,2,3,4,6,12},\, \mid\,)$.}
\label{fig:hasse-12}
\end{figure}

\begin{definition}[Total order]\label{def:total-order}
\index{total order}
A partial order $\preceq$ on $A$ is a \emph{total} (or \emph{linear}) order
if for all $a, b \in A$, either $a \preceq b$ or $b \preceq a$. In this case
$(A, \preceq)$ is called a \emph{chain}.
\end{definition}

\begin{example}\label{ex:usual-order-total}
The usual $\leq$ on $\R$ is a total order. By contrast, divisibility on
$\N^{*}$ is not total: neither $2 \mid 3$ nor $3 \mid 2$.
\end{example}

% ============================================================
\section{Functions}
\label{sec:functions}
\index{function}

\begin{definition}[Function]\label{def:function}
A \emph{function} (or \emph{map}) $f\colon A \to B$ is a relation
$f \subseteq A \times B$ such that for every $a \in A$ there exists
exactly one $b \in B$ with $(a,b) \in f$. We write $b = f(a)$.
The set $A$ is the \emph{domain}, and $B$ is the \emph{codomain}.
The \emph{image} (or \emph{range}) of $f$ is
$\operatorname{im}(f) = \set{f(a) : a \in A}$.
\end{definition}

\begin{definition}[Injection, surjection, bijection]\label{def:inj-surj-bij}
\index{injection}\index{surjection}\index{bijection}
Let $f\colon A \to B$.
\begin{enumerate}[label=(\roman*)]
  \item $f$ is \emph{injective} (one-to-one) if
        $f(a_1) = f(a_2) \implies a_1 = a_2$.
  \item $f$ is \emph{surjective} (onto) if for every $b \in B$ there exists
        $a \in A$ with $f(a) = b$.
  \item $f$ is \emph{bijective} if it is both injective and surjective.
\end{enumerate}
\end{definition}

\begin{figure}[ht]
\centering
\begin{tikzpicture}[scale=0.85, every node/.style={font=\small}]
  % Injection
  \begin{scope}[shift={(-5,0)}]
    \draw[thick, rounded corners] (-0.8,-1.8) rectangle (0.8,1.8);
    \draw[thick, rounded corners] (2.2,-1.8) rectangle (3.8,1.8);
    \node[above] at (0,1.8) {$A$};
    \node[above] at (3,1.8) {$B$};
    \foreach \y/\lbl in {1/a, 0/b, -1/c} {
      \node[fill=blue!20, circle, inner sep=2pt] (L\lbl) at (0,\y) {$\lbl$};
    }
    \foreach \y/\lbl in {1.2/1, 0.4/2, -0.4/3, -1.2/4} {
      \node[fill=red!20, circle, inner sep=2pt] (R\lbl) at (3,\y) {$\lbl$};
    }
    \draw[->, thick] (La) -- (R1);
    \draw[->, thick] (Lb) -- (R3);
    \draw[->, thick] (Lc) -- (R4);
    \node at (1.5,-2.4) {Injective};
  \end{scope}
  % Surjection
  \begin{scope}
    \draw[thick, rounded corners] (-0.8,-1.8) rectangle (0.8,1.8);
    \draw[thick, rounded corners] (2.2,-1.8) rectangle (3.8,1.8);
    \node[above] at (0,1.8) {$A$};
    \node[above] at (3,1.8) {$B$};
    \foreach \y/\lbl in {1.2/a, 0.4/b, -0.4/c, -1.2/d} {
      \node[fill=blue!20, circle, inner sep=2pt] (L\lbl) at (0,\y) {$\lbl$};
    }
    \foreach \y/\lbl in {1/1, 0/2, -1/3} {
      \node[fill=red!20, circle, inner sep=2pt] (R\lbl) at (3,\y) {$\lbl$};
    }
    \draw[->, thick] (La) -- (R1);
    \draw[->, thick] (Lb) -- (R1);
    \draw[->, thick] (Lc) -- (R2);
    \draw[->, thick] (Ld) -- (R3);
    \node at (1.5,-2.4) {Surjective};
  \end{scope}
  % Bijection
  \begin{scope}[shift={(5,0)}]
    \draw[thick, rounded corners] (-0.8,-1.8) rectangle (0.8,1.8);
    \draw[thick, rounded corners] (2.2,-1.8) rectangle (3.8,1.8);
    \node[above] at (0,1.8) {$A$};
    \node[above] at (3,1.8) {$B$};
    \foreach \y/\lbl in {1/a, 0/b, -1/c} {
      \node[fill=blue!20, circle, inner sep=2pt] (L\lbl) at (0,\y) {$\lbl$};
    }
    \foreach \y/\lbl in {1/1, 0/2, -1/3} {
      \node[fill=red!20, circle, inner sep=2pt] (R\lbl) at (3,\y) {$\lbl$};
    }
    \draw[->, thick] (La) -- (R1);
    \draw[->, thick] (Lb) -- (R2);
    \draw[->, thick] (Lc) -- (R3);
    \node at (1.5,-2.4) {Bijective};
  \end{scope}
\end{tikzpicture}
\caption{Diagrams illustrating injection, surjection, and bijection.}
\label{fig:inj-surj-bij}
\end{figure}

\begin{definition}[Composition]\label{def:composition}
\index{composition}
Given $f\colon A \to B$ and $g\colon B \to C$, the \emph{composition}
$g \circ f \colon A \to C$ is defined by $(g \circ f)(a) = g(f(a))$.
\end{definition}

\begin{proposition}\label{prop:composition-preserves}
\begin{enumerate}[label=(\roman*)]
  \item If $f$ and $g$ are injective, then $g \circ f$ is injective.
  \item If $f$ and $g$ are surjective, then $g \circ f$ is surjective.
  \item If $f$ and $g$ are bijective, then $g \circ f$ is bijective.
\end{enumerate}
\end{proposition}

\begin{proof}
(i)\; Suppose $(g \circ f)(a_1) = (g \circ f)(a_2)$. Then
$g(f(a_1)) = g(f(a_2))$. Since $g$ is injective, $f(a_1) = f(a_2)$.
Since $f$ is injective, $a_1 = a_2$.

(ii)\; Let $c \in C$. Since $g$ is surjective, there exists $b \in B$ with
$g(b) = c$. Since $f$ is surjective, there exists $a \in A$ with $f(a) = b$.
Then $(g \circ f)(a) = g(b) = c$.

(iii)\; This follows immediately from (i) and (ii).
\end{proof}

\begin{definition}[Inverse function]\label{def:inverse-function}
\index{inverse function}
Let $f\colon A \to B$ be a bijection. The \emph{inverse} of $f$ is the
function $f^{-1}\colon B \to A$ defined by $f^{-1}(b) = a$ where $a$ is
the unique element of $A$ satisfying $f(a) = b$.
\end{definition}

\begin{proposition}\label{prop:inverse-bijection}
A function $f\colon A \to B$ is bijective if and only if there exists a
function $g\colon B \to A$ such that $g \circ f = \mathrm{id}_A$ and
$f \circ g = \mathrm{id}_B$.
\end{proposition}

\begin{proof}
$(\Rightarrow)$\; Take $g = f^{-1}$. For all $a \in A$,
$(f^{-1} \circ f)(a) = f^{-1}(f(a)) = a$, and for all $b \in B$,
$(f \circ f^{-1})(b) = f(f^{-1}(b)) = b$.

$(\Leftarrow)$\; Suppose such $g$ exists. If $f(a_1) = f(a_2)$, apply $g$:
$a_1 = g(f(a_1)) = g(f(a_2)) = a_2$, so $f$ is injective. Given $b \in B$,
set $a = g(b)$; then $f(a) = f(g(b)) = b$, so $f$ is surjective.
\end{proof}

% ============================================================
\section{Cardinality}
\label{sec:cardinality}
\index{cardinality}

\begin{definition}[Equinumerous sets]\label{def:equinumerous}
\index{equinumerous}
Two sets $A$ and $B$ are \emph{equinumerous}, written $\card{A} = \card{B}$,
if there exists a bijection $f\colon A \to B$.
\end{definition}

\begin{definition}[Countable set]\label{def:countable}
\index{countable set}
A set $A$ is \emph{countably infinite} if $\card{A} = \card{\N}$, i.e.\
there exists a bijection $A \to \N$. A set is \emph{countable} if it is
finite or countably infinite. A set that is not countable is
\emph{uncountable}.
\end{definition}

\begin{proposition}\label{prop:Z-countable}
$\Z$ is countable.
\end{proposition}

\begin{proof}
The function $f\colon \N \to \Z$ defined by
\[
  f(n) = \begin{cases}
    n/2 & \text{if } n \text{ is even},\\
    -(n+1)/2 & \text{if } n \text{ is odd},
  \end{cases}
\]
is a bijection: the sequence of values is $0, -1, 1, -2, 2, -3, 3, \ldots$
\end{proof}

\begin{proposition}\label{prop:Q-countable}
$\Q$ is countable.
\end{proposition}

\begin{proof}[Proof sketch]
List the positive rationals by diagonalising the grid of fractions $p/q$
(with $p, q \geq 1$), skipping those already encountered. This produces
an enumeration of $\Q^{+}$; interleave with the negatives to enumerate
all of $\Q$.
\end{proof}

\begin{theorem}[Cantor's diagonal argument]\label{thm:cantor-diagonal}
\index{Cantor's diagonal argument}\index{uncountable}
The set $\R$ is uncountable. More precisely, the interval $(0,1)$ is
uncountable.
\end{theorem}

\begin{proof}
Suppose for contradiction that $(0,1)$ is countable. Then we can list its
elements as $r_1, r_2, r_3, \ldots$, where each $r_n$ has a decimal
expansion
\[
  r_n = 0.d_{n1}\, d_{n2}\, d_{n3}\, \cdots
\]
(choosing the non-terminating expansion when ambiguous).

Define a real number $r^{*} = 0.d_1^{*}\, d_2^{*}\, d_3^{*}\, \cdots$ by
\[
  d_k^{*} = \begin{cases}
    3 & \text{if } d_{kk} \neq 3, \\
    7 & \text{if } d_{kk} = 3.
  \end{cases}
\]
Then $r^{*} \in (0,1)$ (it has no digits $0$ or $9$, so it lies strictly
between $0$ and $1$ and its expansion does not terminate). For every
$n \in \N^{*}$, the $n$-th digit of $r^{*}$ differs from the $n$-th digit
of $r_n$, so $r^{*} \neq r_n$. This contradicts the assumption that every
element of $(0,1)$ appears in the list. Therefore $(0,1)$ is uncountable.
\end{proof}

\begin{corollary}\label{cor:powerset-N-uncountable}
$\powerset(\N)$ is uncountable.
\end{corollary}

\begin{proof}
The map $S \mapsto \sum_{n \in S} 2^{-(n+1)}$ gives an injection
$\powerset(\N) \hookrightarrow [0,1]$, and the argument above shows
$\card{\powerset(\N)} = \card{\R} > \card{\N}$.
\end{proof}

\begin{theorem}[Cantor]\label{thm:cantor-powerset}
\index{Cantor's theorem}
For any set $A$, there is no surjection $A \to \powerset(A)$.
In particular $\card{A} < \card{\powerset(A)}$.
\end{theorem}

\begin{proof}
Suppose for contradiction that $f\colon A \to \powerset(A)$ is surjective.
Define $D = \set{a \in A : a \notin f(a)}$. Since $D \subseteq A$, we have
$D \in \powerset(A)$, so by surjectivity there exists $d \in A$ with
$f(d) = D$. Now ask: is $d \in D$?

If $d \in D$, then by definition of $D$, $d \notin f(d) = D$, a
contradiction. If $d \notin D$, then $d \notin f(d)$ is false, i.e.\
$d \in f(d) = D$, again a contradiction. Hence no such surjection exists.
\end{proof}

% ============================================================
\section{Exercises}
\label{sec:ch3-exercises}

\begin{exercise}\label{exer:set-identities}
Prove the following set identities for arbitrary sets $A, B, C$:
\begin{enumerate}[label=(\alph*)]
  \item $A \cap (B \cup C) = (A \cap B) \cup (A \cap C)$ \quad
        (distributivity of $\cap$ over $\cup$).
  \item $A \cup (B \cap C) = (A \cup B) \cap (A \cup C)$ \quad
        (distributivity of $\cup$ over $\cap$).
  \item $A \setminus (B \cup C) = (A \setminus B) \cap (A \setminus C)$.
\end{enumerate}
\end{exercise}

\begin{exercise}\label{exer:equiv-rel}
Let $A = \Z$ and define $a \sim b$ iff $3 \mid (a - b)$.
\begin{enumerate}[label=(\alph*)]
  \item Prove that $\sim$ is an equivalence relation.
  \item List all the equivalence classes explicitly.
  \item Draw a diagram showing how the classes partition $\Z$.
\end{enumerate}
\end{exercise}

\begin{exercise}\label{exer:poset-powerset}
Draw the Hasse diagram of $(\powerset(\set{a,b,c}),\, \subseteq)$.
\end{exercise}

\begin{exercise}\label{exer:injection-surjection}
Let $A$ and $B$ be finite sets with $\card{A} = m$ and $\card{B} = n$.
\begin{enumerate}[label=(\alph*)]
  \item Prove that if there exists an injection $A \hookrightarrow B$,
        then $m \leq n$.
  \item Prove that if there exists a surjection $A \twoheadrightarrow B$,
        then $m \geq n$.
  \item Prove that $\card{A} = \card{B}$ if and only if there exists a
        bijection $A \to B$.
\end{enumerate}
\end{exercise}

\begin{exercise}\label{exer:cantor-bernstein}
\textbf{(Cantor--Bernstein--Schr\"oder, statement only.)}
\index{Cantor--Bernstein--Schr\"oder theorem}
Suppose there exist injections $f\colon A \hookrightarrow B$ and
$g\colon B \hookrightarrow A$. Use this to argue informally that
$\card{A} = \card{B}$. (A full proof is beyond the scope of this course.)
\end{exercise}

\begin{exercise}\label{exer:uncountable-irrationals}
Prove that the set of irrational numbers $\R \setminus \Q$ is uncountable.
\emph{Hint:} if $\Q$ and $\R \setminus \Q$ were both countable, what would
follow about $\R$?
\end{exercise}

\begin{exercise}\label{exer:composition-cancel}
Let $f\colon A \to B$ and $g\colon B \to C$.
\begin{enumerate}[label=(\alph*)]
  \item Show that if $g \circ f$ is injective, then $f$ is injective.
  \item Show that if $g \circ f$ is surjective, then $g$ is surjective.
  \item Give an example where $g \circ f$ is bijective but neither $f$ nor
        $g$ is bijective.
\end{enumerate}
\end{exercise}

% ============================================================
\section{Chapter summary}
\label{sec:ch3-summary}

\begin{itemize}[nosep]
  \item A \textbf{set}\index{set} is an unordered collection of distinct
    objects; key operations are union, intersection, difference, and complement.
  \item \textbf{De Morgan's laws}\index{De Morgan's laws} relate complements
    of unions and intersections.
  \item A \textbf{relation}\index{relation} on $A$ is a subset of $A \times A$.
    Equivalence relations (reflexive, symmetric, transitive) produce partitions;
    partial orders (reflexive, antisymmetric, transitive) are visualised with
    Hasse diagrams\index{Hasse diagram}.
  \item A \textbf{function}\index{function} $f\colon A \to B$ is injective,
    surjective, or bijective; bijections have inverses.
  \item \textbf{Cantor's diagonal argument}\index{Cantor's diagonal argument}
    proves $\R$ is uncountable, and Cantor's theorem shows
    $\card{A} < \card{\powerset(A)}$ for every set $A$.
\end{itemize}

\bigskip


% ############################################################
% CHAPTER 4
% ############################################################
\chapter{Combinatorics --- Counting}
\label{ch:combinatorics}
\index{combinatorics}

\begin{quote}
\itshape
``To count is to enumerate without repetition and without omission.''\\
--- traditional
\end{quote}

\medskip
\noindent
Combinatorics is the art and science of counting. In this chapter we develop
the fundamental counting principles, study permutations and combinations,
prove the binomial theorem and Vandermonde's identity, and explore further
topics such as multiset coefficients, stars and bars, and derangements.

% ============================================================
\section{Basic counting principles}
\label{sec:counting-principles}

\begin{theorem}[Addition principle]\label{thm:addition}
\index{addition principle}
If $A_1, A_2, \ldots, A_k$ are pairwise disjoint finite sets, then
\[
  \card{A_1 \cup A_2 \cup \cdots \cup A_k}
    = \card{A_1} + \card{A_2} + \cdots + \card{A_k}.
\]
\end{theorem}

\begin{theorem}[Multiplication principle]\label{thm:multiplication}
\index{multiplication principle}
If a procedure consists of $k$ successive stages, with $n_i$ choices at
stage $i$ (independent of earlier choices), then the total number of outcomes
is
\[
  n_1 \cdot n_2 \cdots n_k.
\]
Equivalently, $\card{A_1 \times A_2 \times \cdots \times A_k}
= \card{A_1} \cdot \card{A_2} \cdots \card{A_k}$.
\end{theorem}

\begin{example}\label{ex:licence-plates}
A licence plate consists of $3$ letters (from $26$) followed by $4$ digits
(from $10$). The number of distinct plates is $26^3 \times 10^4 = 175{,}760{,}000$.
\end{example}

\begin{example}\label{ex:counting-subsets}
Re-deriving Proposition~\ref{prop:powerset-card}: each of the $n$ elements of
a set is either included or excluded from a subset, giving $2^n$ subsets by
the multiplication principle.
\end{example}

% ============================================================
\section{Permutations}
\label{sec:permutations}
\index{permutation}

\begin{definition}[Permutation]\label{def:permutation}
A \emph{permutation} of a set $S$ is a bijection $\sigma\colon S \to S$.
When $\card{S} = n$, the number of permutations of $S$ is denoted $n!$
(``$n$ factorial'').
\end{definition}

\begin{proposition}\label{prop:n-factorial}
For $n \geq 1$, $n! = n \cdot (n-1) \cdot (n-2) \cdots 2 \cdot 1$.
By convention, $0! = 1$.
\end{proposition}

\begin{proof}
There are $n$ choices for the image of the first element, $n - 1$ for the
second (since it must differ), and so on. By the multiplication principle
the total is $n(n-1)(n-2)\cdots 1 = n!$.
\end{proof}

\begin{definition}[$k$-permutation]\label{def:k-permutation}
\index{permutation!$k$-permutation}
A \emph{$k$-permutation} of a set of $n$ elements is an ordered selection
of $k$ distinct elements. Their number is
\[
  P(n, k) = \frac{n!}{(n-k)!} = n(n-1)\cdots(n-k+1).
\]
\end{definition}

\begin{example}\label{ex:podium}
In a race with $10$ runners, the number of possible podium finishes
(gold, silver, bronze) is $P(10, 3) = 10 \cdot 9 \cdot 8 = 720$.
\end{example}

% ============================================================
\section{Combinations}
\label{sec:combinations}
\index{combination}\index{binomial coefficient}

\begin{definition}[Binomial coefficient]\label{def:binomial-coeff}
The number of ways to choose $k$ elements from a set of $n$ elements
(without regard to order) is
\[
  \binom{n}{k} = \frac{n!}{k!\,(n-k)!}, \qquad 0 \leq k \leq n.
\]
We set $\binom{n}{k} = 0$ for $k < 0$ or $k > n$.
\end{definition}

\begin{proposition}[Symmetry]\label{prop:binom-symmetry}
\index{binomial coefficient!symmetry}
$\displaystyle\binom{n}{k} = \binom{n}{n-k}$ for all $0 \leq k \leq n$.
\end{proposition}

\begin{proof}
$\displaystyle\binom{n}{n-k} = \frac{n!}{(n-k)!\,(n-(n-k))!} = \frac{n!}{(n-k)!\,k!} = \binom{n}{k}$.
\end{proof}

\begin{proposition}[Pascal's identity]\label{prop:pascal-identity}
\index{Pascal's identity}
For $1 \leq k \leq n$,
\[
  \binom{n}{k} = \binom{n-1}{k-1} + \binom{n-1}{k}.
\]
\end{proposition}

\begin{proof}
Consider a distinguished element $x$ in a set of $n$ elements. Every
$k$-element subset either contains $x$ or not:
\begin{itemize}[nosep]
  \item those containing $x$: choose the remaining $k-1$ from $n-1$ elements:
        $\binom{n-1}{k-1}$ subsets;
  \item those not containing $x$: choose all $k$ from $n-1$ elements:
        $\binom{n-1}{k}$ subsets.
\end{itemize}
By the addition principle the total is $\binom{n-1}{k-1} + \binom{n-1}{k}$.
\end{proof}

\begin{example}\label{ex:committee}
A committee of $4$ must be chosen from $10$ people. The number of possible
committees is $\binom{10}{4} = \frac{10!}{4!\,6!} = 210$.
\end{example}

\begin{example}\label{ex:lattice-paths}
\index{lattice path}
The number of lattice paths from $(0,0)$ to $(m,n)$ using only unit steps
right ($R$) and up ($U$) is $\binom{m+n}{n}$, since such a path is
determined by choosing which $n$ of the $m + n$ steps are $U$.
\end{example}

% ============================================================
\section{The binomial theorem}
\label{sec:binomial-theorem}

\begin{theorem}[Binomial theorem]\label{thm:binomial}
\index{binomial theorem}
For any $x, y$ and non-negative integer $n$,
\[
  (x + y)^n = \sum_{k=0}^{n} \binom{n}{k} x^{k}\, y^{n-k}.
\]
\end{theorem}

\begin{proof}
We proceed by induction on $n$.

\emph{Base case.}\; For $n = 0$: $(x+y)^0 = 1 = \binom{0}{0} x^0 y^0$.

\emph{Inductive step.}\; Suppose the identity holds for some $n \geq 0$.
Then
\begin{align*}
  (x + y)^{n+1}
    &= (x+y)\sum_{k=0}^{n}\binom{n}{k}x^k y^{n-k} \\
    &= \sum_{k=0}^{n}\binom{n}{k}x^{k+1}y^{n-k}
       + \sum_{k=0}^{n}\binom{n}{k}x^k y^{n+1-k}.
\end{align*}
Re-index the first sum by setting $j = k + 1$:
\[
  = \sum_{j=1}^{n+1}\binom{n}{j-1}x^{j}y^{n+1-j}
    + \sum_{k=0}^{n}\binom{n}{k}x^k y^{n+1-k}.
\]
The $k = n+1$ term in the first sum contributes $\binom{n}{n}x^{n+1}y^0$,
and the $k = 0$ term in the second contributes $\binom{n}{0}x^0 y^{n+1}$.
For $1 \leq k \leq n$ the coefficients combine by Pascal's identity:
\[
  \binom{n}{k-1} + \binom{n}{k} = \binom{n+1}{k}.
\]
Hence $(x+y)^{n+1} = \sum_{k=0}^{n+1}\binom{n+1}{k}x^k y^{n+1-k}$, which
completes the induction.
\end{proof}

\begin{corollary}\label{cor:binom-special}
Setting specific values in the binomial theorem:
\begin{enumerate}[label=(\roman*)]
  \item $x = y = 1$: $\displaystyle\sum_{k=0}^{n}\binom{n}{k} = 2^n$.
  \item $x = 1,\; y = -1$: $\displaystyle\sum_{k=0}^{n}(-1)^k\binom{n}{k} = 0$
        for $n \geq 1$.
\end{enumerate}
\end{corollary}

% ============================================================
\section{Vandermonde's identity}
\label{sec:vandermonde}

\begin{theorem}[Vandermonde's identity]\label{thm:vandermonde}
\index{Vandermonde's identity}
For non-negative integers $m, n, r$ with $r \leq m + n$,
\[
  \binom{m+n}{r} = \sum_{k=0}^{r} \binom{m}{k}\binom{n}{r-k}.
\]
\end{theorem}

\begin{proof}
\textbf{Combinatorial proof.}\; Consider a group of $m + n$ people divided
into two teams: team $A$ of $m$ people and team $B$ of $n$ people. We wish to
choose a committee of $r$ people from the full group.

On the left side, this count is $\binom{m+n}{r}$.

On the right side, we classify committees by the number $k$ of members drawn
from team $A$. If $k$ members come from $A$, then $r - k$ come from $B$.
The number of such committees is $\binom{m}{k}\binom{n}{r-k}$. Summing over
all valid $k$ (from $0$ to $r$, noting that $\binom{m}{k} = 0$ when $k > m$
and $\binom{n}{r-k} = 0$ when $r - k > n$) gives the right-hand side.
\end{proof}

\begin{example}\label{ex:vandermonde-special}
Setting $m = n = r$ gives:
$\displaystyle\binom{2n}{n} = \sum_{k=0}^{n} \binom{n}{k}^2$.
\end{example}

% ============================================================
\section{Pascal's triangle}
\label{sec:pascal-triangle}
\index{Pascal's triangle}

Pascal's triangle\index{Pascal's triangle} arranges the binomial coefficients
in a triangular array where entry $(n, k)$ is $\binom{n}{k}$ and each entry
is the sum of the two entries above it (Pascal's identity,
Proposition~\ref{prop:pascal-identity}).

\begin{figure}[ht]
\centering
\begin{tikzpicture}[scale=0.85]
  \foreach \n in {0,...,7} {
    \foreach \k in {0,...,\n} {
      \pgfmathtruncatemacro{\coeff}{
        \n!/((\k)!*(\n-\k)!)
      }
      \node at (\k - \n/2, -\n*0.85) {\coeff};
    }
  }
\end{tikzpicture}
\caption{Pascal's triangle, rows $n = 0$ through $n = 7$.}
\label{fig:pascal-triangle}
\end{figure}

\begin{remark}\label{rem:pascal-patterns}
Several patterns are visible in Pascal's triangle:
\begin{enumerate}[label=(\roman*)]
  \item Each row is symmetric: $\binom{n}{k} = \binom{n}{n-k}$.
  \item The entries on each edge are $1$.
  \item The sum of row $n$ is $2^n$.
  \item The ``hockey-stick'' identity: $\sum_{i=0}^{r}\binom{n+i}{i} = \binom{n+r+1}{r}$.
\end{enumerate}
\end{remark}

% ============================================================
\section{Multiset coefficients and stars and bars}
\label{sec:stars-bars}
\index{multiset coefficient}\index{stars and bars}

\begin{definition}[Multiset]\label{def:multiset}
\index{multiset}
A \emph{multiset} is a collection of objects where repetition is allowed.
A $k$-element multiset chosen from a set of $n$ types is called a
\emph{$k$-combination with repetition}.
\end{definition}

\begin{theorem}[Stars and bars]\label{thm:stars-bars}
\index{stars and bars}
The number of ways to choose $k$ elements from $n$ types with repetition
allowed (equivalently, the number of solutions in non-negative integers to
$x_1 + x_2 + \cdots + x_n = k$) is
\[
  \binom{n + k - 1}{k} = \binom{n + k - 1}{n - 1}.
\]
\end{theorem}

\begin{proof}
Represent a selection as a string of $k$ stars (\texttt{*}) and $n - 1$ bars
(\texttt{|}). The stars represent the chosen items and the bars separate the
$n$ types. For example, with $n = 4$ types and $k = 6$, the string
\[
  \texttt{** | * | | ***}
\]
encodes $x_1 = 2,\; x_2 = 1,\; x_3 = 0,\; x_4 = 3$.

Each valid string consists of $k + (n-1)$ symbols in total, and is uniquely
determined by choosing which $n - 1$ of these positions are bars (or
equivalently, which $k$ positions are stars). Hence the count is
$\binom{k + n - 1}{n - 1} = \binom{n + k - 1}{k}$.
\end{proof}

\begin{example}\label{ex:stars-bars-coins}
A child wants to buy lollipops at a shop that stocks $4$ flavours. If the
child buys $6$ lollipops, how many flavour combinations are possible?

Using stars and bars with $n = 4$ and $k = 6$:
$\binom{4 + 6 - 1}{6} = \binom{9}{6} = 84$.
\end{example}

\begin{example}\label{ex:positive-solutions}
\index{stars and bars!positive integers}
Find the number of solutions in \emph{positive} integers to
$x_1 + x_2 + x_3 = 10$.

Substitute $y_i = x_i - 1$ (so $y_i \geq 0$) to obtain
$y_1 + y_2 + y_3 = 7$. By stars and bars, the answer is
$\binom{7 + 3 - 1}{7} = \binom{9}{7} = 36$.
\end{example}

\begin{proposition}[Multinomial coefficient]\label{prop:multinomial}
\index{multinomial coefficient}
The number of ways to partition a set of $n$ objects into groups of sizes
$k_1, k_2, \ldots, k_r$ (with $k_1 + \cdots + k_r = n$) is the
\emph{multinomial coefficient}
\[
  \binom{n}{k_1, k_2, \ldots, k_r} = \frac{n!}{k_1!\, k_2!\, \cdots k_r!}.
\]
\end{proposition}

\begin{example}\label{ex:anagram}
\index{anagram}
The number of distinct arrangements of the letters in \textsc{MISSISSIPPI}
is
\[
  \frac{11!}{1!\,4!\,4!\,2!} = 34{,}650.
\]
(There is $1$\,M, $4$\,I's, $4$\,S's, and $2$\,P's.)
\end{example}

% ============================================================
\section{Derangements}
\label{sec:derangements}
\index{derangement}

\begin{definition}[Derangement]\label{def:derangement}
A \emph{derangement} of a set $\set{1, 2, \ldots, n}$ is a permutation
$\sigma$ with no fixed points, i.e.\ $\sigma(i) \neq i$ for all $i$.
The number of derangements of $n$ elements is denoted $D_n$.
\end{definition}

\begin{theorem}\label{thm:derangement-formula}
\index{derangement!formula}
\[
  D_n = n!\sum_{k=0}^{n}\frac{(-1)^k}{k!}.
\]
In particular, $D_n$ is the nearest integer to $n!/e$ for $n \geq 1$.
\end{theorem}

\begin{proof}
Let $A_i$ be the set of permutations of $\set{1, \ldots, n}$ that fix $i$.
By inclusion--exclusion\index{inclusion--exclusion},
\[
  D_n = n! - \card{A_1 \cup A_2 \cup \cdots \cup A_n}
      = \sum_{k=0}^{n}(-1)^k \binom{n}{k}(n-k)!.
\]
For any $k$-element subset $S \subseteq \set{1,\ldots,n}$,
$\card{\bigcap_{i \in S} A_i} = (n-k)!$ (the remaining $n-k$ elements may
be permuted freely). There are $\binom{n}{k}$ such subsets, so
inclusion--exclusion gives
\[
  D_n = \sum_{k=0}^{n}(-1)^k \binom{n}{k}(n-k)!
      = \sum_{k=0}^{n}(-1)^k \frac{n!}{k!}
      = n!\sum_{k=0}^{n}\frac{(-1)^k}{k!}.\qedhere
\]
\end{proof}

\begin{example}\label{ex:derangement-small}
Computing small values:
\begin{center}
\begin{tabular}{c|cccccccc}
\toprule
$n$ & 0 & 1 & 2 & 3 & 4 & 5 & 6 & 7 \\ \midrule
$D_n$ & 1 & 0 & 1 & 2 & 9 & 44 & 265 & 1854 \\
\bottomrule
\end{tabular}
\end{center}
As $n \to \infty$, the probability that a random permutation is a derangement
converges to $1/e \approx 0.3679$.
\end{example}

\begin{proposition}[Recurrence for derangements]\label{prop:derangement-recurrence}
\index{derangement!recurrence}
For $n \geq 2$,
\[
  D_n = (n-1)(D_{n-1} + D_{n-2}),
\]
with $D_0 = 1$ and $D_1 = 0$.
\end{proposition}

\begin{proof}
Consider element $1$. In a derangement, $\sigma(1) = j$ for some
$j \in \set{2, \ldots, n}$, giving $n - 1$ choices. Two cases arise:
\begin{itemize}[nosep]
  \item \emph{Case 1: $\sigma(j) = 1$.}\; Then elements $1$ and $j$ swap,
    and the remaining $n - 2$ elements must be deranged: $D_{n-2}$ ways.
  \item \emph{Case 2: $\sigma(j) \neq 1$.}\; Define a permutation $\tau$ of
    $\set{2, \ldots, n}$ by $\tau(i) = \sigma(i)$ for $i \neq j$ and
    $\tau(j) = \sigma(j)$. The condition $\sigma(j) \neq 1$ together with
    $\sigma(i) \neq i$ for $i \neq 1,j$ means $\tau$ is a derangement of
    $n - 1$ elements: $D_{n-1}$ ways.
\end{itemize}
By the addition and multiplication principles,
$D_n = (n-1)(D_{n-1} + D_{n-2})$.
\end{proof}

% ============================================================
\section{Worked examples}
\label{sec:ch4-worked-examples}

\begin{example}[Committees with restrictions]\label{ex:committee-restricted}
From $8$ men and $6$ women, choose a committee of $5$ that includes at
least $2$ women.

\emph{Solution.}\; Count by the number $w$ of women ($w \geq 2$):
\[
  \sum_{w=2}^{5}\binom{6}{w}\binom{8}{5-w}
  = \binom{6}{2}\binom{8}{3} + \binom{6}{3}\binom{8}{2}
    + \binom{6}{4}\binom{8}{1} + \binom{6}{5}\binom{8}{0}
  = 840 + 560 + 120 + 6 = 1526.
\]
\end{example}

\begin{example}[Distributing identical objects]\label{ex:distribute-identical}
Distribute $15$ identical balls into $4$ distinct boxes such that each box
has at least $2$ balls.

\emph{Solution.}\; Let $x_i$ be the number of balls in box $i$. We need the
number of solutions to $x_1 + x_2 + x_3 + x_4 = 15$ with each $x_i \geq 2$.
Substitute $y_i = x_i - 2$:
\[
  y_1 + y_2 + y_3 + y_4 = 7, \qquad y_i \geq 0.
\]
By stars and bars: $\binom{7 + 3}{3} = \binom{10}{3} = 120$.
\end{example}

\begin{example}[Paths on a grid]\label{ex:grid-paths}
\index{lattice path}
How many shortest lattice paths go from $(0,0)$ to $(5,4)$ passing through
$(2,2)$?

\emph{Solution.}\; A path through $(2,2)$ splits into two independent segments:
$(0,0) \to (2,2)$ and $(2,2) \to (5,4)$. The counts are
$\binom{4}{2} = 6$ and $\binom{5}{2} = 10$, giving $6 \times 10 = 60$ paths.
\end{example}

\begin{example}[Application of the binomial theorem]\label{ex:binomial-application}
Find the coefficient of $x^7$ in $(2x - 3)^{10}$.

\emph{Solution.}\; By the binomial theorem,
\[
  (2x - 3)^{10} = \sum_{k=0}^{10}\binom{10}{k}(2x)^k(-3)^{10-k}.
\]
The $x^7$ term is $\binom{10}{7}(2)^7(-3)^3 = 120 \cdot 128 \cdot (-27) = -414{,}720$.
\end{example}

\begin{example}[Hat problem (derangements)]\label{ex:hat-problem}
\index{hat problem}
Ten guests check their hats. If the hats are returned randomly, what is the
probability that nobody receives their own hat?

\emph{Solution.}\; This is a derangement problem. The probability is
\[
  \frac{D_{10}}{10!} = \sum_{k=0}^{10}\frac{(-1)^k}{k!}
  \approx \frac{1}{e} \approx 0.3679.
\]
\end{example}

% ============================================================
\section{Exercises}
\label{sec:ch4-exercises}

\begin{exercise}\label{exer:counting-words}
How many $5$-letter ``words'' (sequences of letters) can be formed from the
English alphabet if:
\begin{enumerate}[label=(\alph*)]
  \item repetition is allowed?
  \item no letter may be repeated?
  \item the word must begin with a vowel and end with a consonant, with
        repetition allowed?
\end{enumerate}
\end{exercise}

\begin{exercise}\label{exer:binary-strings}
\index{binary string}
How many binary strings of length $12$ contain exactly $5$ ones?
\end{exercise}

\begin{exercise}\label{exer:binomial-identity}
Use the binomial theorem to prove that
$\displaystyle\sum_{k=0}^{n} k\binom{n}{k} = n \cdot 2^{n-1}$.\\
\emph{Hint:} differentiate $(1 + x)^n$ and set $x = 1$.
\end{exercise}

\begin{exercise}\label{exer:stars-bars-upper}
Find the number of solutions in non-negative integers to
$x_1 + x_2 + x_3 + x_4 = 20$ with $x_1 \leq 6$.\\
\emph{Hint:} total solutions minus those with $x_1 \geq 7$.
\end{exercise}

\begin{exercise}\label{exer:derangement-formula-check}
Verify the recurrence $D_n = (n-1)(D_{n-1} + D_{n-2})$ for $n = 4$ and
$n = 5$ by listing all derangements of $\set{1,2,3,4}$.
\end{exercise}

\begin{exercise}\label{exer:vandermonde-application}
Use Vandermonde's identity to evaluate
$\displaystyle\sum_{k=0}^{5}\binom{10}{k}\binom{15}{5-k}$.
\end{exercise}

\begin{exercise}\label{exer:multinomial}
\begin{enumerate}[label=(\alph*)]
  \item How many distinct arrangements are there of the letters in
        \textsc{ARRANGEMENT}?
  \item In how many of these do the two R's appear next to each other?
\end{enumerate}
\end{exercise}

\begin{exercise}\label{exer:committee-president}
A club of $20$ members must choose a president, a secretary, and a committee
of $5$ (which may include the president and secretary). In how many ways can
this be done?
\end{exercise}

\begin{exercise}\label{exer:diagonal-sum}
Prove that $\displaystyle\sum_{k=0}^{n}\binom{n}{k}^2 = \binom{2n}{n}$
using Vandermonde's identity.
\end{exercise}

\begin{exercise}\label{exer:derangement-subfactorial}
Show that $D_n = nD_{n-1} + (-1)^n$ for $n \geq 1$.\\
\emph{Hint:} use the explicit formula
$D_n = n!\sum_{k=0}^{n}(-1)^k/k!$.
\end{exercise}

% ============================================================
\section{Chapter summary}
\label{sec:ch4-summary}

\begin{itemize}[nosep]
  \item The \textbf{addition}\index{addition principle} and
    \textbf{multiplication}\index{multiplication principle} principles are
    the cornerstones of counting.
  \item \textbf{Permutations}\index{permutation}: $n!$ ways to arrange $n$
    objects; $P(n,k)$ for $k$-permutations.
  \item \textbf{Combinations}\index{combination}: $\binom{n}{k}$ counts
    unordered selections; the \textbf{binomial theorem}\index{binomial theorem}
    $(x+y)^n = \sum_k \binom{n}{k}x^k y^{n-k}$.
  \item \textbf{Vandermonde's identity}\index{Vandermonde's identity}:
    $\binom{m+n}{r} = \sum_k \binom{m}{k}\binom{n}{r-k}$.
  \item \textbf{Stars and bars}\index{stars and bars}: the number of
    non-negative integer solutions to $x_1 + \cdots + x_n = k$ is
    $\binom{n+k-1}{k}$.
  \item \textbf{Derangements}\index{derangement}: $D_n = n!\sum_{k=0}^{n}(-1)^k/k!
    \approx n!/e$.
\end{itemize}

% === END OF CHUNK ch3-4 ===
% === START OF CHUNK ch5-6 ===

% ============================================================
%  Chapter 5 — Inclusion-Exclusion and Advanced Counting
% ============================================================

\chapter{Inclusion-Exclusion and Advanced Counting}\label{ch:inclusion-exclusion}

% ------------------------------------------------------------
%  5.0  Motivating introduction
% ------------------------------------------------------------

In the previous chapters we developed the basic tools of enumerative
combinatorics: the addition and multiplication principles, permutations,
combinations, and the binomial theorem.  These tools are remarkably
powerful, yet they share a common limitation: they work best when the
objects being counted fall into \emph{disjoint} categories.  In
practice, the sets we wish to count frequently overlap, and we must
account for the elements that belong to several sets at once.

The \textbf{principle of inclusion-exclusion}\index{inclusion-exclusion}
provides a systematic way to handle such overlaps.  It generalises the
familiar identity $\card{A \cup B} = \card{A} + \card{B} - \card{A \cap B}$
to an arbitrary finite union, and it leads to elegant formulas for
counting surjections, computing the Euler totient function, and
enumerating derangements.

We then turn to two further themes in advanced counting:
\emph{Stirling numbers}\index{Stirling numbers} and
\emph{integer partitions}\index{integer partition}, which refine the
art of distributing objects into groups.


% ============================================================
%  5.1  Inclusion-Exclusion for Two and Three Sets
% ============================================================

\section{Inclusion-Exclusion for Two and Three Sets}\label{sec:ie-two-three}

\begin{theorem}[Inclusion-Exclusion for Two Sets]\label{thm:ie-two}
\index{inclusion-exclusion!two sets}
Let $A$ and $B$ be finite sets.  Then
\[
  \card{A \cup B} \;=\; \card{A} + \card{B} - \card{A \cap B}.
\]
\end{theorem}

\begin{proof}
Every element of $A \cup B$ belongs to at least one of $A$ and $B$.
Write $A \cup B$ as the disjoint union
\[
  A \cup B \;=\; (A \setminus B)\;\cup\;(A \cap B)\;\cup\;(B \setminus A).
\]
Hence $\card{A \cup B} = \card{A \setminus B} + \card{A \cap B} +
\card{B \setminus A}$.  Since $\card{A} = \card{A \setminus B} +
\card{A \cap B}$ and $\card{B} = \card{B \setminus A} +
\card{A \cap B}$, adding gives
\[
  \card{A} + \card{B}
  = \card{A \setminus B} + 2\card{A \cap B} + \card{B \setminus A}
  = \card{A \cup B} + \card{A \cap B},
\]
and the result follows.
\end{proof}

% --- Venn diagram for two sets ---
\begin{center}
\begin{tikzpicture}[scale=1.0]
  \begin{scope}[fill opacity=0.25]
    \fill[blue] (-0.8,0) circle (1.5);
    \fill[red]  ( 0.8,0) circle (1.5);
  \end{scope}
  \draw (-0.8,0) circle (1.5) node[left=1.2cm] {$A$};
  \draw ( 0.8,0) circle (1.5) node[right=1.2cm] {$B$};
  \node at (0,0) {$A \cap B$};
  \node[below] at (0,-2) {$\card{A \cup B} = \card{A} + \card{B} - \card{A \cap B}$};
\end{tikzpicture}
\end{center}

\begin{theorem}[Inclusion-Exclusion for Three Sets]\label{thm:ie-three}
\index{inclusion-exclusion!three sets}
Let $A$, $B$, and $C$ be finite sets.  Then
\begin{align*}
  \card{A \cup B \cup C}
  &= \card{A} + \card{B} + \card{C} \\
  &\quad - \card{A \cap B} - \card{A \cap C} - \card{B \cap C} \\
  &\quad + \card{A \cap B \cap C}.
\end{align*}
\end{theorem}

\begin{proof}
Apply the two-set formula twice:
\begin{align*}
  \card{A \cup B \cup C}
  &= \card{(A \cup B) \cup C} \\
  &= \card{A \cup B} + \card{C} - \card{(A \cup B) \cap C} \\
  &= \card{A \cup B} + \card{C} - \card{(A \cap C) \cup (B \cap C)}.
\end{align*}
Expanding $\card{A \cup B}$ and $\card{(A \cap C) \cup (B \cap C)}$ via
Theorem~\ref{thm:ie-two} and noting that
$(A \cap C) \cap (B \cap C) = A \cap B \cap C$ yields the result.
\end{proof}

% --- Venn diagram for three sets ---
\begin{center}
\begin{tikzpicture}[scale=1.0]
  \begin{scope}[fill opacity=0.2]
    \fill[blue]   (90:1)  circle (1.5);
    \fill[red]    (210:1) circle (1.5);
    \fill[green!70!black]  (330:1) circle (1.5);
  \end{scope}
  \draw (90:1)  circle (1.5) node[above=1.0cm] {$A$};
  \draw (210:1) circle (1.5) node[below left=0.7cm] {$B$};
  \draw (330:1) circle (1.5) node[below right=0.7cm] {$C$};
  \node at (0,0) {\footnotesize $A \!\cap\! B \!\cap\! C$};
\end{tikzpicture}
\end{center}

\begin{example}\label{ex:ie-languages}
In a class of $100$ students, $40$ study French, $35$ study German, and
$20$ study Spanish.  Furthermore, $12$ study both French and German,
$8$ study both French and Spanish, $6$ study both German and Spanish,
and $2$ study all three languages.  How many students study at least one
of the three languages?

Let $F$, $G$, $S$ denote the respective sets.  By inclusion-exclusion,
\[
  \card{F \cup G \cup S}
  = 40 + 35 + 20 - 12 - 8 - 6 + 2
  = 71.
\]
Hence $100 - 71 = 29$ students study none of these languages.
\end{example}


% ============================================================
%  5.2  The General Inclusion-Exclusion Principle
% ============================================================

\section{The General Principle}\label{sec:ie-general}

\begin{theorem}[General Inclusion-Exclusion]\label{thm:ie-general}
\index{inclusion-exclusion!general principle}
Let $A_1, A_2, \dots, A_n$ be finite subsets of a finite set~$U$.  Then
\[
  \card*{\bigcup_{i=1}^{n} A_i}
  \;=\;
  \sum_{k=1}^{n} (-1)^{k+1}
  \!\!\sum_{1 \le i_1 < \cdots < i_k \le n}
  \card{A_{i_1} \cap A_{i_2} \cap \cdots \cap A_{i_k}}.
\]
Equivalently, writing $S_k = \sum_{1\le i_1<\cdots<i_k\le n}
\card{A_{i_1}\cap\cdots\cap A_{i_k}}$,
\[
  \card*{\bigcup_{i=1}^n A_i}
  = S_1 - S_2 + S_3 - \cdots + (-1)^{n+1} S_n.
\]
\end{theorem}

\begin{proof}
We show that each element $x \in \bigcup_{i=1}^{n} A_i$ is counted
exactly once on the right-hand side.  Suppose $x$ belongs to exactly
$m \ge 1$ of the sets $A_1, \dots, A_n$.  Then $x$ is counted once in
each $\binom{m}{k}$ of the intersections of size~$k$, so the
right-hand side counts $x$ a total of
\[
  \sum_{k=1}^{m} (-1)^{k+1}\binom{m}{k}
  = -\sum_{k=1}^{m}(-1)^{k}\binom{m}{k}
  = -\!\Biggl[\sum_{k=0}^{m}(-1)^k \binom{m}{k} - 1\Biggr]
  = -(0 - 1) = 1,
\]
where we used the binomial theorem: $\sum_{k=0}^{m}(-1)^k\binom{m}{k}
= (1-1)^m = 0$ for $m \ge 1$.  Hence every element of the union is
counted exactly once.
\end{proof}

\begin{remark}\label{rem:ie-complement}
It is often more convenient to count the \emph{complement}.  If $A_1,
\dots, A_n \subseteq U$, then
\[
  \card*{U \setminus \bigcup_{i=1}^n A_i}
  = \card{U} - S_1 + S_2 - S_3 + \cdots + (-1)^n S_n,
\]
where $S_k$ is defined as above.  This form is especially useful when
we want to count elements of $U$ that belong to \emph{none} of the
$A_i$.
\end{remark}


% ============================================================
%  5.3  Counting Surjections
% ============================================================

\section{Counting Surjections}\label{sec:surjections}

\begin{theorem}[Number of Surjections]\label{thm:surjections}
\index{surjection!counting}
The number of surjections from a set of $n$ elements onto a set of
$m$ elements (with $n \ge m$) is
\[
  \operatorname{Surj}(n,m)
  = \sum_{k=0}^{m} (-1)^k \binom{m}{k}(m-k)^n.
\]
\end{theorem}

\begin{proof}
Let $X = \set{x_1, \dots, x_n}$ and $Y = \set{y_1, \dots, y_m}$, and
let $U$ be the set of all functions $f\colon X \to Y$, so
$\card{U} = m^n$.  For each $j \in \set{1, \dots, m}$, let $A_j$ be
the set of functions that \emph{miss} $y_j$, i.e.\
$A_j = \set{f \in U : y_j \notin f(X)}$.  A function is surjective if
and only if it lies in none of the~$A_j$, so
\[
  \operatorname{Surj}(n,m) = \card*{U \setminus \bigcup_{j=1}^{m} A_j}.
\]
For any $k$-element subset $\set{j_1, \dots, j_k} \subseteq
\set{1,\dots,m}$, the set $A_{j_1} \cap \cdots \cap A_{j_k}$ consists
of functions whose range avoids $k$ specified elements, hence
$\card{A_{j_1}\cap\cdots\cap A_{j_k}} = (m-k)^n$.  Since there are
$\binom{m}{k}$ such subsets, inclusion-exclusion
(Remark~\ref{rem:ie-complement}) gives
\[
  \operatorname{Surj}(n,m)
  = \sum_{k=0}^{m} (-1)^k \binom{m}{k}(m-k)^n. \qedhere
\]
\end{proof}

\begin{example}\label{ex:surj-4-3}
The number of surjections from $\set{1,2,3,4}$ onto $\set{a,b,c}$ is
\[
  \binom{3}{0}3^4 - \binom{3}{1}2^4 + \binom{3}{2}1^4 - \binom{3}{3}0^4
  = 81 - 48 + 3 - 0 = 36.
\]
\end{example}


% ============================================================
%  5.4  Euler's Totient Function via IE
% ============================================================

\section{Euler's Totient Function}\label{sec:euler-totient}

\begin{definition}[Euler's Totient Function]\label{def:euler-totient}
\index{Euler's totient function}
For a positive integer~$n$, the \textbf{Euler totient function}
$\varphi(n)$\index{phi@$\varphi(n)$} counts the number of integers
$k$ with $1 \le k \le n$ and $\gcd(k, n) = 1$.
\end{definition}

\begin{theorem}[Product Formula for $\varphi$]\label{thm:totient-product}
If $n = p_1^{a_1} p_2^{a_2} \cdots p_r^{a_r}$ is the prime
factorisation of $n \ge 2$, then
\[
  \varphi(n)
  = n \prod_{i=1}^{r}\Bigl(1 - \frac{1}{p_i}\Bigr).
\]
\end{theorem}

\begin{proof}
Let $U = \set{1, 2, \dots, n}$ and, for each $i$, let $A_i$ be the set
of elements of $U$ divisible by~$p_i$.  Then $\varphi(n) = \card{U
\setminus \bigcup_{i=1}^r A_i}$.  We have $\card{A_i} = n/p_i$ and,
more generally, for any subset $\set{i_1,\dots,i_k}\subseteq\set{1,\dots,r}$,
\[
  \card{A_{i_1} \cap \cdots \cap A_{i_k}}
  = \frac{n}{p_{i_1}\cdots p_{i_k}},
\]
since the $p_i$ are distinct primes.  By inclusion-exclusion,
\begin{align*}
  \varphi(n)
  &= \sum_{k=0}^{r} (-1)^k
     \sum_{1\le i_1<\cdots<i_k\le r}
     \frac{n}{p_{i_1}\cdots p_{i_k}} \\
  &= n \sum_{k=0}^{r} (-1)^k
     \sum_{1\le i_1<\cdots<i_k\le r}
     \frac{1}{p_{i_1}\cdots p_{i_k}}
  = n \prod_{i=1}^{r}\Bigl(1-\frac{1}{p_i}\Bigr),
\end{align*}
where the last step expands the product.
\end{proof}

\begin{example}\label{ex:totient-60}
For $n = 60 = 2^2 \cdot 3 \cdot 5$,
\[
  \varphi(60) = 60\Bigl(1-\tfrac{1}{2}\Bigr)\Bigl(1-\tfrac{1}{3}\Bigr)
  \Bigl(1-\tfrac{1}{5}\Bigr) = 60 \cdot \tfrac{1}{2}\cdot \tfrac{2}{3}
  \cdot \tfrac{4}{5} = 16.
\]
\end{example}


% ============================================================
%  5.5  Derangements
% ============================================================

\section{Derangements}\label{sec:derangements}

\begin{definition}[Derangement]\label{def:derangement}
\index{derangement}
A \textbf{derangement} of $\set{1, 2, \dots, n}$ is a permutation
$\sigma$ such that $\sigma(i) \neq i$ for all~$i$.  The number of
derangements of an $n$-element set is denoted $D_n$\index{D@$D_n$}.
\end{definition}

\begin{theorem}[Derangement Formula]\label{thm:derangement}
\index{derangement!formula}
\[
  D_n = n!\sum_{k=0}^{n}\frac{(-1)^k}{k!}.
\]
In particular, $D_n/n! \to 1/e$ as $n \to \infty$, so approximately
a fraction $1/e \approx 36.8\%$ of all permutations are derangements.
\end{theorem}

\begin{proof}
Let $U$ be the set of all permutations of $\set{1,\dots,n}$ and let
$A_i = \set{\sigma \in U : \sigma(i) = i}$ be the set of permutations
that fix~$i$.  A derangement is a permutation in none of the~$A_i$, so
$D_n = \card{U \setminus \bigcup_{i=1}^n A_i}$.

For any $k$-element subset $\set{i_1,\dots,i_k}$, the permutations
fixing all of $i_1, \dots, i_k$ form a set of size $(n-k)!$.
There are $\binom{n}{k}$ such subsets.  By inclusion-exclusion,
\begin{align*}
  D_n &= \sum_{k=0}^{n}(-1)^k \binom{n}{k}(n-k)!
       = \sum_{k=0}^{n}(-1)^k \frac{n!}{k!}
       = n!\sum_{k=0}^{n}\frac{(-1)^k}{k!}. \qedhere
\end{align*}
\end{proof}

\begin{example}\label{ex:derangement-small}
The first several values of $D_n$ are:
\begin{center}
\begin{tabular}{c|ccccccc}
  $n$   & 0 & 1 & 2 & 3 & 4  & 5  & 6 \\ \hline
  $D_n$ & 1 & 0 & 1 & 2 & 9  & 44 & 265
\end{tabular}
\end{center}
For instance, $D_4 = 4!(1 - 1 + \tfrac{1}{2} - \tfrac{1}{6} +
\tfrac{1}{24}) = 24 \cdot \tfrac{3}{8} = 9$.
\end{example}

\begin{remark}\label{rem:derangement-recursive}
Derangements also satisfy the recurrence
\[
  D_n = (n-1)(D_{n-1} + D_{n-2}), \quad n \ge 2,
\]
with $D_0 = 1$ and $D_1 = 0$.  To see this, consider where element~$1$
is sent: say $\sigma(1) = j \neq 1$.  If $\sigma(j) = 1$ (a
``swap''), the remaining $n-2$ elements must be deranged, giving
$D_{n-2}$.  If $\sigma(j) \neq 1$, we must derange $\set{1,\dots,n}
\setminus\set{1}$ with the constraint $\sigma(j)\neq 1$, which is
equivalent to a derangement of $n-1$ elements, giving $D_{n-1}$.
Since there are $n-1$ choices for~$j$, the recurrence follows.
\end{remark}

\begin{exercise}\label{exo:hat-check}
\index{hat-check problem}
\textbf{(Hat-Check Problem.)}
At a party, $n$ guests each check a hat.  On leaving, the hats are
returned at random.  What is the probability that no guest receives
their own hat?  Show that this probability is $D_n / n!$ and compute
it for $n = 5$.
\end{exercise}


% ============================================================
%  5.6  Stirling Numbers of the Second Kind
% ============================================================

\section{Stirling Numbers of the Second Kind}\label{sec:stirling}

\begin{definition}[Stirling Number of the Second Kind]\label{def:stirling-second}
\index{Stirling number!second kind}
The \textbf{Stirling number of the second kind}
$S(n,k)$\index{S(n,k)@$S(n,k)$} (also written $\genfrac{\{}{\}}{0pt}{}{n}{k}$)
is the number of ways to partition a set of $n$ elements into exactly
$k$ non-empty subsets.
\end{definition}

\begin{proposition}[Explicit Formula]\label{prop:stirling-explicit}
\[
  S(n,k) = \frac{1}{k!}\sum_{j=0}^{k}(-1)^j\binom{k}{j}(k-j)^n.
\]
\end{proposition}

\begin{proof}
The number of surjections from an $n$-set to a $k$-set is
$\operatorname{Surj}(n,k) = \sum_{j=0}^k(-1)^j\binom{k}{j}(k-j)^n$
by Theorem~\ref{thm:surjections}.  Each partition into $k$ non-empty
blocks can be labelled in $k!$ ways to give a surjection, so
$\operatorname{Surj}(n,k) = k!\,S(n,k)$.
\end{proof}

\begin{proposition}[Recurrence]\label{prop:stirling-recurrence}
\index{Stirling number!recurrence}
For $n, k \ge 1$,
\[
  S(n,k) = k\,S(n-1,k) + S(n-1,k-1),
\]
with $S(0,0)=1$ and $S(n,0)=S(0,k)=0$ for $n,k\ge 1$.
\end{proposition}

\begin{proof}
Consider element~$n$.  Either it forms a singleton block
($S(n-1,k-1)$ ways to partition the remaining $n-1$ elements into
$k-1$ blocks), or it is added to one of the $k$ existing blocks of a
partition of $\set{1,\dots,n-1}$ into $k$ blocks ($k\,S(n-1,k)$ ways).
\end{proof}

\begin{example}\label{ex:stirling-table}
A table of Stirling numbers $S(n,k)$ for small values:
\begin{center}
\begin{tabular}{c|cccccc}
  $n \setminus k$ & 0 & 1 & 2 & 3 & 4 & 5 \\ \hline
  0 & 1 \\
  1 & 0 & 1 \\
  2 & 0 & 1 & 1 \\
  3 & 0 & 1 & 3 & 1 \\
  4 & 0 & 1 & 7 & 6  & 1 \\
  5 & 0 & 1 & 15& 25 & 10 & 1
\end{tabular}
\end{center}
\end{example}


% ============================================================
%  5.7  Bell Numbers
% ============================================================

\section{Bell Numbers}\label{sec:bell-numbers}

\begin{definition}[Bell Number]\label{def:bell}
\index{Bell number}
The \textbf{Bell number} $B_n$\index{B@$B_n$ (Bell number)} is the
total number of partitions of a set of $n$ elements:
\[
  B_n = \sum_{k=0}^{n} S(n,k).
\]
\end{definition}

\begin{proposition}[Bell Recurrence]\label{prop:bell-recurrence}
\[
  B_{n+1} = \sum_{k=0}^{n}\binom{n}{k} B_k, \quad B_0 = 1.
\]
\end{proposition}

\begin{proof}
Consider element $n+1$.  The block containing $n+1$ has $k$ other
elements (chosen in $\binom{n}{k}$ ways), and the remaining $n-k$
elements are partitioned in $B_{n-k}$ ways.  Summing over $k$ and
re-indexing gives the result.
\end{proof}

\begin{example}\label{ex:bell-values}
The first few Bell numbers are:
$B_0=1$, $B_1=1$, $B_2=2$, $B_3=5$, $B_4=15$, $B_5=52$, $B_6=203$.
\end{example}


% ============================================================
%  5.8  Integer Partitions
% ============================================================

\section{Integer Partitions}\label{sec:integer-partitions}

\begin{definition}[Integer Partition]\label{def:integer-partition}
\index{integer partition}
A \textbf{partition} of a positive integer $n$ is a way of writing $n$
as a sum of positive integers, where the order of the summands does not
matter.  Each summand is called a \textbf{part}\index{part}.  The number
of partitions of $n$ is denoted $p(n)$\index{p(n)@$p(n)$}.
\end{definition}

\begin{example}\label{ex:partitions-5}
The partitions of $5$ are:
\[
  5, \quad 4+1, \quad 3+2, \quad 3+1+1, \quad 2+2+1, \quad 2+1+1+1,
  \quad 1+1+1+1+1,
\]
so $p(5) = 7$.
\end{example}

\begin{notation}\label{not:partition-lambda}
We write a partition as $\lambda = (\lambda_1, \lambda_2, \dots,
\lambda_\ell)$ with $\lambda_1 \ge \lambda_2 \ge \cdots \ge
\lambda_\ell \ge 1$ and $\lambda_1 + \cdots + \lambda_\ell = n$.  The
number~$\ell$ is the \textbf{length}\index{partition!length} of
$\lambda$.
\end{notation}

\begin{theorem}[Generating Function for $p(n)$]\label{thm:partition-gf}
\index{integer partition!generating function}
\[
  \sum_{n=0}^{\infty} p(n)\,x^n
  = \prod_{k=1}^{\infty} \frac{1}{1-x^k}.
\]
\end{theorem}

\begin{proof}
Each factor $\frac{1}{1-x^k} = 1 + x^k + x^{2k} + \cdots$ encodes
choosing how many parts equal to~$k$ appear.  Multiplying over all
$k \ge 1$ and collecting the coefficient of $x^n$ counts the number
of partitions of~$n$.
\end{proof}

\begin{theorem}[Euler's Distinct--Odd Theorem]\label{thm:euler-distinct-odd}
\index{Euler's distinct--odd theorem}
The number of partitions of $n$ into \emph{distinct} parts equals
the number of partitions of $n$ into \emph{odd} parts.
\end{theorem}

\begin{proof}
The generating function for distinct-part partitions is
$\prod_{k=1}^{\infty}(1+x^k)$, and for odd-part partitions it is
$\prod_{k=0}^{\infty}\frac{1}{1-x^{2k+1}}$.  We verify they are equal:
\[
  \prod_{k=1}^{\infty}(1+x^k)
  = \prod_{k=1}^{\infty}\frac{1-x^{2k}}{1-x^k}
  = \frac{\prod_{k=1}^{\infty}(1-x^{2k})}{\prod_{k=1}^{\infty}(1-x^k)}
  = \frac{1}{\prod_{j\text{ odd}}(1-x^j)},
\]
since the even factors cancel.
\end{proof}


% ============================================================
%  5.9  Exercises
% ============================================================

\section{Exercises}\label{sec:ie-exercises}

\begin{exercise}\label{exo:ie-divisibility}
How many integers in $\set{1, 2, \dots, 1000}$ are divisible by
neither $3$, $5$, nor $7$?
\end{exercise}

\begin{exercise}\label{exo:ie-surj-5-3}
Compute the number of surjections from $\set{1,\dots,5}$ onto
$\set{a,b,c}$.
\end{exercise}

\begin{exercise}\label{exo:totient-properties}
Prove that $\sum_{d \mid n}\varphi(d) = n$ for every positive
integer~$n$.
\end{exercise}

\begin{exercise}\label{exo:derangement-nearest-int}
Show that $D_n$ is the nearest integer to $n!/e$ for all $n \ge 1$.
\end{exercise}

\begin{exercise}\label{exo:stirling-identity}
Prove that $\sum_{k=0}^{n} S(n,k)\,x^{\underline{k}} = x^n$,
where $x^{\underline{k}} = x(x-1)\cdots(x-k+1)$ is the falling
factorial.
\end{exercise}

\begin{exercise}\label{exo:bell-triangle}
\index{Bell triangle}
The \textbf{Bell triangle} is constructed by placing $B_0 = 1$ at the
top and filling each row so that each entry is the sum of its left and
upper-left neighbours.  Construct rows $0$ through~$5$ and verify that
the left-most entry in each row gives $B_n$.
\end{exercise}

\begin{exercise}\label{exo:partitions-10}
List all partitions of $8$ into distinct parts and all partitions of
$8$ into odd parts, thereby verifying Euler's theorem for $n=8$.
\end{exercise}

\begin{exercise}\label{exo:menage}
\index{probl\`eme des m\'enages}
\textbf{(Probl\`eme des M\'enages.)}
$n$ married couples sit around a circular table so that men and women
alternate and no husband sits next to his wife.  Using
inclusion-exclusion, show that the number of such seatings is
\[
  2\,n!\sum_{k=0}^{n}(-1)^k \frac{2n}{2n-k}\binom{2n-k}{k}(n-k)!.
\]
\end{exercise}


% ============================================================
%  5.10  Chapter Summary
% ============================================================

\section{Chapter Summary}\label{sec:ie-summary}

\begin{itemize}
  \item \textbf{Inclusion-exclusion} (Theorem~\ref{thm:ie-general})
    computes $\card{\bigcup A_i}$ by alternately adding and subtracting
    the sizes of intersections.

  \item \textbf{Surjections:} $\operatorname{Surj}(n,m) = \sum_{k=0}^m
    (-1)^k\binom{m}{k}(m-k)^n$ (Theorem~\ref{thm:surjections}).

  \item \textbf{Euler's totient:} $\varphi(n) = n\prod_{p\mid n}
    (1-1/p)$ (Theorem~\ref{thm:totient-product}).

  \item \textbf{Derangements:} $D_n = n!\sum_{k=0}^n(-1)^k/k!$
    (Theorem~\ref{thm:derangement}).

  \item \textbf{Stirling numbers} $S(n,k)$ count set partitions into
    $k$ blocks; they satisfy
    $S(n,k) = kS(n\!-\!1,k) + S(n\!-\!1,k\!-\!1)$.

  \item \textbf{Bell numbers} $B_n = \sum_k S(n,k)$ count all set
    partitions; they satisfy $B_{n+1}=\sum_k\binom{n}{k}B_k$.

  \item \textbf{Integer partitions:} the generating function
    $\prod_{k\ge 1}(1-x^k)^{-1}$ encodes $p(n)$; distinct-part
    partitions are equinumerous with odd-part partitions.
\end{itemize}


% ============================================================
% ============================================================
%  Chapter 6 — Ordinary and Exponential Generating Functions
% ============================================================
% ============================================================

\chapter{Ordinary and Exponential Generating Functions}\label{ch:generating-functions}

% ------------------------------------------------------------
%  6.0  Motivating introduction
% ------------------------------------------------------------

Generating functions are one of the most versatile tools in
combinatorics.  The idea is deceptively simple: encode a sequence
$(a_0, a_1, a_2, \dots)$ as the coefficients of a formal power series,
and then use the algebraic and analytic properties of the series to
extract information about the sequence.

\index{generating function}
George P\'olya\index{P\'olya, George} called generating functions
``a clothesline on which we hang up a sequence of numbers for display.''
Herb Wilf\index{Wilf, Herbert} more vividly wrote that a generating
function is ``a device somewhat similar to a bag: instead of carrying
many little objects, we put them all in a bag, and then we have only
one object to carry.''

In this chapter we study two principal types:
\begin{itemize}
  \item \textbf{Ordinary generating functions (OGFs)}, suited to
    problems of selection and combination;
  \item \textbf{Exponential generating functions (EGFs)}, suited to
    problems of arrangement and permutation.
\end{itemize}


% ============================================================
%  6.1  Formal Power Series
% ============================================================

\section{Formal Power Series}\label{sec:formal-power-series}

Before we do combinatorics, we set up the algebraic framework in which
generating functions live.

\begin{definition}[Formal Power Series]\label{def:fps}
\index{formal power series}
A \textbf{formal power series} over a commutative ring $R$ is a
sequence $(a_0, a_1, a_2, \dots)$ of elements of $R$, written
\[
  f(x) = \sum_{n=0}^{\infty} a_n x^n.
\]
The set of all such series is denoted $R[[x]]$\index{R[[x]]@$R[[x]]$}.
\end{definition}

\begin{definition}[Operations on Formal Power Series]\label{def:fps-ops}
Let $f(x) = \sum a_n x^n$ and $g(x) = \sum b_n x^n$ be elements of
$R[[x]]$.
\begin{enumerate}[label=(\roman*)]
  \item \textbf{Addition:}\index{formal power series!addition}
    $f + g = \sum_{n \ge 0}(a_n + b_n)x^n$.
  \item \textbf{Multiplication:}\index{formal power series!multiplication}
    $f \cdot g = \sum_{n \ge 0} c_n x^n$ where
    $c_n = \sum_{k=0}^n a_k b_{n-k}$
    (the \emph{Cauchy product}\index{Cauchy product}).
  \item \textbf{Coefficient extraction:} We write $[x^n]f(x) = a_n$.
\end{enumerate}
\end{definition}

\begin{proposition}[Multiplicative Inverses]\label{prop:fps-inverse}
\index{formal power series!inverse}
A formal power series $f(x) = \sum a_n x^n \in R[[x]]$ has a
multiplicative inverse if and only if $a_0$ is a unit in $R$.  When
$R$ is a field, this means $a_0 \neq 0$.
\end{proposition}

\begin{proof}
Suppose $f(x)g(x) = 1$ with $g(x) = \sum b_n x^n$.  Comparing
constant terms gives $a_0 b_0 = 1$, so $a_0$ must be a unit.
Conversely, if $a_0$ is a unit, define $b_0 = a_0^{-1}$ and
recursively
\[
  b_n = -a_0^{-1}\sum_{k=1}^{n}a_k b_{n-k}, \quad n \ge 1.
\]
One verifies that $f(x)g(x) = 1$.
\end{proof}

\begin{remark}\label{rem:fps-convergence}
In the theory of formal power series, \emph{convergence is
irrelevant}.  We never substitute a numerical value for~$x$; the
variable~$x$ is merely a bookkeeping device.  This makes the algebra
entirely rigorous without any appeal to analysis.  Of course, when
the series also converges (e.g.\ for $\abs{x} < R$), we gain
additional tools from analysis.
\end{remark}


% ============================================================
%  6.2  Ordinary Generating Functions
% ============================================================

\section{Ordinary Generating Functions}\label{sec:ogf}

\begin{definition}[OGF]\label{def:ogf}
\index{ordinary generating function}
\index{OGF|see{ordinary generating function}}
The \textbf{ordinary generating function} (OGF) of a sequence
$(a_n)_{n \ge 0}$ is
\[
  A(x) = \sum_{n=0}^{\infty} a_n x^n.
\]
\end{definition}

\subsection{Basic OGFs}\label{ssec:basic-ogfs}

We collect some standard OGFs that will be used repeatedly.

\begin{proposition}[Catalogue of OGFs]\label{prop:ogf-catalogue}
\index{ordinary generating function!catalogue}
\begin{enumerate}[label=(\roman*)]
  \item $\displaystyle\frac{1}{1-x} = \sum_{n\ge 0}x^n$
    \quad (constant sequence $a_n = 1$).
  \item $\displaystyle\frac{1}{(1-x)^2} = \sum_{n\ge 0}(n+1)x^n$.
  \item $\displaystyle\frac{1}{(1-x)^k} = \sum_{n\ge 0}
    \binom{n+k-1}{k-1}x^n$ \quad for $k \ge 1$.
  \item $\displaystyle\frac{x}{(1-x)^2} = \sum_{n\ge 0}n\,x^n$.
  \item $(1+x)^m = \sum_{n=0}^{m}\binom{m}{n}x^n$ \quad for
    $m \in \N$.
  \item $\displaystyle\frac{1}{1-x-x^2}$ is the OGF of the
    Fibonacci numbers (see Section~\ref{sec:ogf-fibonacci}).
\end{enumerate}
\end{proposition}

\subsection{Operations on OGFs}\label{ssec:ogf-operations}

\begin{proposition}[OGF Operations]\label{prop:ogf-ops}
\index{ordinary generating function!operations}
Let $A(x) = \sum a_n x^n$ and $B(x) = \sum b_n x^n$.
\begin{enumerate}[label=(\roman*)]
  \item \textbf{Scaling:} $A(cx) = \sum a_n c^n x^n$.
  \item \textbf{Right shift:} $x^k A(x) = \sum_{n\ge k} a_{n-k} x^n$.
  \item \textbf{Differentiation:} $A'(x) = \sum_{n\ge 0}(n+1)a_{n+1}x^n$.
  \item \textbf{Convolution:} $A(x)B(x) = \sum_{n\ge 0}c_n x^n$
    where $c_n = \sum_{k=0}^n a_k b_{n-k}$.
  \item \textbf{Partial sums:} $\frac{A(x)}{1-x} = \sum_{n\ge 0}
    \bigl(\sum_{k=0}^n a_k\bigr)x^n$.
\end{enumerate}
\end{proposition}


% ============================================================
%  6.3  Solving Recurrences with OGFs
% ============================================================

\section{Solving Recurrences with OGFs}\label{sec:ogf-recurrences}

Generating functions provide a systematic three-step method for solving
linear recurrences with constant coefficients:
\begin{enumerate}
  \item Multiply the recurrence by $x^n$ and sum over all valid~$n$ to
    obtain a functional equation for the OGF.
  \item Solve algebraically for the OGF.
  \item Extract coefficients using partial fractions or known series.
\end{enumerate}

\subsection{Worked Example: Fibonacci Numbers}\label{sec:ogf-fibonacci}

\begin{theorem}[OGF of Fibonacci Numbers]\label{thm:fib-ogf}
\index{Fibonacci numbers!OGF}
Define $(F_n)$ by $F_0 = 0$, $F_1 = 1$, and $F_n = F_{n-1}+F_{n-2}$
for $n \ge 2$.  Then
\[
  F(x) = \sum_{n=0}^{\infty}F_n x^n = \frac{x}{1-x-x^2}.
\]
\end{theorem}

\begin{proof}
Let $F(x) = \sum_{n\ge 0}F_n x^n$.  Multiplying $F_n = F_{n-1}+F_{n-2}$
by $x^n$ and summing for $n \ge 2$:
\[
  \sum_{n\ge 2}F_n x^n
  = \sum_{n\ge 2}F_{n-1}x^n + \sum_{n\ge 2}F_{n-2}x^n.
\]
The left side is $F(x) - F_0 - F_1 x = F(x) - x$.  The right side is
$x(F(x)-F_0) + x^2 F(x) = xF(x) + x^2 F(x)$.  Hence
\[
  F(x) - x = xF(x) + x^2 F(x)
  \implies
  F(x)(1 - x - x^2) = x
  \implies
  F(x) = \frac{x}{1-x-x^2}. \qedhere
\]
\end{proof}

\begin{corollary}[Binet's Formula]\label{cor:binet}
\index{Binet's formula}
Let $\phi = (1+\sqrt{5})/2$ and $\hat\phi = (1-\sqrt{5})/2$.  Then
\[
  F_n = \frac{\phi^n - \hat\phi^n}{\sqrt{5}}.
\]
\end{corollary}

\begin{proof}
Factor $1 - x - x^2 = -(x - 1/\phi)(x - 1/\hat\phi)$ and decompose
$F(x)$ into partial fractions:
\[
  F(x) = \frac{x}{1-x-x^2}
  = \frac{1}{\sqrt{5}}\left(\frac{1}{1-\phi x} - \frac{1}{1-\hat\phi x}\right).
\]
Expanding each geometric series and extracting $[x^n]$ gives the result.
\end{proof}


% ============================================================
%  6.4  Catalan Numbers
% ============================================================

\section{Catalan Numbers}\label{sec:catalan}

\begin{definition}[Catalan Numbers]\label{def:catalan}
\index{Catalan numbers}
The \textbf{Catalan numbers} $C_n$\index{C@$C_n$ (Catalan)} are
defined by $C_0 = 1$ and the recurrence
\[
  C_{n+1} = \sum_{k=0}^{n} C_k\,C_{n-k}, \quad n \ge 0.
\]
\end{definition}

\begin{theorem}[OGF and Closed Form for Catalan Numbers]\label{thm:catalan-ogf}
\index{Catalan numbers!generating function}
The OGF $C(x) = \sum_{n\ge 0}C_n x^n$ satisfies
\[
  C(x) = \frac{1 - \sqrt{1-4x}}{2x},
\]
and the closed form is
\[
  C_n = \frac{1}{n+1}\binom{2n}{n}.
\]
\end{theorem}

\begin{proof}
The recurrence $C_{n+1} = \sum_{k=0}^n C_k C_{n-k}$ says that the
sequence $(C_1, C_2, \dots)$ is the Cauchy product of $(C_0, C_1,
\dots)$ with itself.  In terms of generating functions,
\[
  \frac{C(x) - 1}{x} = C(x)^2,
\]
i.e.\ $xC(x)^2 - C(x) + 1 = 0$.  By the quadratic formula,
\[
  C(x) = \frac{1 \pm \sqrt{1 - 4x}}{2x}.
\]
Since $C(0) = C_0 = 1$, we need the branch with the minus sign (the
plus sign gives $C(x) \to \infty$ as $x \to 0$).

For the closed form, use the generalised binomial theorem:
\[
  \sqrt{1-4x} = \sum_{n=0}^{\infty}\binom{1/2}{n}(-4x)^n.
\]
One computes
$\binom{1/2}{n}(-4)^n = \frac{(-1)^{n-1}\cdot 2}{n}\binom{2(n-1)}{n-1}$
for $n \ge 1$.  After simplification,
\[
  C(x) = \sum_{n=0}^{\infty}\frac{1}{n+1}\binom{2n}{n}\,x^n. \qedhere
\]
\end{proof}

\begin{example}\label{ex:catalan-values}
The first several Catalan numbers:
$C_0=1$, $C_1=1$, $C_2=2$, $C_3=5$, $C_4=14$, $C_5=42$, $C_6=132$.
\end{example}

\begin{remark}[Catalan Interpretations]\label{rem:catalan-interpretations}
\index{Catalan numbers!interpretations}
The Catalan number $C_n$ counts, among many other things:
\begin{enumerate}[label=(\roman*)]
  \item the number of \textbf{ballot paths}\index{ballot path}
    (lattice paths from $(0,0)$ to $(2n,0)$ using steps $(1,1)$ and
    $(1,-1)$ that never go below the $x$-axis);
  \item the number of ways to \textbf{triangulate}\index{triangulation}
    a convex $(n+2)$-gon;
  \item the number of strings of $n$ pairs of \textbf{matched
    parentheses}\index{parenthesisation};
  \item the number of \textbf{full binary trees}\index{binary tree}
    with $n$ internal nodes.
\end{enumerate}
\end{remark}

% --- TikZ: Ballot paths for C_3 = 5 ---
\begin{center}
\begin{tikzpicture}[scale=0.55, every node/.style={font=\tiny}]
  % We draw the 5 ballot paths from (0,0) to (6,0) for C_3
  \foreach \col/\xshift/\path in {%
    blue!70!black/0/{1,1,1,-1,-1,-1},%
    red!70!black/7.5/{1,1,-1,1,-1,-1},%
    green!50!black/15/{1,1,-1,-1,1,-1},%
    orange!80!black/22.5/{1,-1,1,1,-1,-1},%
    purple!70!black/30/{1,-1,1,-1,1,-1}%
  }{%
    \begin{scope}[shift={(\xshift,0)}]
      \draw[gray!40, thin] (0,0) grid (6,3);
      \draw[thick, ->] (-0.3,0) -- (6.5,0);
      \draw[\col, very thick] (0,0)
        \foreach \s [count=\i from 0] in \path {%
          -- ++(1,\s)
        };
    \end{scope}
  }
  \node[below] at (3,-0.6) {Path 1};
  \node[below] at (10.5,-0.6) {Path 2};
  \node[below] at (18,-0.6) {Path 3};
  \node[below] at (25.5,-0.6) {Path 4};
  \node[below] at (33,-0.6) {Path 5};
  \node[below] at (16.5,-1.5) {The $C_3 = 5$ ballot paths from $(0,0)$ to $(6,0)$};
\end{tikzpicture}
\end{center}


% ============================================================
%  6.5  Exponential Generating Functions
% ============================================================

\section{Exponential Generating Functions}\label{sec:egf}

When we count \emph{labelled} structures (permutations, arrangements,
labelled graphs), the natural generating function weights the $n$th
term by $1/n!$ instead of $1$.

\begin{definition}[EGF]\label{def:egf}
\index{exponential generating function}
\index{EGF|see{exponential generating function}}
The \textbf{exponential generating function} (EGF) of a sequence
$(a_n)_{n \ge 0}$ is
\[
  \hat{A}(x) = \sum_{n=0}^{\infty} a_n \frac{x^n}{n!}.
\]
\end{definition}

\begin{proposition}[Basic EGFs]\label{prop:egf-catalogue}
\index{exponential generating function!catalogue}
\begin{enumerate}[label=(\roman*)]
  \item $e^x = \sum_{n\ge 0}\frac{x^n}{n!}$
    \quad (sequence $a_n = 1$).
  \item $e^{cx} = \sum_{n\ge 0}c^n\frac{x^n}{n!}$
    \quad (sequence $a_n = c^n$).
  \item $\frac{1}{1-x} = \sum_{n\ge 0}n!\,\frac{x^n}{n!}$
    \quad (sequence $a_n = n!$, i.e.\ permutations).
  \item $\displaystyle -\ln(1-x) = \sum_{n\ge 1}(n-1)!\,\frac{x^n}{n!}$
    \quad (sequence $a_n = (n-1)!$ for $n\ge 1$, i.e.\ cyclic
    permutations).
\end{enumerate}
\end{proposition}

\subsection{The Exponential Formula}\label{ssec:exponential-formula}

The key property of EGFs is that \emph{multiplication corresponds to
shuffling labels}.

\begin{theorem}[Product of EGFs]\label{thm:egf-product}
\index{exponential generating function!product}
If $\hat{A}(x) = \sum a_n x^n/n!$ and $\hat{B}(x) = \sum b_n x^n/n!$,
then
\[
  \hat{A}(x)\,\hat{B}(x)
  = \sum_{n=0}^{\infty}\Biggl(\sum_{k=0}^{n}\binom{n}{k}a_k\,b_{n-k}\Biggr)
  \frac{x^n}{n!}.
\]
Thus the coefficient of $x^n/n!$ in the product counts ways to split
$\set{1,\dots,n}$ into an ordered pair $(S,T)$ with $\card{S}=k$ and
$\card{T}=n-k$, place an ``$A$-structure'' on $S$ and a
``$B$-structure'' on~$T$, summed over all~$k$.
\end{theorem}

\begin{proof}
Direct computation of the Cauchy product:
\[
  \hat{A}(x)\hat{B}(x)
  = \sum_{n\ge 0}\sum_{k=0}^{n}\frac{a_k}{k!}\cdot\frac{b_{n-k}}{(n-k)!}\,x^n
  = \sum_{n\ge 0}\frac{1}{n!}\sum_{k=0}^{n}\binom{n}{k}a_k b_{n-k}\,x^n.
  \qedhere
\]
\end{proof}


% ============================================================
%  6.6  Permutations and Derangements via EGF
% ============================================================

\section{Permutations and Derangements via EGF}\label{sec:egf-permutations}

\begin{example}[EGF for Permutations]\label{ex:egf-permutations}
\index{permutation!EGF}
The number of permutations of an $n$-set is $n!$, so the EGF is
\[
  \sum_{n=0}^{\infty}n!\,\frac{x^n}{n!}
  = \sum_{n=0}^{\infty}x^n = \frac{1}{1-x}.
\]
\end{example}

\begin{theorem}[EGF for Derangements]\label{thm:derangement-egf}
\index{derangement!EGF}
The EGF for the derangement numbers $D_n$ is
\[
  \hat{D}(x) = \sum_{n=0}^{\infty}D_n\frac{x^n}{n!}
  = \frac{e^{-x}}{1-x}.
\]
\end{theorem}

\begin{proof}
From Theorem~\ref{thm:derangement}, $D_n = n!\sum_{k=0}^n(-1)^k/k!$.
This is the binomial convolution $D_n = \sum_{k=0}^n \binom{n}{k}(-1)^k
(n-k)!$, which by Theorem~\ref{thm:egf-product} corresponds to
\[
  \hat{D}(x) = e^{-x}\cdot\frac{1}{1-x}. \qedhere
\]
\end{proof}

\begin{remark}\label{rem:derangement-egf-alt}
Alternatively, write $D_n/n! = \sum_{k=0}^n(-1)^k/k!$ and recognise
that $D_n/n! \to e^{-1}$ as $n \to \infty$.  In fact,
$\hat{D}(x) = e^{-x}/(1-x)$ gives the elegant interpretation: a
permutation is a derangement times a set of fixed points.  That is,
we partition $\set{1,\dots,n}$ into a set of fixed points (whose EGF
is $e^x$) and a derangement on the rest.  Since the EGF for all
permutations is $1/(1-x)$, we get $1/(1-x) = e^x \hat{D}(x)$.
\end{remark}


% ============================================================
%  6.7  Bell Numbers via EGF
% ============================================================

\section{Bell Numbers via EGF}\label{sec:egf-bell}

\begin{theorem}[EGF for Bell Numbers]\label{thm:bell-egf}
\index{Bell number!EGF}
\[
  \sum_{n=0}^{\infty}B_n\frac{x^n}{n!} = e^{e^x - 1}.
\]
\end{theorem}

\begin{proof}
A partition of $\set{1,\dots,n}$ is a set of non-empty subsets (the
blocks).  The EGF for a single non-empty subset (``at least one
element'') is $e^x - 1$.  By the compositional formula for EGFs (the
``exponential formula''\index{exponential formula}), the EGF for
\emph{sets of non-empty subsets} is
\[
  \exp\bigl(e^x - 1\bigr). \qedhere
\]
\end{proof}

\begin{remark}\label{rem:bell-Dobinski}
\index{Dobi\'nski's formula}
Expanding $e^{e^x-1} = e^{-1}\sum_{k=0}^{\infty}\frac{e^{kx}}{k!}$
and extracting $[x^n/n!]$ gives \textbf{Dobi\'nski's formula}:
\[
  B_n = \frac{1}{e}\sum_{k=0}^{\infty}\frac{k^n}{k!}.
\]
This remarkable identity expresses $B_n$ as a convergent infinite sum.
\end{remark}


% ============================================================
%  6.8  Worked Example: A Recurrence via OGF
% ============================================================

\section{Worked Example: A General Recurrence}\label{sec:ogf-worked}

\begin{example}[Solving $a_n = 5a_{n-1} - 6a_{n-2}$]\label{ex:ogf-recurrence}
Let $a_0 = 1$, $a_1 = 4$, and $a_n = 5a_{n-1} - 6a_{n-2}$ for $n\ge2$.
Find a closed form for $a_n$.

\textbf{Step 1.}  Let $A(x) = \sum a_n x^n$.  Multiply the recurrence
by $x^n$ and sum for $n \ge 2$:
\[
  A(x) - 1 - 4x = 5x(A(x) - 1) - 6x^2 A(x).
\]
Solving: $A(x)(1 - 5x + 6x^2) = 1 - x$, so
\[
  A(x) = \frac{1-x}{1-5x+6x^2} = \frac{1-x}{(1-2x)(1-3x)}.
\]

\textbf{Step 2.}  Partial fractions: write
$\frac{1-x}{(1-2x)(1-3x)} = \frac{\alpha}{1-2x}+\frac{\beta}{1-3x}$.
Setting $x = 1/2$: $\alpha = (1-1/2)/(1-3/2) = (1/2)/(-1/2) = -1$.
Setting $x = 1/3$: $\beta = (1-1/3)/(1-2/3) = (2/3)/(1/3) = 2$.
Hence
\[
  A(x) = \frac{-1}{1-2x} + \frac{2}{1-3x}
  = -\sum_{n\ge 0}2^n x^n + 2\sum_{n\ge 0}3^n x^n.
\]

\textbf{Step 3.}  Extract coefficients:
$a_n = 2\cdot 3^n - 2^n$.
We verify: $a_0 = 2-1 = 1$, $a_1 = 6-2 = 4$, $a_2 = 18-4 = 14 =
5(4)-6(1) = 14$.\checkmark
\end{example}


% ============================================================
%  6.9  Worked Example: Derangements via OGF recurrence
% ============================================================

\section{Worked Example: Derangements Revisited}\label{sec:ogf-derangements}

\begin{example}[Derangement OGF from the Recurrence]\label{ex:derangement-ogf}
\index{derangement!OGF approach}
The recurrence $D_n = (n-1)(D_{n-1}+D_{n-2})$ (Remark~\ref{rem:derangement-recursive})
does not have constant coefficients, so the OGF approach is less
direct.  Instead, we use the EGF.

From $D_n = n!\sum_{k=0}^n(-1)^k/k!$, define $d_n = D_n/n!$.  Then
$d_n = \sum_{k=0}^n(-1)^k/k!$, so $d_n - d_{n-1} = (-1)^n/n!$.
The EGF $\hat{D}(x) = \sum d_n n!\,\frac{x^n}{n!} = \sum D_n x^n/n!$
but it is more illuminating to note:

The relation $D_n = nD_{n-1} + (-1)^n$ gives
\[
  \frac{D_n}{n!} = \frac{D_{n-1}}{(n-1)!}\cdot\frac{1}{n}\cdot n + \frac{(-1)^n}{n!}
\]
which simplifies to $\frac{D_n}{n!} - \frac{D_{n-1}}{(n-1)!} = \frac{(-1)^n}{n!}$.
Summing from $n=1$ to $N$:
\[
  \frac{D_N}{N!} = \sum_{n=0}^{N}\frac{(-1)^n}{n!},
\]
confirming $\hat{D}(x) = e^{-x}/(1-x)$ once more.
\end{example}


% ============================================================
%  6.10  Composition and the Exponential Formula
% ============================================================

\section{Composition and the Exponential Formula}\label{sec:exponential-formula}

\begin{theorem}[Exponential Formula]\label{thm:exponential-formula}
\index{exponential formula}
Let $\hat{C}(x) = \sum_{n\ge 1}c_n\,x^n/n!$ be the EGF of a class of
connected (or indecomposable) labelled structures, with $c_0 = 0$.
Then the EGF for \emph{sets of such structures} is
\[
  \hat{A}(x) = \exp\!\bigl(\hat{C}(x)\bigr)
  = \sum_{n=0}^{\infty}a_n\frac{x^n}{n!}.
\]
Conversely, $\hat{C}(x) = \ln\hat{A}(x)$.
\end{theorem}

\begin{example}\label{ex:exp-formula-permutations}
A permutation is a set of cycles.  The EGF for a single cycle of length
$n\ge 1$ is $(n-1)!\,x^n/n! = x^n/n$, so
\[
  \hat{C}(x) = \sum_{n=1}^{\infty}\frac{x^n}{n} = -\ln(1-x).
\]
The exponential formula gives
\[
  \exp\!\bigl(-\ln(1-x)\bigr) = \frac{1}{1-x},
\]
which is indeed the EGF for $n!$ (the number of permutations).
\end{example}

\begin{example}\label{ex:exp-formula-involutions}
\index{involution}
An \textbf{involution} is a permutation $\sigma$ with $\sigma^2 = \mathrm{id}$,
i.e.\ every cycle has length $1$ or $2$.  The EGF for fixed points
(cycles of length $1$) is $x$, and for transpositions (cycles of length
$2$) it is $x^2/2$.  So the EGF for involutions is
\[
  \exp\!\Bigl(x + \frac{x^2}{2}\Bigr).
\]
Writing $I_n = [x^n/n!]\exp(x+x^2/2)$, we get $I_0=1$, $I_1=1$,
$I_2=2$, $I_3=4$, $I_4=10$, $I_5=26$.
\end{example}


% ============================================================
%  6.11  Further Applications
% ============================================================

\section{Further Applications}\label{sec:gf-further}

\subsection{Counting with Restrictions}\label{ssec:gf-restrictions}

\begin{example}[Compositions with Restricted Parts]\label{ex:ogf-restricted}
\index{composition!restricted parts}
A \textbf{composition}\index{composition} of $n$ into parts from a set
$S \subseteq \N_{\ge 1}$ is an ordered tuple $(s_1,\dots,s_k)$ with
$s_i \in S$ and $s_1+\cdots+s_k = n$.  The OGF for the number of such
compositions is
\[
  \frac{1}{1 - \sum_{s\in S}x^s}.
\]
For instance, if $S = \set{1,2}$, the OGF is $\frac{1}{1-x-x^2}$,
giving the Fibonacci numbers as the number of compositions of $n$ into
parts $1$ and $2$ (offset by one).
\end{example}

\subsection{The Snake Oil Method}\label{ssec:snake-oil}

\begin{example}[A Binomial Identity]\label{ex:snake-oil}
\index{snake oil method}
We prove the \textbf{Vandermonde identity}\index{Vandermonde identity}
\[
  \sum_{k=0}^{r}\binom{m}{k}\binom{n}{r-k} = \binom{m+n}{r}
\]
using generating functions.  The left side is the coefficient of $x^r$
in $(1+x)^m(1+x)^n = (1+x)^{m+n}$, which is $\binom{m+n}{r}$.
\end{example}


% ============================================================
%  6.12  Exercises
% ============================================================

\section{Exercises}\label{sec:gf-exercises}

\begin{exercise}\label{exo:ogf-tower}
Find a closed form for $\sum_{n=0}^{\infty}n^2 x^n$ by
differentiating $\sum x^n = 1/(1-x)$.
\end{exercise}

\begin{exercise}\label{exo:ogf-recurrence-1}
Solve the recurrence $a_n = 3a_{n-1} - 2a_{n-2}$ with $a_0=0$,
$a_1=1$ using OGFs.
\end{exercise}

\begin{exercise}\label{exo:catalan-parentheses}
Show directly (without generating functions) that the number of valid
strings of $n$ pairs of parentheses satisfies the Catalan recurrence.
\end{exercise}

\begin{exercise}\label{exo:catalan-reflection}
\index{reflection principle}
\textbf{(Reflection Principle.)}
Show that the number of lattice paths from $(0,0)$ to $(2n,0)$ with
steps $(1,1)$ and $(1,-1)$ that touch or cross the $x$-axis negatively
is $\binom{2n}{n-1}$, and deduce $C_n = \binom{2n}{n}-\binom{2n}{n-1}
= \frac{1}{n+1}\binom{2n}{n}$.
\end{exercise}

\begin{exercise}\label{exo:egf-fixed-points}
\index{permutation!with $k$ fixed points}
Using EGFs, show that the number of permutations of $\set{1,\dots,n}$
with exactly $k$ fixed points is $\binom{n}{k}D_{n-k}$.
\end{exercise}

\begin{exercise}\label{exo:egf-connected-graphs}
\index{connected graph}
Let $g_n = 2^{\binom{n}{2}}$ be the number of labelled graphs on $n$
vertices, and let $c_n$ be the number of \emph{connected} labelled
graphs.  Using the exponential formula, write down the relation between
$\sum g_n x^n/n!$ and $\sum c_n x^n/n!$, and compute $c_1$, $c_2$,
$c_3$, $c_4$.
\end{exercise}

\begin{exercise}\label{exo:ogf-partition-distinct}
Show that the OGF for the number of partitions of $n$ into distinct
parts is $\prod_{k=1}^{\infty}(1+x^k)$, and expand it to find the
coefficients up to $n = 10$.
\end{exercise}

\begin{exercise}\label{exo:fibonacci-convolution}
Use the OGF of the Fibonacci numbers to prove that
$\sum_{k=0}^{n}F_k F_{n-k} = \frac{1}{5}\bigl((2n+1)F_n +
2nF_{n-1}\bigr)$.
\textit{Hint:} square $F(x) = x/(1-x-x^2)$ and use partial fractions.
\end{exercise}

\begin{exercise}\label{exo:bell-egf-verify}
Verify the EGF $e^{e^x-1}$ for Bell numbers by expanding through $n=5$
and checking against the known values $B_0 = 1, \dots, B_5 = 52$.
\end{exercise}

\begin{exercise}\label{exo:involution-recurrence}
\index{involution!recurrence}
Show that the number of involutions satisfies $I_n = I_{n-1}+(n-1)I_{n-2}$
for $n \ge 2$, with $I_0 = I_1 = 1$.  Derive this both combinatorially
and from the EGF $\exp(x+x^2/2)$.
\end{exercise}


% ============================================================
%  6.13  Chapter Summary
% ============================================================

\section{Chapter Summary}\label{sec:gf-summary}

\begin{itemize}
  \item A \textbf{formal power series} $\sum a_n x^n \in R[[x]]$
    encodes a sequence; convergence is not required.  The ring $R[[x]]$
    admits inverses whenever $a_0$ is a unit.

  \item \textbf{Ordinary generating functions} $A(x) = \sum a_n x^n$
    are the tool of choice for selections (subsets, multisets,
    compositions, partitions).  Key operations: convolution encodes
    sums of indices; partial fractions extract coefficients.

  \item \textbf{Exponential generating functions}
    $\hat{A}(x) = \sum a_n x^n/n!$ are suited to labelled structures.
    Their product encodes binomial convolution, i.e.\ shuffling labels.

  \item \textbf{Fibonacci numbers:} $F(x) = x/(1-x-x^2)$;
    Binet's formula follows from partial fractions
    (Theorem~\ref{thm:fib-ogf}, Corollary~\ref{cor:binet}).

  \item \textbf{Catalan numbers:} $C(x) = (1-\sqrt{1-4x})/(2x)$;
    closed form $C_n = \binom{2n}{n}/(n+1)$
    (Theorem~\ref{thm:catalan-ogf}).

  \item \textbf{Derangements:} EGF is $e^{-x}/(1-x)$
    (Theorem~\ref{thm:derangement-egf}).

  \item \textbf{Bell numbers:} EGF is $e^{e^x-1}$
    (Theorem~\ref{thm:bell-egf}); Dobi\'nski's formula gives
    $B_n = e^{-1}\sum k^n/k!$.

  \item The \textbf{exponential formula}
    (Theorem~\ref{thm:exponential-formula}): if $\hat{C}(x)$ is the
    EGF of connected structures, then $\exp(\hat{C}(x))$ is the EGF
    for sets of such structures.
\end{itemize}

% === END OF CHUNK ch5-6 ===
% === START OF CHUNK ch7-8 ===

% ============================================================
%  CHAPTER 7 — Recurrences
% ============================================================
\chapter{Recurrences}\label{ch:recurrences}
\index{recurrence}

A \emph{recurrence relation}\index{recurrence!relation} (or simply a
\emph{recurrence}) defines each term of a sequence in terms of one or
more preceding terms together with an initial condition (or several).
Recurrences appear naturally in counting problems, algorithm analysis,
and many other areas of discrete mathematics.  In this chapter we develop
systematic methods for solving them.

% ------------------------------------------------------------
\section{Basic Notions}\label{sec:rec-basic}
% ------------------------------------------------------------

\begin{definition}[Recurrence relation]\label{def:recurrence}
\index{recurrence!relation}
A \emph{recurrence relation} for a sequence $(a_n)_{n\ge 0}$ is an
equation that expresses $a_n$ in terms of one or more of the previous
terms $a_{n-1}, a_{n-2}, \dots$ for all $n$ greater than or equal to
some fixed integer $n_0$.  The values $a_0, a_1, \dots, a_{n_0-1}$ that
are specified independently are called \emph{initial conditions}.
\end{definition}

\begin{example}[Simple recurrence]\label{ex:simple-rec}
The relation $a_n = 2a_{n-1}$ with $a_0 = 3$ defines the sequence
$3, 6, 12, 24, \dots$  Its closed-form solution is $a_n = 3 \cdot 2^n$.
\end{example}

The \emph{order}\index{recurrence!order} of a recurrence is the
difference between the largest and smallest indices that appear.  For
instance $a_n = a_{n-1} + a_{n-2}$ has order~$2$.

% ------------------------------------------------------------
\section{Linear Homogeneous Recurrences}\label{sec:rec-homogeneous}
\index{recurrence!linear homogeneous}
% ------------------------------------------------------------

\begin{definition}[Linear homogeneous recurrence with constant coefficients]\label{def:LHCC}
A recurrence of the form
\[
   a_n = c_1\, a_{n-1} + c_2\, a_{n-2} + \cdots + c_k\, a_{n-k},
   \qquad c_k \neq 0,
\]
where $c_1, \dots, c_k$ are constants, is called a \emph{linear
homogeneous recurrence with constant coefficients} (LHCC) of order~$k$.
\end{definition}

\subsection{The characteristic equation}\label{ssec:char-eq}
\index{characteristic equation}

The key idea is to guess a solution of the form $a_n = r^n$ for some
constant $r \neq 0$.  Substituting into the recurrence yields
\[
   r^n = c_1\, r^{n-1} + c_2\, r^{n-2} + \cdots + c_k\, r^{n-k}.
\]
Dividing through by $r^{n-k}$ gives the \emph{characteristic equation}
\begin{equation}\label{eq:char-eq}
   r^k - c_1\, r^{k-1} - c_2\, r^{k-2} - \cdots - c_k = 0.
\end{equation}

\begin{theorem}[Solution of a second-order LHCC]\label{thm:LHCC-order2}
\index{recurrence!second-order solution}
Consider the recurrence $a_n = c_1\, a_{n-1} + c_2\, a_{n-2}$ with
characteristic equation $r^2 - c_1\, r - c_2 = 0$.
\begin{enumerate}
\item If the characteristic equation has two distinct roots $r_1 \neq r_2$,
      then the general solution is
      \[
         a_n = \alpha\, r_1^n + \beta\, r_2^n,
      \]
      where $\alpha, \beta$ are determined by the initial conditions.
\item If the characteristic equation has a repeated root $r_1 = r_2 = r_0$,
      then the general solution is
      \[
         a_n = (\alpha + \beta\, n)\, r_0^n.
      \]
\end{enumerate}
\end{theorem}

\begin{proof}
\textbf{Case 1.} Since $r_1$ and $r_2$ are roots of the characteristic
equation, both $(r_1^n)$ and $(r_2^n)$ satisfy the recurrence.  Because
the recurrence is linear and homogeneous, any linear combination
$a_n = \alpha\, r_1^n + \beta\, r_2^n$ also satisfies it.  We need to
show that every solution has this form.  Given initial conditions
$a_0, a_1$, we must solve
\[
   \alpha + \beta = a_0, \qquad
   \alpha\, r_1 + \beta\, r_2 = a_1.
\]
The determinant of this system is $r_2 - r_1 \neq 0$ (since the roots
are distinct), so there is a unique solution for $\alpha, \beta$.
Because the recurrence together with initial conditions determines the
sequence uniquely, this must be the solution.

\textbf{Case 2.} When $r_1 = r_2 = r_0$ we have
$r_0 = c_1/2$ and $c_2 = -r_0^2$.  One checks directly that
$a_n = n\, r_0^n$ satisfies the recurrence:
\[
   c_1\, (n-1)\, r_0^{n-1} + c_2\, (n-2)\, r_0^{n-2}
   = r_0^{n-2}\bigl[2r_0(n-1) - r_0^2(n-2)\bigr]\, r_0^{n-2}
\]
simplifies to $n\, r_0^n$ after using $c_1 = 2r_0$ and $c_2 = -r_0^2$.
Hence $\{r_0^n, n\, r_0^n\}$ form a basis, and the general solution is
$a_n = (\alpha + \beta\, n)\, r_0^n$.
\end{proof}

\begin{theorem}[General LHCC of order $k$]\label{thm:LHCC-general}
\index{recurrence!general LHCC}
If the characteristic equation \eqref{eq:char-eq} has $t$ distinct roots
$r_1, \dots, r_t$ with multiplicities $m_1, \dots, m_t$ (so that
$m_1 + \cdots + m_t = k$), then the general solution is
\[
   a_n = \sum_{i=1}^{t}\,
         \bigl(\alpha_{i,0} + \alpha_{i,1}\, n + \cdots
               + \alpha_{i,m_i - 1}\, n^{m_i - 1}\bigr)\, r_i^n,
\]
where the $k$ constants $\alpha_{i,j}$ are determined by the $k$ initial
conditions.
\end{theorem}

\begin{example}[Distinct roots]\label{ex:distinct-roots}
Solve $a_n = 5a_{n-1} - 6a_{n-2}$, with $a_0 = 1$, $a_1 = 4$.

The characteristic equation is $r^2 - 5r + 6 = 0$, whose roots are
$r_1 = 2$ and $r_2 = 3$.  Thus $a_n = \alpha \cdot 2^n + \beta \cdot 3^n$.
The initial conditions give
\[
   \alpha + \beta = 1, \qquad 2\alpha + 3\beta = 4
   \implies \alpha = -1,\; \beta = 2.
\]
Hence $a_n = -2^n + 2 \cdot 3^n = 2 \cdot 3^n - 2^n$.
\end{example}

\begin{example}[Repeated root]\label{ex:repeated-root}
Solve $a_n = 4a_{n-1} - 4a_{n-2}$, with $a_0 = 1$, $a_1 = 6$.

The characteristic equation $r^2 - 4r + 4 = 0$ has the double root
$r_0 = 2$.  So $a_n = (\alpha + \beta n)\, 2^n$.  From $a_0 = 1$ we get
$\alpha = 1$; from $a_1 = 6$ we get $2(\alpha + \beta) = 6$, so
$\beta = 2$.  Therefore $a_n = (1 + 2n)\, 2^n$.
\end{example}

% ------------------------------------------------------------
\section{Non-Homogeneous Linear Recurrences}\label{sec:rec-nonhomogeneous}
\index{recurrence!non-homogeneous}
% ------------------------------------------------------------

\begin{definition}[Non-homogeneous linear recurrence]\label{def:nonhom}
A recurrence of the form
\[
   a_n = c_1\, a_{n-1} + \cdots + c_k\, a_{n-k} + f(n),
\]
where $f(n) \not\equiv 0$, is called a \emph{non-homogeneous} linear
recurrence.
\end{definition}

The general solution equals the general solution of the associated
homogeneous recurrence \emph{plus} any particular solution of the
non-homogeneous recurrence:
\begin{equation}\label{eq:nonhom-structure}
   a_n = a_n^{(\mathrm{h})} + a_n^{(\mathrm{p})}.
\end{equation}

\begin{theorem}[Method of undetermined coefficients]\label{thm:undetermined-coeff}
\index{undetermined coefficients}
If the non-homogeneous term has the form $f(n) = (b_d\, n^d + \cdots + b_0)\, s^n$
for constants $b_i, s$, then one may seek a particular solution of the form
\[
   a_n^{(\mathrm{p})}
   = n^m\, (B_d\, n^d + \cdots + B_0)\, s^n,
\]
where $m$ is the multiplicity of $s$ as a root of the characteristic
equation (set $m = 0$ if $s$ is not a root).
\end{theorem}

\begin{example}[Non-homogeneous recurrence]\label{ex:nonhom}
Solve $a_n = 3a_{n-1} + 4^n$, with $a_0 = 1$.

\emph{Homogeneous part.}  The characteristic equation of $a_n = 3a_{n-1}$
is $r - 3 = 0$, so $a_n^{(\mathrm{h})} = \alpha \cdot 3^n$.

\emph{Particular solution.}  Since $f(n) = 4^n$ and $s = 4$ is not a
root, we try $a_n^{(\mathrm{p})} = A \cdot 4^n$.  Substituting:
$A \cdot 4^n = 3A \cdot 4^{n-1} + 4^n$.  Dividing by $4^{n-1}$:
$4A = 3A + 4$, giving $A = 4$.

\emph{General solution.}  $a_n = \alpha \cdot 3^n + 4 \cdot 4^n
= \alpha \cdot 3^n + 4^{n+1}$.
From $a_0 = 1$: $\alpha + 4 = 1$, so $\alpha = -3$.
Thus $a_n = -3^{n+1} + 4^{n+1}$.
\end{example}

% ------------------------------------------------------------
\section{The Fibonacci Sequence and Binet's Formula}\label{sec:fibonacci}
\index{Fibonacci sequence}\index{Binet's formula}
% ------------------------------------------------------------

\begin{definition}[Fibonacci sequence]\label{def:fibonacci}
The \emph{Fibonacci sequence} $(F_n)_{n \ge 0}$ is defined by
\[
   F_0 = 0, \qquad F_1 = 1, \qquad F_n = F_{n-1} + F_{n-2}
   \quad (n \ge 2).
\]
\end{definition}

Its first terms are $0, 1, 1, 2, 3, 5, 8, 13, 21, 34, 55, \dots$

\begin{theorem}[Binet's formula]\label{thm:binet}
\index{Binet's formula}
For every $n \ge 0$,
\[
   F_n = \frac{\varphi^n - \psi^n}{\sqrt{5}},
\]
where $\varphi = \dfrac{1+\sqrt{5}}{2}$ (the \emph{golden ratio}\index{golden ratio})
and $\psi = \dfrac{1-\sqrt{5}}{2}$.
\end{theorem}

\begin{proof}
The characteristic equation of $F_n = F_{n-1} + F_{n-2}$ is
$r^2 - r - 1 = 0$, whose roots are
\[
   \varphi = \frac{1+\sqrt{5}}{2}, \qquad
   \psi   = \frac{1-\sqrt{5}}{2}.
\]
By Theorem~\ref{thm:LHCC-order2} the general solution is
$F_n = \alpha\, \varphi^n + \beta\, \psi^n$.
The initial conditions give the system
\[
   \alpha + \beta = 0, \qquad
   \alpha\, \varphi + \beta\, \psi = 1.
\]
From the first equation $\beta = -\alpha$.  Substituting into the second:
\[
   \alpha(\varphi - \psi) = 1
   \implies \alpha = \frac{1}{\varphi - \psi} = \frac{1}{\sqrt{5}}.
\]
Therefore $\beta = -1/\sqrt{5}$, and
$F_n = \dfrac{\varphi^n - \psi^n}{\sqrt{5}}$.
\end{proof}

\begin{remark}\label{rem:binet-approx}
Since $\abs{\psi} = \abs{(1-\sqrt{5})/2} \approx 0.618 < 1$,
the term $\psi^n / \sqrt{5} \to 0$ as $n \to \infty$.
Hence $F_n$ is the nearest integer to
$\varphi^n / \sqrt{5}$ for all $n \ge 0$.
\end{remark}

% ------------------------------------------------------------
\section{Classic Recurrences}\label{sec:classic-rec}
% ------------------------------------------------------------

\subsection{The Tower of Hanoi}\label{ssec:hanoi}
\index{Tower of Hanoi}

The Tower of Hanoi puzzle asks for the minimum number of moves $T_n$
needed to transfer $n$ disks from one peg to another, moving one disk at
a time and never placing a larger disk on a smaller one.

\begin{proposition}[Tower of Hanoi recurrence]\label{prop:hanoi}
The sequence $(T_n)$ satisfies
\[
   T_0 = 0, \qquad T_n = 2\, T_{n-1} + 1 \quad (n \ge 1).
\]
Its closed form is $T_n = 2^n - 1$.
\end{proposition}

\begin{proof}
To move $n$ disks from peg~$A$ to peg~$C$:
\begin{enumerate}
\item Move the top $n-1$ disks from $A$ to $B$ ($T_{n-1}$ moves).
\item Move disk $n$ from $A$ to $C$ ($1$ move).
\item Move the $n-1$ disks from $B$ to $C$ ($T_{n-1}$ moves).
\end{enumerate}
This gives $T_n = 2T_{n-1} + 1$.

For the closed form, observe that this is a non-homogeneous first-order
recurrence.  Setting $S_n = T_n + 1$ yields $S_n = 2S_{n-1}$ with
$S_0 = 1$, so $S_n = 2^n$ and $T_n = 2^n - 1$.
\end{proof}

\subsection{Catalan numbers}\label{ssec:catalan}
\index{Catalan numbers}

\begin{definition}[Catalan numbers]\label{def:catalan}
The $n$-th \emph{Catalan number} is
\[
   C_n = \frac{1}{n+1}\binom{2n}{n}, \qquad n \ge 0.
\]
\end{definition}

The first values are $1, 1, 2, 5, 14, 42, 132, 429, \dots$

\begin{proposition}[Catalan recurrence]\label{prop:catalan-rec}
The Catalan numbers satisfy
\[
   C_0 = 1, \qquad
   C_{n+1} = \sum_{i=0}^{n} C_i\, C_{n-i} \quad (n \ge 0).
\]
Equivalently, $C_{n+1} = \dfrac{2(2n+1)}{n+2}\, C_n$.
\end{proposition}

\begin{proof}
Consider the number of ways to fully parenthesise a product of $n+2$
factors.  Splitting after the $i$-th factor (for $i = 1, \dots, n+1$)
gives $C_i \cdot C_{n+1-i}$ arrangements for the left and right parts.
Summing and reindexing yields the convolution.

For the ratio formula, divide the closed form for $C_{n+1}$ by that
for $C_n$:
\[
   \frac{C_{n+1}}{C_n}
   = \frac{1}{n+2}\binom{2n+2}{n+1}
     \cdot \frac{n+1}{\binom{2n}{n}}
   = \frac{n+1}{n+2}\cdot\frac{(2n+2)!/(n+1)!^2}{(2n)!/n!^2}
   = \frac{2(2n+1)}{n+2}. \qedhere
\]
\end{proof}

\begin{remark}\label{rem:catalan-interpretations}
Catalan numbers count many combinatorial objects, including: the number
of Dyck paths\index{Dyck path} of length $2n$, the number of different
triangulations of a convex $(n+2)$-gon, and the number of full binary
trees with $n+1$ leaves.
\end{remark}

% ------------------------------------------------------------
\section{The Master Theorem}\label{sec:master}
\index{master theorem}
% ------------------------------------------------------------

Many divide-and-conquer algorithms yield recurrences of the form
$T(n) = a\, T(n/b) + f(n)$.  The master theorem gives their asymptotic
behaviour in three cases.

\begin{theorem}[Master theorem]\label{thm:master}
Let $a \ge 1$, $b > 1$ be constants, and let $f(n)$ be a non-negative
function.  Define $T(n)$ by
\[
   T(n) = a\, T\!\left(\frac{n}{b}\right) + f(n),
\]
with $T(1) = \Theta(1)$.  Let $c = \log_b a$.  Then:
\begin{enumerate}
\item If $f(n) = O(n^{c-\epsilon})$ for some $\epsilon > 0$, then
      $T(n) = \Theta(n^c)$.
\item If $f(n) = \Theta(n^c \log^k n)$ for some $k \ge 0$, then
      $T(n) = \Theta(n^c \log^{k+1} n)$.
\item If $f(n) = \Omega(n^{c+\epsilon})$ for some $\epsilon > 0$,
      and if $a\, f(n/b) \le \delta\, f(n)$ for some $\delta < 1$ and
      all sufficiently large $n$, then $T(n) = \Theta(f(n))$.
\end{enumerate}
\end{theorem}

\begin{example}[Merge sort]\label{ex:merge-sort}
\index{merge sort}
Merge sort satisfies $T(n) = 2T(n/2) + \Theta(n)$.
Here $a = 2$, $b = 2$, $c = \log_2 2 = 1$, and $f(n) = \Theta(n) =
\Theta(n^c)$.  Case~2 with $k = 0$ gives $T(n) = \Theta(n \log n)$.
\end{example}

\begin{example}[Binary search]\label{ex:binary-search}
\index{binary search}
Binary search satisfies $T(n) = T(n/2) + \Theta(1)$.
Here $a = 1$, $b = 2$, $c = \log_2 1 = 0$, and
$f(n) = \Theta(1) = \Theta(n^0)$.  Case~2 with $k = 0$ gives
$T(n) = \Theta(\log n)$.
\end{example}

\begin{example}[Strassen's algorithm]\label{ex:strassen}
\index{Strassen's algorithm}
Strassen's matrix multiplication satisfies $T(n) = 7T(n/2) + \Theta(n^2)$.
Here $c = \log_2 7 \approx 2.807$, and $f(n) = \Theta(n^2) =
O(n^{c - \epsilon})$ with $\epsilon \approx 0.807$.  Case~1 gives
$T(n) = \Theta(n^{\log_2 7})$.
\end{example}

% ------------------------------------------------------------
\section{Solving Recurrences via Generating Functions}\label{sec:rec-gf}
\index{generating function!solving recurrences}
% ------------------------------------------------------------

Generating functions provide a powerful algebraic approach to solving
recurrences, especially when the characteristic-equation method is
cumbersome.

\begin{definition}[Ordinary generating function]\label{def:ogf}
\index{generating function!ordinary}
The \emph{ordinary generating function} (OGF) of a sequence
$(a_n)_{n \ge 0}$ is the formal power series
\[
   A(x) = \sum_{n=0}^{\infty} a_n\, x^n.
\]
\end{definition}

\paragraph{Strategy.}
Given a recurrence for $(a_n)$:
\begin{enumerate}
\item Multiply both sides by $x^n$ and sum over the valid range of $n$.
\item Express each sum in terms of $A(x)$.
\item Solve the resulting equation for $A(x)$.
\item Extract $a_n$ by expanding $A(x)$ in partial fractions or by
      recognising known series.
\end{enumerate}

\begin{example}[Fibonacci via generating functions]\label{ex:fib-gf}
\index{Fibonacci sequence!generating function}
We derive the OGF of the Fibonacci sequence.  Starting from
$F_n = F_{n-1} + F_{n-2}$ for $n \ge 2$, multiply by $x^n$ and sum:
\[
   \sum_{n=2}^{\infty} F_n x^n
   = \sum_{n=2}^{\infty} F_{n-1} x^n
   + \sum_{n=2}^{\infty} F_{n-2} x^n.
\]
The left side is $F(x) - F_0 - F_1 x = F(x) - x$.  The right side is
$x\bigl(F(x) - F_0\bigr) + x^2 F(x) = x\, F(x) + x^2 F(x)$.
Therefore
\[
   F(x) - x = x\, F(x) + x^2 F(x)
   \implies F(x) = \frac{x}{1 - x - x^2}.
\]
Partial-fraction decomposition with roots $\varphi, \psi$ of
$1 - x - x^2 = 0$ recovers Binet's formula.
\end{example}

\begin{example}[Catalan OGF]\label{ex:catalan-gf}
\index{Catalan numbers!generating function}
Let $C(x) = \sum_{n \ge 0} C_n x^n$.  The convolution recurrence
$C_{n+1} = \sum_{i=0}^{n} C_i C_{n-i}$ translates to
\[
   \frac{C(x) - 1}{x} = C(x)^2
   \implies x\, C(x)^2 - C(x) + 1 = 0.
\]
By the quadratic formula,
\[
   C(x) = \frac{1 - \sqrt{1 - 4x}}{2x}.
\]
(We choose the minus sign so that $C(0) = 1$.)  Expanding $(1-4x)^{1/2}$
via the generalised binomial theorem yields the closed form
$C_n = \frac{1}{n+1}\binom{2n}{n}$.
\end{example}

% ------------------------------------------------------------
\section{Exercises}\label{sec:rec-exercises}
% ------------------------------------------------------------

\begin{exercise}\label{exo:rec1}
Solve the recurrence $a_n = 7a_{n-1} - 10a_{n-2}$ with $a_0 = 2$ and
$a_1 = 1$.
\end{exercise}

\begin{exercise}\label{exo:rec2}
Solve $a_n = 6a_{n-1} - 9a_{n-2}$ with $a_0 = 0$ and $a_1 = 3$.
\end{exercise}

\begin{exercise}\label{exo:rec3}
Find a closed-form solution for $a_n = 2a_{n-1} + 3^n$ with $a_0 = 1$.
\end{exercise}

\begin{exercise}\label{exo:rec4}
\index{Tribonacci numbers}
The \emph{Tribonacci numbers} are defined by $T_0 = 0$, $T_1 = 0$,
$T_2 = 1$, and $T_n = T_{n-1} + T_{n-2} + T_{n-3}$ for $n \ge 3$.
Write down the characteristic equation and find its roots numerically.
\end{exercise}

\begin{exercise}\label{exo:rec5}
Use the master theorem to determine the asymptotic behaviour of
$T(n) = 4T(n/2) + n$.
\end{exercise}

\begin{exercise}\label{exo:rec6}
Use the master theorem for $T(n) = 2T(n/4) + \sqrt{n}$.
\end{exercise}

\begin{exercise}\label{exo:rec7}
Find the ordinary generating function for the sequence defined by
$a_n = 3a_{n-1}$, $a_0 = 2$.  Verify by expanding as a power series.
\end{exercise}

\begin{exercise}\label{exo:rec8}
Prove by induction that the Tower of Hanoi requires at least $2^n - 1$
moves (i.e.\ the solution is optimal).
\end{exercise}

\begin{exercise}\label{exo:rec9}
Show that the Catalan numbers satisfy
$C_n = \dfrac{4n-2}{n+1}\, C_{n-1}$ for $n \ge 1$, starting from
$C_0 = 1$.
\end{exercise}

\begin{exercise}\label{exo:rec10}
A staircase has $n$ steps.  You can climb $1$ or $2$ steps at a time.
Let $s_n$ be the number of distinct ways to climb the staircase.
Set up and solve the recurrence for $s_n$.
\end{exercise}

% ------------------------------------------------------------
\section{Chapter Summary}\label{sec:rec-summary}
% ------------------------------------------------------------

\begin{enumerate}
\item A \textbf{linear homogeneous recurrence with constant coefficients}
      is solved via the \textbf{characteristic equation}.  Distinct roots
      give rise to exponential terms; repeated roots introduce polynomial
      factors.
\item A \textbf{non-homogeneous} recurrence is solved by combining the
      general homogeneous solution with a particular solution found by
      undetermined coefficients.
\item \textbf{Binet's formula} gives a closed form for the Fibonacci
      numbers using the golden ratio.
\item Classic recurrences include the \textbf{Tower of Hanoi}
      ($T_n = 2^n - 1$) and the \textbf{Catalan numbers}
      ($C_n = \frac{1}{n+1}\binom{2n}{n}$).
\item The \textbf{master theorem} gives asymptotic solutions to
      divide-and-conquer recurrences $T(n) = aT(n/b) + f(n)$.
\item \textbf{Generating functions} translate recurrences into algebraic
      equations for formal power series, from which closed forms can be
      extracted.
\end{enumerate}

\newpage

% ============================================================
%  CHAPTER 8 — Graph Theory: Basics
% ============================================================
\chapter{Graph Theory --- Basics}\label{ch:graphs}
\index{graph theory}

Graph theory is one of the most widely applicable branches of discrete
mathematics, with uses ranging from network design and social-network
analysis to chemistry, linguistics, and theoretical computer science.
We begin with the historical problem that sparked the subject, then
develop the fundamental definitions and results.

% ------------------------------------------------------------
\section{The K\"onigsberg Bridge Problem}\label{sec:konigsberg}
\index{K\"onigsberg bridges}
% ------------------------------------------------------------

In 1736, Leonhard Euler\index{Euler, Leonhard} considered whether one
could walk through the city of K\"onigsberg (now Kaliningrad) crossing
each of its seven bridges exactly once and return to the starting point.
He proved that no such walk exists by modelling the landmasses as
\emph{vertices} and the bridges as \emph{edges}, creating what is
regarded as the first theorem in graph theory.

\begin{center}
\begin{tikzpicture}[
  vertex/.style={circle, draw, fill=black!15, minimum size=6mm, inner sep=0pt},
  every edge/.style={draw, thick}
]
  \node[vertex] (A) at (0,0)  {$A$};
  \node[vertex] (B) at (3,1.5)  {$B$};
  \node[vertex] (C) at (3,-1.5) {$C$};
  \node[vertex] (D) at (6,0)  {$D$};

  \draw[thick] (A) to[bend left=20]  (B);
  \draw[thick] (A) to[bend right=20] (B);
  \draw[thick] (A) to[bend left=20]  (C);
  \draw[thick] (A) to[bend right=20] (C);
  \draw[thick] (A) -- (D);
  \draw[thick] (B) -- (D);
  \draw[thick] (C) -- (D);
\end{tikzpicture}
\end{center}

% ------------------------------------------------------------
\section{Fundamental Definitions}\label{sec:graph-def}
% ------------------------------------------------------------

\begin{definition}[Graph]\label{def:graph}
\index{graph!definition}
A \emph{graph} $G = (V, E)$ consists of a finite non-empty set $V$ of
\emph{vertices}\index{vertex} (also called \emph{nodes}) and a set~$E$
of \emph{edges}\index{edge}, where each edge is associated with an
unordered pair of vertices.
\end{definition}

\begin{definition}[Adjacency and incidence]\label{def:adjacency}
\index{adjacency}\index{incidence}
Two vertices $u, v \in V$ are \emph{adjacent} if there is an edge
$e = \{u,v\} \in E$.  In this case $e$ is \emph{incident} with both $u$
and~$v$, and $u$ and $v$ are called the \emph{endpoints} of~$e$.
\end{definition}

\begin{definition}[Loop and multi-edge]\label{def:loop-multiedge}
\index{loop}\index{multi-edge}
An edge whose two endpoints coincide is called a \emph{loop}.  Two or
more edges with the same pair of endpoints are called \emph{multi-edges}
(or \emph{parallel edges}).
\end{definition}

% ------------------------------------------------------------
\section{Types of Graphs}\label{sec:graph-types}
\index{graph!types}
% ------------------------------------------------------------

\begin{definition}[Simple graph]\label{def:simple-graph}
\index{graph!simple}
A \emph{simple graph} is a graph with no loops and no multi-edges.
Every edge is uniquely determined by its pair of distinct endpoints, so
we may identify $E$ with a subset of $\binom{V}{2}$.
\end{definition}

\begin{definition}[Multigraph]\label{def:multigraph}
\index{multigraph}
A \emph{multigraph} allows multi-edges (but usually not loops).  The
K\"onigsberg graph above is a multigraph.
\end{definition}

\begin{definition}[Directed graph]\label{def:digraph}
\index{directed graph}\index{digraph}
A \emph{directed graph} (or \emph{digraph}) $G = (V, A)$ consists of a
set $V$ of vertices and a set $A$ of \emph{arcs}\index{arc} (directed
edges), where each arc is an ordered pair $(u, v)$ with $u$ the
\emph{tail} and $v$ the \emph{head}.
\end{definition}

\begin{center}
\begin{tikzpicture}[
  vertex/.style={circle, draw, fill=black!15, minimum size=6mm, inner sep=0pt},
  >={stealth}
]
  \node[vertex] (u) at (0,0) {$u$};
  \node[vertex] (v) at (2.5,0) {$v$};
  \node[vertex] (w) at (1.25,1.8) {$w$};

  \draw[->, thick] (u) -- (v);
  \draw[->, thick] (v) -- (w);
  \draw[->, thick] (w) -- (u);
\end{tikzpicture}
\captionof{figure}{A directed graph on three vertices.}\label{fig:digraph-example}
\end{center}

\begin{definition}[Weighted graph]\label{def:weighted-graph}
\index{graph!weighted}
A \emph{weighted graph} is a graph in which each edge $e$ is assigned a
numerical value $w(e)$, called its \emph{weight}.
\end{definition}

\begin{center}
\begin{tikzpicture}[
  vertex/.style={circle, draw, fill=black!15, minimum size=6mm, inner sep=0pt}
]
  \node[vertex] (a) at (0,0)  {$a$};
  \node[vertex] (b) at (3,0)  {$b$};
  \node[vertex] (c) at (1.5,2) {$c$};

  \draw[thick] (a) -- node[below] {$5$} (b);
  \draw[thick] (b) -- node[right] {$3$} (c);
  \draw[thick] (c) -- node[left]  {$7$} (a);
\end{tikzpicture}
\captionof{figure}{A weighted graph.}\label{fig:weighted-example}
\end{center}

% ------------------------------------------------------------
\section{Graph Representations}\label{sec:graph-rep}
\index{graph!representation}
% ------------------------------------------------------------

\subsection{Adjacency matrix}\label{ssec:adj-matrix}
\index{adjacency matrix}

\begin{definition}[Adjacency matrix]\label{def:adj-matrix}
Let $G$ be a simple graph with vertices $v_1, \dots, v_n$.  Its
\emph{adjacency matrix} is the $n \times n$ matrix $\mathbf{A}$ with
\[
   A_{ij} =
   \begin{cases}
     1 & \text{if } \{v_i, v_j\} \in E, \\
     0 & \text{otherwise.}
   \end{cases}
\]
For a simple undirected graph, $\mathbf{A}$ is symmetric with zeros on
the diagonal.
\end{definition}

\begin{example}[Adjacency matrix]\label{ex:adj-matrix}
Consider the following graph and its adjacency matrix.

\begin{minipage}{0.45\textwidth}
\centering
\begin{tikzpicture}[
  vertex/.style={circle, draw, fill=black!15, minimum size=6mm, inner sep=0pt}
]
  \node[vertex] (1) at (0,1.5)  {$1$};
  \node[vertex] (2) at (2,1.5)  {$2$};
  \node[vertex] (3) at (2,0)    {$3$};
  \node[vertex] (4) at (0,0)    {$4$};

  \draw[thick] (1) -- (2);
  \draw[thick] (1) -- (4);
  \draw[thick] (2) -- (3);
  \draw[thick] (3) -- (4);
  \draw[thick] (2) -- (4);
\end{tikzpicture}
\end{minipage}%
\begin{minipage}{0.45\textwidth}
\[
   \mathbf{A} =
   \begin{pmatrix}
     0 & 1 & 0 & 1 \\
     1 & 0 & 1 & 1 \\
     0 & 1 & 0 & 1 \\
     1 & 1 & 1 & 0
   \end{pmatrix}
\]
\end{minipage}
\end{example}

\subsection{Incidence matrix}\label{ssec:inc-matrix}
\index{incidence matrix}

\begin{definition}[Incidence matrix]\label{def:inc-matrix}
Let $G$ have $n$ vertices and $m$ edges.  The \emph{incidence matrix} is
the $n \times m$ matrix $\mathbf{B}$ where
\[
   B_{ij} =
   \begin{cases}
     1 & \text{if vertex $v_i$ is an endpoint of edge $e_j$}, \\
     0 & \text{otherwise.}
   \end{cases}
\]
Each column of $\mathbf{B}$ has exactly two ones (for a simple graph
without loops).
\end{definition}

\subsection{Adjacency list}\label{ssec:adj-list}
\index{adjacency list}

\begin{definition}[Adjacency list]\label{def:adj-list}
An \emph{adjacency list} representation stores, for each vertex $v$, the
list of all vertices adjacent to~$v$.  This representation is often more
space-efficient than an adjacency matrix for sparse graphs.
\end{definition}

\begin{example}[Adjacency list]\label{ex:adj-list}
For the graph in Example~\ref{ex:adj-matrix}:
\[
\begin{aligned}
   1 &: \{2, 4\}, \\
   2 &: \{1, 3, 4\}, \\
   3 &: \{2, 4\}, \\
   4 &: \{1, 2, 3\}.
\end{aligned}
\]
\end{example}

\begin{remark}\label{rem:representation-complexity}
For a graph with $n$ vertices and $m$ edges, the adjacency matrix
uses $\Theta(n^2)$ space, while the adjacency list uses $\Theta(n + m)$
space.  Checking whether a specific edge exists takes $O(1)$ time with
the matrix but $O(\deg v)$ time with the list.
\end{remark}

% ------------------------------------------------------------
\section{Vertex Degree and the Handshaking Lemma}\label{sec:degree}
% ------------------------------------------------------------

\begin{definition}[Degree]\label{def:degree}
\index{degree (of a vertex)}
The \emph{degree} of a vertex $v$ in a graph $G$, denoted $\deg(v)$, is
the number of edges incident with $v$ (with loops counted twice).  A
vertex of degree~$0$ is called \emph{isolated}\index{isolated vertex}.
\end{definition}

\begin{theorem}[Handshaking lemma]\label{thm:handshaking}
\index{handshaking lemma}
In any graph $G = (V, E)$,
\begin{equation}\label{eq:handshaking}
   \sum_{v \in V} \deg(v) = 2\, \card{E}.
\end{equation}
\end{theorem}

\begin{proof}
Each edge $e = \{u, v\}$ contributes exactly $1$ to $\deg(u)$ and $1$
to $\deg(v)$, hence $2$ to the total sum.  Summing over all edges gives
the result.
\end{proof}

\begin{corollary}\label{cor:even-odd-vertices}
In any graph, the number of vertices with odd degree is even.
\end{corollary}

\begin{proof}
Let $V_{\mathrm{odd}}$ and $V_{\mathrm{even}}$ denote the sets of
vertices of odd and even degree, respectively.  By
Theorem~\ref{thm:handshaking},
\[
   \sum_{v \in V_{\mathrm{odd}}} \deg(v)
   + \sum_{v \in V_{\mathrm{even}}} \deg(v) = 2\card{E}.
\]
The second sum is even (each term is even), and $2\card{E}$ is even, so
the first sum must be even.  Since each term in the first sum is odd,
the number of terms must be even.
\end{proof}

\begin{definition}[In-degree and out-degree]\label{def:in-out-degree}
\index{in-degree}\index{out-degree}
In a digraph, the \emph{in-degree} $\deg^-(v)$ of vertex $v$ is the
number of arcs with head~$v$, and the \emph{out-degree} $\deg^+(v)$ is
the number of arcs with tail~$v$.  For every digraph,
\[
   \sum_{v \in V} \deg^-(v) = \sum_{v \in V} \deg^+(v) = \card{A}.
\]
\end{definition}

% ------------------------------------------------------------
\section{Walks, Paths, and Cycles}\label{sec:paths-cycles}
% ------------------------------------------------------------

\begin{definition}[Walk]\label{def:walk}
\index{walk}
A \emph{walk} in a graph $G$ is a sequence of vertices
$v_0, v_1, \dots, v_k$ such that $\{v_{i-1}, v_i\} \in E$ for each
$i = 1, \dots, k$.  The number $k$ is the \emph{length} of the walk.
\end{definition}

\begin{definition}[Trail and path]\label{def:trail-path}
\index{trail}\index{path}
A walk in which all \emph{edges} are distinct is called a \emph{trail}.
A walk in which all \emph{vertices} are distinct is called a \emph{path}.
(Every path is a trail, but not conversely.)
\end{definition}

\begin{definition}[Cycle]\label{def:cycle}
\index{cycle}
A \emph{cycle} (or \emph{closed path}) is a walk $v_0, v_1, \dots, v_k$
with $v_0 = v_k$, $k \ge 3$, and all vertices $v_0, v_1, \dots, v_{k-1}$
distinct.
\end{definition}

\begin{definition}[Connectivity]\label{def:connectivity}
\index{connectivity}\index{connected graph}
A graph~$G$ is \emph{connected} if for every pair of vertices $u, v$
there exists a path from $u$ to $v$.  A maximal connected subgraph
of~$G$ is called a \emph{connected component}\index{connected component}.
\end{definition}

\begin{proposition}\label{prop:connected-edges}
A connected graph on $n$ vertices has at least $n - 1$ edges.
\end{proposition}

\begin{proof}
We proceed by induction on $n$.  The base case $n = 1$ is immediate
(zero edges suffice).  For the inductive step, let $G$ be connected
with $n \ge 2$ vertices.  Choose a vertex $v$ and consider the graph
$G' = G - v$.  If $G'$ has $c$ connected components, then $G'$ has at
least $n - 1 - c$ edges by induction applied to each component.  Since
$G$ is connected, $v$ must be adjacent to at least one vertex in each
component of~$G'$, contributing at least $c$ additional edges.  Thus
$\card{E(G)} \ge (n - 1 - c) + c = n - 1$.
\end{proof}

\begin{center}
\begin{tikzpicture}[
  vertex/.style={circle, draw, fill=black!15, minimum size=6mm, inner sep=0pt}
]
  % Component 1
  \node[vertex] (a) at (0,0)   {$a$};
  \node[vertex] (b) at (1.5,1) {$b$};
  \node[vertex] (c) at (1.5,-1){$c$};
  \node[vertex] (d) at (3,0)   {$d$};
  \draw[thick] (a) -- (b) -- (d) -- (c) -- (a);
  \draw[thick] (b) -- (c);

  % Component 2
  \node[vertex] (e) at (5,0.5) {$e$};
  \node[vertex] (f) at (6.5,0.5){$f$};
  \draw[thick] (e) -- (f);

  % Component 3
  \node[vertex] (g) at (5,-1)  {$g$};
\end{tikzpicture}
\captionof{figure}{A graph with three connected components.}\label{fig:components}
\end{center}

% ------------------------------------------------------------
\section{Important Families of Graphs}\label{sec:graph-families}
\index{graph!families}
% ------------------------------------------------------------

\subsection{Complete graphs $K_n$}\label{ssec:complete}
\index{complete graph}\index{$K_n$}

\begin{definition}[Complete graph]\label{def:complete-graph}
The \emph{complete graph} $K_n$ is the simple graph on $n$ vertices in
which every pair of distinct vertices is joined by an edge.  It has
$\binom{n}{2} = \frac{n(n-1)}{2}$ edges.
\end{definition}

\begin{center}
\begin{tikzpicture}[
  vertex/.style={circle, draw, fill=black!15, minimum size=5mm,
                 inner sep=0pt, font=\footnotesize}
]
  % K_3
  \begin{scope}[shift={(0,0)}]
    \node[vertex] (a1) at (90:1)  {$1$};
    \node[vertex] (a2) at (210:1) {$2$};
    \node[vertex] (a3) at (330:1) {$3$};
    \draw[thick] (a1) -- (a2) -- (a3) -- (a1);
    \node at (0,-1.6) {$K_3$};
  \end{scope}
  % K_4
  \begin{scope}[shift={(3.5,0)}]
    \node[vertex] (b1) at (45:1)  {$1$};
    \node[vertex] (b2) at (135:1) {$2$};
    \node[vertex] (b3) at (225:1) {$3$};
    \node[vertex] (b4) at (315:1) {$4$};
    \draw[thick] (b1) -- (b2) -- (b3) -- (b4) -- (b1);
    \draw[thick] (b1) -- (b3);
    \draw[thick] (b2) -- (b4);
    \node at (0,-1.6) {$K_4$};
  \end{scope}
  % K_5
  \begin{scope}[shift={(7.5,0)}]
    \foreach \i in {1,...,5}
      \node[vertex] (c\i) at ({90+(\i-1)*72}:1) {$\i$};
    \foreach \i in {1,...,5}
      \foreach \j in {\i,...,5}
        \draw[thick] (c\i) -- (c\j);
    \node at (0,-1.6) {$K_5$};
  \end{scope}
\end{tikzpicture}
\captionof{figure}{Complete graphs $K_3$, $K_4$, and $K_5$.}\label{fig:complete-graphs}
\end{center}

\subsection{Complete bipartite graphs $K_{m,n}$}\label{ssec:bipartite}
\index{bipartite graph}\index{complete bipartite graph}\index{$K_{m,n}$}

\begin{definition}[Bipartite graph]\label{def:bipartite}
A graph $G = (V, E)$ is \emph{bipartite} if $V$ can be partitioned into
two disjoint sets $X$ and $Y$ such that every edge has one endpoint
in~$X$ and one in~$Y$.
\end{definition}

\begin{definition}[Complete bipartite graph]\label{def:complete-bipartite}
The \emph{complete bipartite graph} $K_{m,n}$ has parts of sizes $m$
and $n$, with every vertex in one part adjacent to every vertex in the
other.  It has $mn$ edges.
\end{definition}

\begin{center}
\begin{tikzpicture}[
  vertex/.style={circle, draw, fill=black!15, minimum size=5mm,
                 inner sep=0pt, font=\footnotesize}
]
  % K_{2,3}
  \begin{scope}[shift={(0,0)}]
    \node[vertex] (x1) at (0,1.2)  {$x_1$};
    \node[vertex] (x2) at (0,0)    {$x_2$};
    \node[vertex] (y1) at (2.5,1.8){$y_1$};
    \node[vertex] (y2) at (2.5,0.6){$y_2$};
    \node[vertex] (y3) at (2.5,-0.6){$y_3$};
    \foreach \i in {1,2}
      \foreach \j in {1,2,3}
        \draw[thick] (x\i) -- (y\j);
    \node at (1.25,-1.4) {$K_{2,3}$};
  \end{scope}
  % K_{3,3}
  \begin{scope}[shift={(5.5,0)}]
    \node[vertex] (u1) at (0,1.8)  {$u_1$};
    \node[vertex] (u2) at (0,0.6)  {$u_2$};
    \node[vertex] (u3) at (0,-0.6) {$u_3$};
    \node[vertex] (w1) at (2.5,1.8){$w_1$};
    \node[vertex] (w2) at (2.5,0.6){$w_2$};
    \node[vertex] (w3) at (2.5,-0.6){$w_3$};
    \foreach \i in {1,2,3}
      \foreach \j in {1,2,3}
        \draw[thick] (u\i) -- (w\j);
    \node at (1.25,-1.4) {$K_{3,3}$};
  \end{scope}
\end{tikzpicture}
\captionof{figure}{Complete bipartite graphs $K_{2,3}$ and $K_{3,3}$.}\label{fig:bipartite-graphs}
\end{center}

\subsection{Cycle graphs $C_n$ and path graphs $P_n$}\label{ssec:cycle-path}
\index{cycle graph}\index{$C_n$}\index{path graph}\index{$P_n$}

\begin{definition}[Cycle graph and path graph]\label{def:cycle-path-graph}
The \emph{cycle graph} $C_n$ ($n \ge 3$) is the graph consisting of a
single cycle on $n$ vertices.  The \emph{path graph} $P_n$ ($n \ge 1$)
is obtained from $C_n$ by removing one edge (equivalently, it is a path
with $n$ vertices and $n-1$ edges).
\end{definition}

\begin{center}
\begin{tikzpicture}[
  vertex/.style={circle, draw, fill=black!15, minimum size=5mm,
                 inner sep=0pt, font=\footnotesize}
]
  % C_5
  \begin{scope}[shift={(0,0)}]
    \foreach \i in {1,...,5}
      \node[vertex] (c\i) at ({90+(\i-1)*72}:1) {};
    \draw[thick] (c1) -- (c2) -- (c3) -- (c4) -- (c5) -- (c1);
    \node at (0,-1.6) {$C_5$};
  \end{scope}
  % C_6
  \begin{scope}[shift={(3.5,0)}]
    \foreach \i in {1,...,6}
      \node[vertex] (d\i) at ({90+(\i-1)*60}:1) {};
    \draw[thick] (d1) -- (d2) -- (d3) -- (d4) -- (d5) -- (d6) -- (d1);
    \node at (0,-1.6) {$C_6$};
  \end{scope}
  % P_5
  \begin{scope}[shift={(7.5,0)}]
    \foreach \i in {1,...,5}
      \node[vertex] (p\i) at ({(\i-1)*0.9},0) {};
    \draw[thick] (p1) -- (p2) -- (p3) -- (p4) -- (p5);
    \node at (1.8,-1.6) {$P_5$};
  \end{scope}
\end{tikzpicture}
\captionof{figure}{Cycle graphs $C_5$, $C_6$ and path graph $P_5$.}\label{fig:cycle-path-graphs}
\end{center}

\subsection{The Petersen graph}\label{ssec:petersen}
\index{Petersen graph}

\begin{definition}[Petersen graph]\label{def:petersen}
The \emph{Petersen graph} is a $3$-regular graph on $10$ vertices and
$15$ edges.  It serves as a counterexample to many optimistic conjectures
in graph theory.
\end{definition}

\begin{center}
\begin{tikzpicture}[
  vertex/.style={circle, draw, fill=black!15, minimum size=5mm,
                 inner sep=0pt}
]
  % Outer cycle
  \foreach \i in {1,...,5}
    \node[vertex] (o\i) at ({90+(\i-1)*72}:2.2) {};
  \draw[thick] (o1) -- (o2) -- (o3) -- (o4) -- (o5) -- (o1);

  % Inner pentagram
  \foreach \i in {1,...,5}
    \node[vertex] (in\i) at ({90+(\i-1)*72}:1) {};
  \draw[thick] (in1) -- (in3) -- (in5) -- (in2) -- (in4) -- (in1);

  % Spokes
  \foreach \i in {1,...,5}
    \draw[thick] (o\i) -- (in\i);
\end{tikzpicture}
\captionof{figure}{The Petersen graph.}\label{fig:petersen}
\end{center}

\begin{proposition}\label{prop:petersen-properties}
The Petersen graph is
\begin{enumerate}
\item $3$-regular (every vertex has degree~$3$),
\item connected,
\item has girth~$5$ (the shortest cycle has length~$5$),
\item has diameter~$2$ (every two vertices are at distance at most~$2$),
\item not Hamiltonian\index{Hamiltonian} (there is no cycle visiting
      every vertex exactly once).
\end{enumerate}
\end{proposition}

% ------------------------------------------------------------
\section{More Graph Illustrations}\label{sec:graph-illustrations}
% ------------------------------------------------------------

Visualising graphs is essential for building intuition.  Here we collect
further examples.

\begin{example}[A graph and its complement]\label{ex:complement}
\index{complement (of a graph)}
The \emph{complement} $\overline{G}$ of a simple graph $G = (V,E)$ is
the graph on the same vertex set where $\{u,v\} \in E(\overline{G})$ if
and only if $\{u,v\} \notin E(G)$.

\begin{center}
\begin{tikzpicture}[
  vertex/.style={circle, draw, fill=black!15, minimum size=5mm,
                 inner sep=0pt, font=\footnotesize}
]
  % G
  \begin{scope}[shift={(0,0)}]
    \node[vertex] (a) at (0,1.5)   {$1$};
    \node[vertex] (b) at (1.5,1.5) {$2$};
    \node[vertex] (c) at (1.5,0)   {$3$};
    \node[vertex] (d) at (0,0)     {$4$};
    \draw[thick] (a) -- (b);
    \draw[thick] (b) -- (c);
    \draw[thick] (c) -- (d);
    \node at (0.75,-0.8) {$G$};
  \end{scope}
  % Complement
  \begin{scope}[shift={(4.5,0)}]
    \node[vertex] (a2) at (0,1.5)   {$1$};
    \node[vertex] (b2) at (1.5,1.5) {$2$};
    \node[vertex] (c2) at (1.5,0)   {$3$};
    \node[vertex] (d2) at (0,0)     {$4$};
    \draw[thick] (a2) -- (c2);
    \draw[thick] (a2) -- (d2);
    \draw[thick] (b2) -- (d2);
    \node at (0.75,-0.8) {$\overline{G}$};
  \end{scope}
\end{tikzpicture}
\end{center}
Together, $G$ and $\overline{G}$ have all $\binom{4}{2} = 6$ possible
edges.
\end{example}

\begin{example}[Regular graphs]\label{ex:regular}
\index{regular graph}
A graph is \emph{$k$-regular} if every vertex has degree~$k$.

\begin{center}
\begin{tikzpicture}[
  vertex/.style={circle, draw, fill=black!15, minimum size=5mm,
                 inner sep=0pt}
]
  % 2-regular (C_6)
  \begin{scope}[shift={(0,0)}]
    \foreach \i in {1,...,6}
      \node[vertex] (a\i) at ({90+(\i-1)*60}:1.1) {};
    \draw[thick] (a1)--(a2)--(a3)--(a4)--(a5)--(a6)--(a1);
    \node at (0,-1.7) {$2$-regular};
  \end{scope}
  % 3-regular (K_{3,3})
  \begin{scope}[shift={(4,0)}]
    \foreach \i in {1,...,3}{
      \node[vertex] (b\i) at ({(\i-1)*1.0 - 1.0}, 1) {};
      \node[vertex] (c\i) at ({(\i-1)*1.0 - 1.0},-1) {};
    }
    \foreach \i in {1,...,3}
      \foreach \j in {1,...,3}
        \draw[thick] (b\i) -- (c\j);
    \node at (0,-1.7) {$3$-regular};
  \end{scope}
\end{tikzpicture}
\end{center}
\end{example}

\begin{example}[The cube graph $Q_3$]\label{ex:cube-graph}
\index{cube graph}\index{$Q_3$}
The \emph{$3$-dimensional cube graph} (or hypercube) $Q_3$ has $8$
vertices labelled by binary strings of length~$3$, with two vertices
adjacent if and only if their labels differ in exactly one bit.

\begin{center}
\begin{tikzpicture}[
  vertex/.style={circle, draw, fill=black!15, minimum size=5mm,
                 inner sep=0pt, font=\tiny},
  scale=1.3
]
  % Front face
  \node[vertex] (000) at (0,0)   {000};
  \node[vertex] (001) at (2,0)   {001};
  \node[vertex] (011) at (2,1.8) {011};
  \node[vertex] (010) at (0,1.8) {010};
  \draw[thick] (000) -- (001) -- (011) -- (010) -- (000);

  % Back face
  \node[vertex] (100) at (0.8,0.7)   {100};
  \node[vertex] (101) at (2.8,0.7)   {101};
  \node[vertex] (111) at (2.8,2.5)   {111};
  \node[vertex] (110) at (0.8,2.5)   {110};
  \draw[thick] (100) -- (101) -- (111) -- (110) -- (100);

  % Connecting edges
  \draw[thick] (000) -- (100);
  \draw[thick] (001) -- (101);
  \draw[thick] (010) -- (110);
  \draw[thick] (011) -- (111);
\end{tikzpicture}
\captionof{figure}{The cube graph $Q_3$.}\label{fig:cube-graph}
\end{center}
\end{example}

\begin{example}[A tree]\label{ex:tree-illustration}
\index{tree}
A \emph{tree} is a connected acyclic graph.  Every tree on $n$ vertices
has exactly $n - 1$ edges.

\begin{center}
\begin{tikzpicture}[
  vertex/.style={circle, draw, fill=black!15, minimum size=5mm,
                 inner sep=0pt, font=\footnotesize},
  level distance=1.2cm,
  sibling distance=1.5cm
]
  \node[vertex] {$r$}
    child { node[vertex] {$a$}
      child { node[vertex] {$d$} }
      child { node[vertex] {$e$} }
    }
    child { node[vertex] {$b$} }
    child { node[vertex] {$c$}
      child { node[vertex] {$f$} }
    };
\end{tikzpicture}
\captionof{figure}{A rooted tree with root $r$.}\label{fig:tree-example}
\end{center}
\end{example}

% ------------------------------------------------------------
\section{Subgraphs and Graph Isomorphism}\label{sec:subgraphs-iso}
% ------------------------------------------------------------

\begin{definition}[Subgraph]\label{def:subgraph}
\index{subgraph}
A graph $H = (V', E')$ is a \emph{subgraph} of $G = (V, E)$ if
$V' \subseteq V$ and $E' \subseteq E$.  If $V' = V$, then $H$ is a
\emph{spanning subgraph}\index{spanning subgraph}.
\end{definition}

\begin{definition}[Induced subgraph]\label{def:induced-subgraph}
\index{induced subgraph}
Given $S \subseteq V$, the subgraph \emph{induced} by $S$, denoted
$G[S]$, has vertex set $S$ and includes all edges of $G$ whose both
endpoints lie in~$S$.
\end{definition}

\begin{definition}[Graph isomorphism]\label{def:isomorphism}
\index{graph!isomorphism}
Two graphs $G_1 = (V_1, E_1)$ and $G_2 = (V_2, E_2)$ are
\emph{isomorphic}, written $G_1 \cong G_2$, if there exists a bijection
$\phi \colon V_1 \to V_2$ such that
\[
   \{u, v\} \in E_1 \iff \{\phi(u), \phi(v)\} \in E_2.
\]
\end{definition}

\begin{remark}\label{rem:iso-invariants}
Isomorphic graphs share all structural properties: number of vertices,
number of edges, degree sequence, connectivity, etc.  These serve as
\emph{isomorphism invariants}\index{isomorphism invariant}: if any
invariant differs, the graphs are not isomorphic.  However, matching
invariants does not guarantee isomorphism.
\end{remark}

% ------------------------------------------------------------
\section{Exercises}\label{sec:graph-exercises}
% ------------------------------------------------------------

\begin{exercise}\label{exo:graph1}
Draw the graphs $K_6$ and $K_{2,4}$.  How many edges does each have?
\end{exercise}

\begin{exercise}\label{exo:graph2}
Write the adjacency matrix and an adjacency list for the cycle graph
$C_5$ (label the vertices $1, 2, 3, 4, 5$).
\end{exercise}

\begin{exercise}\label{exo:graph3}
A graph has $6$ vertices with degree sequence $(3, 3, 3, 3, 2, 2)$.
How many edges does it have?  Draw one such graph.
\end{exercise}

\begin{exercise}\label{exo:graph4}
Prove that a graph is bipartite if and only if it contains no odd cycle.
\end{exercise}

\begin{exercise}\label{exo:graph5}
Show that $K_n$ has $\frac{n!}{2n}$ distinct Hamiltonian cycles for
$n \ge 3$.
\end{exercise}

\begin{exercise}\label{exo:graph6}
For which values of $n$ is the cycle graph $C_n$ bipartite?
\end{exercise}

\begin{exercise}\label{exo:graph7}
\index{self-complementary graph}
A graph $G$ is \emph{self-complementary} if $G \cong \overline{G}$.
Show that a self-complementary graph on $n$ vertices has exactly
$\binom{n}{2}/2$ edges.  For which values of $n$ is this an integer?
\end{exercise}

\begin{exercise}\label{exo:graph8}
Let $G$ be a simple graph on $n$ vertices.  Show that if $G$ has more
than $\binom{n-1}{2}$ edges, then $G$ is connected.
\end{exercise}

\begin{exercise}\label{exo:graph9}
Determine the number of vertices, edges, and the degree sequence of the
Petersen graph directly from its definition.  Verify the handshaking lemma.
\end{exercise}

\begin{exercise}\label{exo:graph10}
The \emph{$n$-dimensional hypercube} $Q_n$\index{hypercube graph}
has $2^n$ vertices labelled by binary strings of length~$n$, with two
vertices adjacent when their labels differ in exactly one position.
\begin{enumerate}
\item How many edges does $Q_n$ have?
\item What is the degree of each vertex?
\item Show that $Q_n$ is bipartite.
\end{enumerate}
\end{exercise}

\begin{exercise}\label{exo:graph11}
Draw the Petersen graph and identify a Hamiltonian path (visiting all
$10$ vertices exactly once).  Explain why no Hamiltonian \emph{cycle}
exists.
\end{exercise}

\begin{exercise}\label{exo:graph12}
Let $G$ be a connected graph with exactly two vertices of odd degree.
Show that $G$ has an Eulerian trail\index{Eulerian trail} (a trail using
every edge exactly once) whose endpoints are the two odd-degree vertices.
\end{exercise}

% ------------------------------------------------------------
\section{Chapter Summary}\label{sec:graph-summary}
% ------------------------------------------------------------

\begin{enumerate}
\item A \textbf{graph} $G = (V, E)$ consists of vertices and edges.
      Variations include simple graphs, multigraphs, digraphs, and
      weighted graphs.
\item Graphs can be represented by an \textbf{adjacency matrix},
      \textbf{incidence matrix}, or \textbf{adjacency list}, each with
      different space and time trade-offs.
\item The \textbf{handshaking lemma} states that
      $\sum_{v} \deg(v) = 2\card{E}$, implying that the number of
      odd-degree vertices is always even.
\item A \textbf{walk} is a sequence of adjacent vertices; a \textbf{path}
      has no repeated vertices; a \textbf{cycle} is a closed path.
      A graph is \textbf{connected} if every pair of vertices is linked
      by a path.
\item Important graph families include the complete graphs $K_n$, the
      complete bipartite graphs $K_{m,n}$, the cycle graphs $C_n$, the
      path graphs $P_n$, and the \textbf{Petersen graph}.
\item \textbf{Isomorphic} graphs are structurally identical; invariants
      such as the degree sequence help distinguish non-isomorphic graphs.
\end{enumerate}

% === END OF CHUNK ch7-8 ===
% === START OF CHUNK ch9-11+end ===

%%%%%%%%%%%%%%%%%%%%%%%%%%%%%%%%%%%%%%%%%%%%%%%%%%%%%%%%%%%%%
%  CHAPTER 9 — Trees, Eulerian and Hamiltonian Graphs
%%%%%%%%%%%%%%%%%%%%%%%%%%%%%%%%%%%%%%%%%%%%%%%%%%%%%%%%%%%%%
\chapter{Trees, Eulerian and Hamiltonian Graphs}
\label{ch:trees-euler-hamilton}
\index{tree|(}

Trees are among the simplest yet most important structures in graph
theory and computer science.  They appear naturally in hierarchical
data, search algorithms, network design, and combinatorial enumeration.
In this chapter we characterize trees, study spanning trees and
minimum-weight spanning trees, explore Eulerian and Hamiltonian
graphs, and present the classical theorems of Dirac and Ore.

%------------------------------------------------------------------
\section{Trees: Definitions and Characterization}
\label{sec:tree-def}

\begin{definition}[Tree]
\label{def:tree}
\index{tree!definition}
A \emph{tree} is a connected acyclic graph.  A \emph{forest} is an
acyclic graph (not necessarily connected); each connected component
of a forest is a tree.
\end{definition}

\begin{definition}[Leaf]
\label{def:leaf}
\index{leaf}
A vertex of degree~$1$ in a tree is called a \emph{leaf}
(or \emph{pendant vertex}).
\end{definition}

\begin{lemma}
\label{lem:tree-leaf}
Every tree with at least two vertices has at least two leaves.
\end{lemma}

\begin{proof}
Let $T$ be a tree on $n\ge 2$ vertices.  Take a longest path
$v_0,v_1,\dots,v_k$ in~$T$.  Since $T$ is acyclic, $v_0$ cannot be
adjacent to any $v_i$ with $i\ge 2$ (otherwise a cycle would form),
and by maximality of the path $v_0$ has no neighbor outside the path.
Hence $\deg(v_0)=1$.  The same argument applies to~$v_k$.
\end{proof}

\begin{theorem}[Characterization of trees]
\label{thm:tree-char}
\index{tree!characterization}
Let $G$ be a graph on $n$ vertices.  The following are equivalent:
\begin{enumerate}
  \item $G$ is a tree (connected and acyclic).
  \item $G$ is connected and has exactly $n-1$ edges.
  \item $G$ is acyclic and has exactly $n-1$ edges.
  \item For every pair of vertices $u,v$ there is a unique path
        from $u$ to~$v$.
  \item $G$ is connected, but removing any edge disconnects it.
  \item $G$ is acyclic, but adding any edge creates exactly one cycle.
\end{enumerate}
\end{theorem}

\begin{proof}
We prove the cycle of implications
$(1)\!\Rightarrow\!(2)\!\Rightarrow\!(3)\!\Rightarrow\!(1)$ and then
show $(1)\!\Leftrightarrow\!(4)\!\Leftrightarrow\!(5)
\!\Leftrightarrow\!(6)$.

\medskip\noindent
$(1)\Rightarrow(2)$:\; By induction on~$n$.  For $n=1$ the tree has
$0=n-1$ edges.  Suppose every tree on fewer than $n\ge 2$ vertices has
$n-2$ edges.  By Lemma~\ref{lem:tree-leaf}, $G$ has a leaf~$v$.
Removing $v$ and its incident edge gives a tree $G'$ on $n-1$ vertices,
which by induction has $n-2$ edges.  Hence $G$ has $n-1$ edges.
Since $G$ is a tree it is connected.

\medskip\noindent
$(2)\Rightarrow(3)$:\; Suppose $G$ is connected with $n-1$ edges but
contains a cycle~$C$.  Removing an edge of~$C$ keeps $G$ connected
(both endpoints remain linked via the rest of the cycle) but yields a
connected graph on $n$ vertices with $n-2$ edges.  Repeating the
argument as long as cycles remain we eventually reach a connected
acyclic graph (a tree) on $n$ vertices with fewer than $n-1$ edges,
contradicting $(1)\Rightarrow(2)$.

\medskip\noindent
$(3)\Rightarrow(1)$:\; Suppose $G$ is acyclic with $n-1$ edges.  Let
$G$ have $k$ components $T_1,\dots,T_k$ with $n_i$ vertices each.
Each $T_i$ is a tree (connected and acyclic), so has $n_i-1$ edges.
The total number of edges is $\sum(n_i-1)=n-k$.  But $G$ has $n-1$
edges, so $k=1$ and $G$ is connected.

\medskip\noindent
$(1)\Leftrightarrow(4)$:\; In a connected graph, a path between every
pair exists.  If two distinct paths joined $u$ to~$v$, their union
would contain a cycle, contradicting acyclicity.  Conversely, unique
paths imply connectivity and acyclicity.

\medskip\noindent
$(1)\Leftrightarrow(5)$:\; A tree is connected.  If removing edge
$\{u,v\}$ leaves the graph connected, an alternative $u$--$v$ path
together with the edge $\{u,v\}$ would form a cycle.  Conversely, a
connected graph in which every edge is a bridge is acyclic.

\medskip\noindent
$(1)\Leftrightarrow(6)$:\; A tree is acyclic.  Adding edge $\{u,v\}$
creates a cycle using the unique $u$--$v$ path.  Conversely, an acyclic
graph where every non-edge addition creates exactly one cycle is
connected (if $u,v$ were in different components, adding $\{u,v\}$
would create no cycle).
\end{proof}

\begin{example}[Small trees]
\label{ex:small-trees}
Up to isomorphism, there is one tree on $1,2,3$ vertices, two trees on
$4$ vertices (the path $P_4$ and the star $K_{1,3}$), and three trees
on $5$ vertices.

\begin{center}
\begin{tikzpicture}[every node/.style={circle,draw,inner sep=1.5pt,fill=black},
                    scale=0.85]
  % P4
  \node (a1) at (0,0) {};
  \node (a2) at (0.8,0) {};
  \node (a3) at (1.6,0) {};
  \node (a4) at (2.4,0) {};
  \draw (a1)--(a2)--(a3)--(a4);
  \node[draw=none,fill=none] at (1.2,-0.5) {$P_4$};
  % K_{1,3}
  \node (b0) at (4.2,0) {};
  \node (b1) at (3.4,0.7) {};
  \node (b2) at (4.2,0.9) {};
  \node (b3) at (5.0,0.7) {};
  \draw (b0)--(b1) (b0)--(b2) (b0)--(b3);
  \node[draw=none,fill=none] at (4.2,-0.5) {$K_{1,3}$};
\end{tikzpicture}
\end{center}
\end{example}

%------------------------------------------------------------------
\section{Spanning Trees}
\label{sec:spanning-tree}
\index{spanning tree}

\begin{definition}[Spanning tree]
\label{def:spanning-tree}
A \emph{spanning tree} of a connected graph~$G$ is a subgraph that is
a tree and includes every vertex of~$G$.
\end{definition}

\begin{proposition}
\label{prop:spanning-exists}
Every connected graph has a spanning tree.
\end{proposition}

\begin{proof}
If $G$ is connected and contains a cycle, remove an edge of that cycle;
the resulting graph is still connected.  Repeat until no cycle remains.
The result is a connected acyclic spanning subgraph, hence a spanning tree.
\end{proof}

\subsection{Cayley's Formula}
\label{ssec:cayley}
\index{Cayley's formula}

\begin{theorem}[Cayley, 1889]
\label{thm:cayley}
The number of labeled trees on the vertex set $\{1,2,\dots,n\}$ is $n^{n-2}$.
\end{theorem}

We state without proof the classical result, noting that elegant proofs
exist via Pr\"ufer sequences and the matrix-tree theorem.
\index{Prufer sequence@Pr\"ufer sequence}

\begin{example}
For $n=3$ there are $3^{3-2}=3$ labeled trees on $\{1,2,3\}$: the
edges are $\{1,2\},\{2,3\}$; or $\{1,3\},\{2,3\}$; or $\{1,2\},\{1,3\}$.
\end{example}

%------------------------------------------------------------------
\section{Minimum Spanning Trees}
\label{sec:mst}
\index{minimum spanning tree}

Let $G=(V,E)$ be a connected graph with edge-weight function
$w\colon E\to\R$.  A \emph{minimum spanning tree} (MST) is a spanning
tree $T$ that minimizes $\sum_{e\in E(T)}w(e)$.

\subsection{Kruskal's Algorithm}
\label{ssec:kruskal}
\index{Kruskal's algorithm}

\begin{enumerate}
  \item Sort the edges in non-decreasing order of weight.
  \item Initialize $T=\emptyset$.
  \item For each edge $e$ (in sorted order): add $e$ to $T$ if it does
        not create a cycle.
  \item Return $T$.
\end{enumerate}

\begin{theorem}
\label{thm:kruskal-correct}
Kruskal's algorithm produces a minimum spanning tree.
\end{theorem}

\begin{proof}[Proof sketch]
At each step the lightest available edge that does not create a cycle is
added.  By the cut property (every lightest edge crossing a cut belongs
to some MST), each chosen edge belongs to an MST.  Since we select
exactly $n-1$ edges and form a spanning tree, the result is an MST.
\end{proof}

\subsection{Prim's Algorithm}
\label{ssec:prim}
\index{Prim's algorithm}

\begin{enumerate}
  \item Start with an arbitrary vertex $v_0$; set $S=\{v_0\}$, $T=\emptyset$.
  \item While $S\ne V$: among all edges with one endpoint in $S$ and the
        other in $V\setminus S$, choose one of minimum weight, say $\{u,v\}$
        with $u\in S$.  Add $v$ to $S$ and $\{u,v\}$ to~$T$.
  \item Return $T$.
\end{enumerate}

\begin{theorem}
\label{thm:prim-correct}
Prim's algorithm produces a minimum spanning tree.
\end{theorem}

\begin{example}[MST computation]
\label{ex:mst}
Consider the weighted graph below.  Both Kruskal's and Prim's algorithms
yield the MST with edges $\{a,b\},\{b,c\},\{c,d\}$ of total weight~$6$.

\begin{center}
\begin{tikzpicture}[
  vtx/.style={circle,draw,minimum size=16pt,inner sep=0pt},
  scale=1.1]
  \node[vtx] (a) at (0,0) {$a$};
  \node[vtx] (b) at (2,0) {$b$};
  \node[vtx] (c) at (2,-1.5) {$c$};
  \node[vtx] (d) at (0,-1.5) {$d$};
  \draw (a)--(b) node[midway,above]{1};
  \draw (b)--(c) node[midway,right]{2};
  \draw (c)--(d) node[midway,below]{3};
  \draw (a)--(d) node[midway,left]{4};
  \draw (a)--(c) node[midway,above right]{5};
\end{tikzpicture}
\end{center}
\end{example}

%------------------------------------------------------------------
\section{Eulerian Graphs}
\label{sec:euler}
\index{Eulerian graph|(}

\begin{definition}[Euler circuit and Euler path]
\label{def:euler-circuit}
\index{Euler circuit}
\index{Euler path}
An \emph{Euler circuit} in a graph $G$ is a closed walk that traverses
every edge exactly once.  An \emph{Euler path} is a walk that
traverses every edge exactly once (not necessarily closed).  A graph
possessing an Euler circuit is called \emph{Eulerian}.
\end{definition}

\begin{theorem}[Euler, 1736]
\label{thm:euler-char}
\index{Eulerian graph!characterization}
A connected graph $G$ has an Euler circuit if and only if every vertex
of $G$ has even degree.  It has an Euler path (but no Euler circuit) if
and only if it has exactly two vertices of odd degree.
\end{theorem}

\begin{proof}
$(\Rightarrow)$\; If $G$ has an Euler circuit, each visit to a vertex
uses one edge to enter and one to leave, contributing $2$ to the degree.
Hence every vertex has even degree.

$(\Leftarrow)$\; We proceed by strong induction on $\abs{E(G)}$.  If
$\abs{E(G)}=0$ the empty walk suffices.  Otherwise, since every vertex
has even degree $\ge 2$, $G$ contains a cycle~$C$ (start at any vertex
and walk without repeating a vertex; since every visited vertex has even
degree, we must eventually return).  Remove the edges of~$C$ from~$G$
to obtain $G'$.  Every vertex of $G'$ still has even degree.  Each
connected component of $G'$ with edges is Eulerian by induction.  We
splice: walk along~$C$, and whenever we reach a vertex that starts a
component of $G'$ with edges, we detour through that component's Euler
circuit before continuing along~$C$.  The result is an Euler circuit
of~$G$.

For the Euler path statement: if $G$ has exactly two odd-degree vertices
$u$ and~$v$, add edge $\{u,v\}$ to make all degrees even; find an Euler
circuit; remove $\{u,v\}$ to obtain an Euler path from $u$ to~$v$.
\end{proof}

\subsection{Fleury's Algorithm}
\label{ssec:fleury}
\index{Fleury's algorithm}

Given an Eulerian graph~$G$:
\begin{enumerate}
  \item Start at any vertex.
  \item At each step, choose an edge to traverse.  Prefer a non-bridge
        edge (i.e., do not cross a bridge unless no other option exists).
  \item Remove the traversed edge.  Repeat until all edges are traversed.
\end{enumerate}

\begin{remark}
Fleury's algorithm, while elegant, runs in $O(\abs{E}^2)$ time because
of bridge-checking.  Hierholzer's algorithm achieves $O(\abs{E})$.
\end{remark}

\begin{example}[K\"onigsberg bridges]
\label{ex:konigsberg}
\index{Konigsberg bridges@K\"onigsberg bridges}
The original K\"onigsberg bridge graph has four vertices with degrees
$3,3,3,5$ --- all odd.  Hence it has neither an Euler path nor an Euler
circuit, which was Euler's celebrated 1736 result.

\begin{center}
\begin{tikzpicture}[
  vtx/.style={circle,draw,minimum size=14pt,inner sep=0pt,fill=white},
  scale=1]
  \node[vtx] (A) at (0,1) {$A$};
  \node[vtx] (B) at (2,2) {$B$};
  \node[vtx] (C) at (2,0) {$C$};
  \node[vtx] (D) at (4,1) {$D$};
  \draw (A) to[bend left=20] (B);
  \draw (A) to[bend right=20] (B);
  \draw (A) to[bend left=20] (C);
  \draw (A) to[bend right=20] (C);
  \draw (B)--(D);
  \draw (C)--(D);
  \draw (A)--(D);
\end{tikzpicture}
\end{center}
\end{example}

\index{Eulerian graph|)}

%------------------------------------------------------------------
\section{Hamiltonian Graphs}
\label{sec:hamilton}
\index{Hamiltonian graph|(}

\begin{definition}[Hamilton cycle and Hamilton path]
\label{def:hamilton-cycle}
\index{Hamilton cycle}
\index{Hamilton path}
A \emph{Hamilton cycle} in a graph $G$ is a cycle that visits every
vertex exactly once.  A \emph{Hamilton path} visits every vertex
exactly once.  A graph possessing a Hamilton cycle is called
\emph{Hamiltonian}.
\end{definition}

\begin{remark}
Unlike the Eulerian case, no simple necessary and sufficient condition
for the existence of a Hamilton cycle is known.  In fact, deciding
whether a graph is Hamiltonian is NP-complete.
\index{NP-complete}
\end{remark}

\begin{theorem}[Dirac, 1952]
\label{thm:dirac}
\index{Dirac's theorem}
If $G$ is a simple graph on $n\ge 3$ vertices such that
$\deg(v)\ge n/2$ for every vertex~$v$, then $G$ is Hamiltonian.
\end{theorem}

\begin{proof}
Suppose for contradiction that $G$ satisfies the degree condition but is
not Hamiltonian.  Add edges to $G$ one at a time (choosing non-adjacent
pairs) until obtaining a graph $G^*$ that is not Hamiltonian but adding
any further edge would create a Hamilton cycle.  Such a maximal
non-Hamiltonian graph $G^*$ still satisfies $\deg(v)\ge n/2$.

Let $u,v$ be non-adjacent in $G^*$.  Then $G^*+\{u,v\}$ has a Hamilton
cycle, which must use the edge $\{u,v\}$.  Hence $G^*$ has a Hamilton
path $v=v_1,v_2,\dots,v_n=u$.

Define the sets:
\[
  S=\set{i : v_i \text{ is adjacent to } u},\qquad
  T=\set{i : v_{i+1} \text{ is adjacent to } v}.
\]
Then $\abs{S}\ge n/2$ and $\abs{T}\ge n/2$, and both $S,T\subseteq
\{1,\dots,n-1\}$.  By pigeonhole, there exists $i\in S\cap T$.  Then
$v_i$ is adjacent to $u=v_n$ and $v_{i+1}$ is adjacent to $v=v_1$.
The cycle
\[
  v_1,v_2,\dots,v_i,v_n,v_{n-1},\dots,v_{i+1},v_1
\]
is a Hamilton cycle in $G^*$, a contradiction.
\end{proof}

\begin{theorem}[Ore, 1960]
\label{thm:ore}
\index{Ore's theorem}
If $G$ is a simple graph on $n\ge 3$ vertices such that for every
pair of non-adjacent vertices $u,v$ we have
$\deg(u)+\deg(v)\ge n$, then $G$ is Hamiltonian.
\end{theorem}

\begin{proof}
The proof is virtually identical to Dirac's.  We form the maximal
non-Hamiltonian graph $G^*$ and obtain a Hamilton path
$v_1,v_2,\dots,v_n$ with $v_1,v_n$ non-adjacent.  Ore's condition
gives $\deg(v_1)+\deg(v_n)\ge n$.  The same pigeonhole argument as
above yields a Hamilton cycle, contradicting maximality.
\end{proof}

\begin{corollary}
\label{cor:dirac-from-ore}
Dirac's theorem follows from Ore's theorem.
\end{corollary}

\begin{proof}
If $\deg(v)\ge n/2$ for all $v$, then $\deg(u)+\deg(v)\ge n$ for every
pair $u,v$, so Ore's condition is satisfied.
\end{proof}

\begin{example}[Hamiltonian vs.\ non-Hamiltonian]
\label{ex:hamilton}
The Petersen graph is $3$-regular on $10$ vertices.  Since $3<10/2=5$,
Dirac's theorem does not apply, and indeed the Petersen graph is not
Hamiltonian (though it does have a Hamilton path).

\begin{center}
\begin{tikzpicture}[
  vtx/.style={circle,draw,fill=black,inner sep=1.5pt},
  scale=0.9]
  % outer pentagon
  \foreach \i in {0,...,4}{
    \node[vtx] (o\i) at ({90+72*\i}:2) {};
  }
  % inner pentagram
  \foreach \i in {0,...,4}{
    \node[vtx] (i\i) at ({90+72*\i}:1) {};
  }
  % outer edges
  \foreach \i [evaluate=\i as \j using {int(mod(\i+1,5))}] in {0,...,4}{
    \draw (o\i)--(o\j);
  }
  % inner edges (pentagram)
  \foreach \i [evaluate=\i as \j using {int(mod(\i+2,5))}] in {0,...,4}{
    \draw (i\i)--(i\j);
  }
  % spokes
  \foreach \i in {0,...,4}{
    \draw (o\i)--(i\i);
  }
\end{tikzpicture}
\end{center}
\end{example}

\index{Hamiltonian graph|)}

%------------------------------------------------------------------
\section{Exercises}
\label{sec:ex-ch9}

\begin{exercise}
\label{exo:tree-edges}
Prove that a forest on $n$ vertices and $k$ components has exactly
$n-k$ edges.
\end{exercise}

\begin{exercise}
\label{exo:prufer}
\index{Prufer sequence@Pr\"ufer sequence}
Use Pr\"ufer sequences to show that the complete graph $K_n$ has
$n^{n-2}$ spanning trees.
\end{exercise}

\begin{exercise}
\label{exo:kruskal-run}
Apply Kruskal's algorithm to the complete graph $K_4$ with edge weights
$w(1,2)=3$, $w(1,3)=1$, $w(1,4)=6$, $w(2,3)=4$, $w(2,4)=2$,
$w(3,4)=5$.  List the edges in the order they are added.
\end{exercise}

\begin{exercise}
\label{exo:euler-degree}
Let $G$ be a connected graph with exactly $2k$ vertices of odd degree
($k\ge 1$).  Show that the edges of $G$ can be partitioned into exactly
$k$ edge-disjoint paths (none of which is closed).
\end{exercise}

\begin{exercise}
\label{exo:hamilton-cube}
Show that the $n$-dimensional hypercube graph $Q_n$ is Hamiltonian for
$n\ge 2$.
\end{exercise}

\begin{exercise}
\label{exo:ore-application}
Let $G$ be a simple graph on $7$ vertices such that every vertex has
degree at least $4$.  Prove that $G$ is Hamiltonian using Ore's theorem.
\end{exercise}

\index{tree|)}


%%%%%%%%%%%%%%%%%%%%%%%%%%%%%%%%%%%%%%%%%%%%%%%%%%%%%%%%%%%%%
%  CHAPTER 10 — Coloring, Planarity, Bipartite Graphs
%%%%%%%%%%%%%%%%%%%%%%%%%%%%%%%%%%%%%%%%%%%%%%%%%%%%%%%%%%%%%
\chapter{Coloring, Planarity, and Bipartite Graphs}
\label{ch:coloring-planarity}

Graph coloring, planarity, and bipartiteness are central themes in
graph theory with wide-ranging applications in scheduling, map design,
register allocation, and combinatorial optimization.

%------------------------------------------------------------------
\section{Vertex Coloring}
\label{sec:vertex-coloring}
\index{graph coloring|(}
\index{vertex coloring}

\begin{definition}[Proper coloring, chromatic number]
\label{def:proper-coloring}
\index{chromatic number}
A \emph{proper $k$-coloring} of a graph $G=(V,E)$ is a function
$c\colon V\to\{1,2,\dots,k\}$ such that $c(u)\ne c(v)$ whenever
$\{u,v\}\in E$.  The \emph{chromatic number} $\chi(G)$ is the smallest
$k$ for which $G$ admits a proper $k$-coloring.
\end{definition}

\begin{example}
\label{ex:chi-basic}
\begin{itemize}
  \item $\chi(K_n)=n$.
  \item $\chi(C_{2k})=2$ and $\chi(C_{2k+1})=3$ for $k\ge 1$.
  \item A graph is bipartite if and only if $\chi(G)\le 2$.
\end{itemize}
\end{example}

\subsection{Greedy Coloring}
\label{ssec:greedy-coloring}
\index{greedy coloring}

\begin{proposition}[Greedy bound]
\label{prop:greedy-bound}
For any graph $G$, $\chi(G)\le\Delta(G)+1$, where $\Delta(G)$ is the
maximum degree.
\end{proposition}

\begin{proof}
Order the vertices $v_1,\dots,v_n$ arbitrarily.  Assign to $v_i$ the
smallest color not used by its already-colored neighbors.  Since $v_i$
has at most $\Delta(G)$ neighbors, at most $\Delta(G)$ colors are
forbidden, and a color among $\{1,\dots,\Delta(G)+1\}$ is always
available.
\end{proof}

\begin{theorem}[Brooks, 1941]
\label{thm:brooks}
\index{Brooks' theorem}
If $G$ is a connected graph that is neither a complete graph nor an odd
cycle, then $\chi(G)\le\Delta(G)$.
\end{theorem}

\begin{remark}[Four Color Theorem]
\label{rem:four-color}
\index{Four Color Theorem}
Every planar graph is $4$-colorable.  This was proved by Appel and Haken
(1976) with extensive computer assistance, and later verified by
Robertson, Sanders, Seymour, and Thomas (1997).
\end{remark}

%------------------------------------------------------------------
\section{Chromatic Polynomial}
\label{sec:chromatic-poly}
\index{chromatic polynomial|(}

\begin{definition}[Chromatic polynomial]
\label{def:chromatic-poly}
The \emph{chromatic polynomial} $P(G,k)$ counts the number of proper
$k$-colorings of~$G$.  It is a polynomial in~$k$.
\end{definition}

\begin{example}
\label{ex:chromatic-tree}
For a tree $T$ on $n$ vertices, $P(T,k)=k(k-1)^{n-1}$.
\end{example}

\begin{proof}
Root the tree.  The root has $k$ choices.  Each subsequent vertex
(processed parent before children) has $k-1$ choices (any color except
its parent's).
\end{proof}

\subsection{Deletion--Contraction}
\label{ssec:del-contract}
\index{deletion-contraction}

\begin{theorem}[Deletion--Contraction]
\label{thm:del-contract}
For any edge $e=\{u,v\}$ of~$G$:
\[
  P(G,k)=P(G-e,\,k)-P(G/e,\,k),
\]
where $G-e$ is obtained by deleting~$e$ and $G/e$ by contracting~$e$
(merging $u$ and~$v$ and removing resulting multi-edges).
\end{theorem}

\begin{proof}
The colorings of $G-e$ can be split into two classes: those where
$c(u)\ne c(v)$, which number $P(G,k)$, and those where $c(u)=c(v)$,
which are in bijection with colorings of $G/e$ (the merged vertex
receives the common color).  Hence
$P(G-e,k)=P(G,k)+P(G/e,k)$.
\end{proof}

\begin{example}
\label{ex:chromatic-C4}
For the cycle $C_4$, contracting an edge gives $C_3$ and deleting it
gives $P_4$:
\[
  P(C_4,k)=P(P_4,k)-P(C_3,k)=k(k-1)^3-k(k-1)(k-2)
  =(k-1)^4+(k-1).
\]
\end{example}

\index{chromatic polynomial|)}

%------------------------------------------------------------------
\section{Planar Graphs}
\label{sec:planar}
\index{planar graph|(}

\begin{definition}[Planar graph]
\label{def:planar}
A graph is \emph{planar} if it can be drawn in the plane with no edge
crossings.  Such a drawing is called a \emph{planar embedding}.  The
connected regions of the complement of the drawing are called
\emph{faces}.
\index{face!of planar graph}
\end{definition}

\begin{theorem}[Euler's formula for planar graphs]
\label{thm:euler-formula}
\index{Euler's formula!planar graphs}
If $G$ is a connected planar graph with $V$ vertices, $E$ edges, and
$F$ faces (including the outer face), then
\[
  V-E+F=2.
\]
\end{theorem}

\begin{proof}
By induction on~$E$.  If $E=0$, then $V=1$ and $F=1$ (only the outer
face), so $1-0+1=2$.

If $G$ is a tree, then $E=V-1$ and $F=1$, so $V-(V-1)+1=2$.

If $G$ contains a cycle, let $e$ be an edge on a cycle.  Removing $e$
merges two faces into one, so $F$ decreases by~$1$ while $V$ stays the
same and $E$ decreases by~$1$.  By induction on the smaller graph,
$V-(E-1)+(F-1)=2$, which gives $V-E+F=2$.
\end{proof}

\begin{corollary}
\label{cor:edge-bound}
\index{planar graph!edge bound}
If $G$ is a simple connected planar graph with $V\ge 3$ vertices, then
$E\le 3V-6$.
\end{corollary}

\begin{proof}
Each face is bounded by at least $3$ edges, and each edge borders at
most $2$ faces, so $2E\ge 3F$.  Thus $F\le 2E/3$.  Substituting into
Euler's formula: $2=V-E+F\le V-E+2E/3=V-E/3$, giving $E\le 3V-6$.
\end{proof}

\begin{corollary}
\label{cor:edge-bound-triangle-free}
If $G$ is a simple connected planar graph with $V\ge 3$ and no
triangles, then $E\le 2V-4$.
\end{corollary}

\begin{proof}
Without triangles, each face is bounded by at least $4$ edges, so
$2E\ge 4F$, hence $F\le E/2$.  Then $2=V-E+F\le V-E/2$, giving
$E\le 2V-4$.
\end{proof}

\begin{theorem}[$K_5$ is non-planar]
\label{thm:K5-nonplanar}
\index{K5 non-planar@$K_5$ non-planar}
The complete graph $K_5$ is not planar.
\end{theorem}

\begin{proof}
$K_5$ has $V=5$ and $E=10$.  But $3V-6=9<10=E$, violating
Corollary~\ref{cor:edge-bound}.
\end{proof}

\begin{theorem}[$K_{3,3}$ is non-planar]
\label{thm:K33-nonplanar}
\index{K33 non-planar@$K_{3,3}$ non-planar}
The complete bipartite graph $K_{3,3}$ is not planar.
\end{theorem}

\begin{proof}
$K_{3,3}$ has $V=6$ and $E=9$.  Since $K_{3,3}$ is bipartite it
contains no triangles, so Corollary~\ref{cor:edge-bound-triangle-free}
applies: $E\le 2V-4=8<9=E$, a contradiction.
\end{proof}

\begin{theorem}[Kuratowski, 1930]
\label{thm:kuratowski}
\index{Kuratowski's theorem}
A graph is planar if and only if it contains no subdivision of $K_5$ or
$K_{3,3}$ as a subgraph.
\end{theorem}

\begin{example}[Planar embedding of $K_4$]
\label{ex:planar-K4}

\begin{center}
\begin{tikzpicture}[
  vtx/.style={circle,draw,minimum size=14pt,inner sep=0pt},
  scale=1]
  \node[vtx] (1) at (0,2)   {1};
  \node[vtx] (2) at (-1.5,0){2};
  \node[vtx] (3) at (1.5,0) {3};
  \node[vtx] (4) at (0,0.8) {4};
  \draw (1)--(2)--(3)--(1);
  \draw (1)--(4) (2)--(4) (3)--(4);
  \node at (0,-0.7) {$V=4,\ E=6,\ F=4$};
  \node at (0,-1.2) {$4-6+4=2\;\checkmark$};
\end{tikzpicture}
\end{center}
\end{example}

\index{planar graph|)}

%------------------------------------------------------------------
\section{Bipartite Graphs}
\label{sec:bipartite}
\index{bipartite graph|(}

\begin{definition}[Bipartite graph]
\label{def:bipartite}
A graph $G=(V,E)$ is \emph{bipartite} if $V$ can be partitioned into
two sets $A$ and $B$ such that every edge has one endpoint in $A$ and
one in~$B$.
\end{definition}

\begin{theorem}[Odd-cycle characterization]
\label{thm:bipartite-char}
\index{bipartite graph!odd cycle characterization}
A graph is bipartite if and only if it contains no odd cycle.
\end{theorem}

\begin{proof}
$(\Rightarrow)$\; In a bipartite graph with parts $A,B$, any walk
alternates between $A$ and $B$.  A cycle that starts in $A$ must return
to $A$, requiring an even number of steps.

$(\Leftarrow)$\; Suppose $G$ has no odd cycle.  We may assume $G$ is
connected (handle components independently).  Fix a vertex $s$ and
define
\[
  A=\set{v\in V : d(s,v)\text{ is even}},\qquad
  B=\set{v\in V : d(s,v)\text{ is odd}},
\]
where $d(s,v)$ is the distance from $s$ to~$v$.  We claim every edge
joins a vertex in $A$ to one in~$B$.  Suppose $\{u,v\}\in E$ with
$u,v\in A$ (both at even distance from~$s$).  Then a shortest $s$--$u$
path, the edge $\{u,v\}$, and a shortest $v$--$s$ path form a closed
walk of odd length, which must contain an odd cycle, contradicting
our hypothesis.  A symmetric argument handles $u,v\in B$.
\end{proof}

%------------------------------------------------------------------
\section{Matchings}
\label{sec:matching}
\index{matching|(}

\begin{definition}[Matching]
\label{def:matching}
A \emph{matching} in a graph $G$ is a set of pairwise disjoint edges.
A matching is \emph{perfect} if it covers every vertex.
\end{definition}

\begin{theorem}[K\"onig, 1931]
\label{thm:konig}
\index{Konig's theorem@K\"onig's theorem}
In a bipartite graph, the maximum size of a matching equals the minimum
size of a vertex cover.
\end{theorem}

\begin{theorem}[Hall's marriage theorem]
\label{thm:hall}
\index{Hall's marriage theorem}
Let $G=(A\cup B,E)$ be a bipartite graph.  There exists a matching that
covers every vertex in $A$ if and only if for every $S\subseteq A$,
\[
  \abs{N(S)}\ge\abs{S},
\]
where $N(S)=\set{b\in B:\{a,b\}\in E\text{ for some }a\in S}$ is the
neighborhood of~$S$.
\end{theorem}

\begin{proof}
$(\Rightarrow)$\; If a matching covers $A$, then for every
$S\subseteq A$ the matched partners of $S$ lie in $N(S)$, so
$\abs{N(S)}\ge\abs{S}$.

$(\Leftarrow)$\; By induction on $\abs{A}$.  If $\abs{A}=1$ the
condition $\abs{N(S)}\ge 1$ ensures $A$ has a neighbor, and the single
edge is a matching.

\textbf{Case 1:} $\abs{N(S)}\ge\abs{S}+1$ for every non-empty proper
subset $S\subsetneq A$.  Pick any $a\in A$ and any neighbor $b\in N(a)$.
Match $a$ to~$b$.  In the graph $G'=G-(a,b)$ (removing $a,b$ and their
edges), for any $S'\subseteq A\setminus\{a\}$:
$\abs{N_{G'}(S')}\ge\abs{N_G(S')}-1\ge(\abs{S'}+1)-1=\abs{S'}$.
By induction, $G'$ has a matching covering $A\setminus\{a\}$.

\textbf{Case 2:} There exists a non-empty proper subset $S\subsetneq A$
with $\abs{N(S)}=\abs{S}$.  By induction, the subgraph induced by
$S\cup N(S)$ has a matching $M_1$ covering~$S$.  Now consider the
subgraph $G''$ induced by $(A\setminus S)\cup(B\setminus N(S))$.  For
any $T\subseteq A\setminus S$:
\[
  \abs{N_{G''}(T)}=\abs{N_G(S\cup T)}-\abs{N(S)}
  \ge\abs{S\cup T}-\abs{S}=\abs{T}.
\]
By induction, $G''$ has a matching $M_2$ covering $A\setminus S$.  Then
$M_1\cup M_2$ covers all of~$A$.
\end{proof}

\begin{example}[Matching in a bipartite graph]
\label{ex:matching}

\begin{center}
\begin{tikzpicture}[
  vtx/.style={circle,draw,minimum size=14pt,inner sep=0pt},
  matched/.style={very thick,red},
  scale=1]
  % Part A
  \node[vtx] (a1) at (0,2)   {$a_1$};
  \node[vtx] (a2) at (0,1)   {$a_2$};
  \node[vtx] (a3) at (0,0)   {$a_3$};
  % Part B
  \node[vtx] (b1) at (3,2)   {$b_1$};
  \node[vtx] (b2) at (3,1)   {$b_2$};
  \node[vtx] (b3) at (3,0)   {$b_3$};
  % Edges
  \draw (a1)--(b1); \draw (a1)--(b2);
  \draw (a2)--(b1); \draw (a2)--(b3);
  \draw (a3)--(b2); \draw (a3)--(b3);
  % Matching
  \draw[matched] (a1)--(b2);
  \draw[matched] (a2)--(b1);
  \draw[matched] (a3)--(b3);
  \node at (1.5,-0.8) {Perfect matching in red};
\end{tikzpicture}
\end{center}
\end{example}

\index{matching|)}
\index{bipartite graph|)}
\index{graph coloring|)}

%------------------------------------------------------------------
\section{Exercises}
\label{sec:ex-ch10}

\begin{exercise}
\label{exo:chi-petersen}
\index{Petersen graph}
Show that the chromatic number of the Petersen graph is~$3$.
\end{exercise}

\begin{exercise}
\label{exo:chromatic-cycle}
Compute the chromatic polynomial $P(C_n,k)$ for the cycle $C_n$ using
deletion--contraction.
\end{exercise}

\begin{exercise}
\label{exo:euler-formula-verify}
Verify Euler's formula $V-E+F=2$ for the cube graph $Q_3$ drawn in the
plane.
\end{exercise}

\begin{exercise}
\label{exo:K23-planar}
Show that $K_{2,3}$ is planar by giving an explicit planar embedding.
\end{exercise}

\begin{exercise}
\label{exo:hall-application}
A dormitory has $n$ students, each listing $3$ acceptable rooms from
among $n$ rooms.  If every set of $k$ students collectively list at
least $k$ acceptable rooms, show that every student can be assigned an
acceptable room.
\end{exercise}

\begin{exercise}
\label{exo:konig-verify}
In the bipartite graph of Example~\ref{ex:matching}, find a minimum
vertex cover and verify K\"onig's theorem.
\end{exercise}

\begin{exercise}
\label{exo:bipartite-degree}
Let $G=(A\cup B,E)$ be a bipartite graph where every vertex in $A$ has
degree $d\ge 1$ and every vertex in $B$ has degree at most~$d$.  Prove
that $G$ has a matching covering~$A$.
\end{exercise}


%%%%%%%%%%%%%%%%%%%%%%%%%%%%%%%%%%%%%%%%%%%%%%%%%%%%%%%%%%%%%
%  CHAPTER 11 — Introduction to Coding Theory
%%%%%%%%%%%%%%%%%%%%%%%%%%%%%%%%%%%%%%%%%%%%%%%%%%%%%%%%%%%%%
\chapter{Introduction to Coding Theory}
\label{ch:coding-theory}
\index{coding theory|(}

Digital communication channels are inherently noisy.  Coding theory
provides mathematical tools to detect and correct errors introduced
during transmission or storage.  This chapter introduces the basic
concepts of error-correcting codes, with emphasis on linear codes and
the elegant family of Hamming codes.

%------------------------------------------------------------------
\section{Basic Concepts}
\label{sec:codes-basic}

We work over the binary field $\mathbb{F}_2=\{0,1\}$ with arithmetic
modulo~$2$.  A \emph{binary code} of length~$n$ is a subset
$\mathcal{C}\subseteq\mathbb{F}_2^n$.  Each element of $\mathcal{C}$
is a \emph{codeword}.
\index{codeword}
\index{binary code}

\begin{definition}[Hamming distance]
\label{def:hamming-dist}
\index{Hamming distance}
The \emph{Hamming distance} $d(\mathbf{x},\mathbf{y})$ between two
vectors $\mathbf{x},\mathbf{y}\in\mathbb{F}_2^n$ is the number of
positions in which they differ:
\[
  d(\mathbf{x},\mathbf{y})
  =\card\set{i:x_i\ne y_i}.
\]
\end{definition}

\begin{definition}[Hamming weight]
\label{def:hamming-weight}
\index{Hamming weight}
The \emph{Hamming weight} $\operatorname{wt}(\mathbf{x})$ of a vector
$\mathbf{x}$ is $d(\mathbf{x},\mathbf{0})$, the number of nonzero
entries.
\end{definition}

\begin{proposition}
\label{prop:hamming-metric}
The Hamming distance is a metric on $\mathbb{F}_2^n$.
\end{proposition}

\begin{definition}[Minimum distance]
\label{def:min-dist}
\index{minimum distance}
The \emph{minimum distance} of a code $\mathcal{C}$ is
\[
  d(\mathcal{C})=\min\set{d(\mathbf{x},\mathbf{y})
    :\mathbf{x},\mathbf{y}\in\mathcal{C},\;\mathbf{x}\ne\mathbf{y}}.
\]
An $(n,M,d)$-code has length~$n$, $M$ codewords, and minimum
distance~$d$.
\end{definition}

\begin{theorem}[Error detection and correction]
\label{thm:error-capability}
\index{error detection}
\index{error correction}
A code with minimum distance $d$ can:
\begin{enumerate}
  \item detect up to $d-1$ errors,
  \item correct up to $\lfloor(d-1)/2\rfloor$ errors.
\end{enumerate}
\end{theorem}

\begin{proof}
(1) If at most $d-1$ errors occur, the received word differs from the
sent codeword in at most $d-1$ positions, hence is not a codeword
(since all codewords are at distance $\ge d$).

(2) If at most $t=\lfloor(d-1)/2\rfloor$ errors occur, the received
word $\mathbf{r}$ is within distance $t$ of the sent codeword
$\mathbf{c}$.  For any other codeword $\mathbf{c}'$:
\[
  d(\mathbf{c},\mathbf{c}')\le d(\mathbf{c},\mathbf{r})
  +d(\mathbf{r},\mathbf{c}')
  \implies d(\mathbf{r},\mathbf{c}')\ge d-t>t.
\]
So $\mathbf{c}$ is the unique closest codeword to~$\mathbf{r}$.
\end{proof}

\begin{example}[Repetition code]
\label{ex:repetition}
\index{repetition code}
The binary repetition code of length~$n$ is
$\mathcal{C}=\{00\cdots0,\,11\cdots1\}$ with $d=n$.  It can correct
$\lfloor(n-1)/2\rfloor$ errors but has very low rate ($1/n$).
\end{example}

%------------------------------------------------------------------
\section{Linear Codes}
\label{sec:linear-codes}
\index{linear code|(}

\begin{definition}[Linear code]
\label{def:linear-code}
A \emph{linear code} $\mathcal{C}$ is a $k$-dimensional subspace of
$\mathbb{F}_2^n$.  We say $\mathcal{C}$ is an $[n,k]$-code, or
$[n,k,d]$-code if the minimum distance is~$d$.
\end{definition}

\begin{proposition}
\label{prop:linear-min-weight}
For a linear code, the minimum distance equals the minimum weight of a
nonzero codeword:
\[
  d(\mathcal{C})=\min\set{\operatorname{wt}(\mathbf{c})
    :\mathbf{c}\in\mathcal{C},\;\mathbf{c}\ne\mathbf{0}}.
\]
\end{proposition}

\begin{proof}
Since $\mathcal{C}$ is a subspace, $\mathbf{x}-\mathbf{y}\in\mathcal{C}$
for all $\mathbf{x},\mathbf{y}\in\mathcal{C}$.  Over $\mathbb{F}_2$,
$d(\mathbf{x},\mathbf{y})=\operatorname{wt}(\mathbf{x}+\mathbf{y})$,
so the minimum distance is the minimum weight of a nonzero codeword.
\end{proof}

\subsection{Generator and Parity-Check Matrices}
\label{ssec:generator-parity}
\index{generator matrix}
\index{parity-check matrix}

\begin{definition}[Generator matrix]
\label{def:generator}
A \emph{generator matrix} of an $[n,k]$-code $\mathcal{C}$ is a
$k\times n$ matrix $G$ whose rows form a basis of~$\mathcal{C}$.
Thus $\mathcal{C}=\set{\mathbf{u}G:\mathbf{u}\in\mathbb{F}_2^k}$.
\end{definition}

\begin{definition}[Systematic form]
\label{def:systematic}
\index{systematic form}
A generator matrix is in \emph{systematic form} (or \emph{standard
form}) if $G=[I_k\mid A]$, where $I_k$ is the $k\times k$ identity and
$A$ is a $k\times(n-k)$ matrix.
\end{definition}

\begin{definition}[Parity-check matrix]
\label{def:parity-check}
A \emph{parity-check matrix} of an $[n,k]$-code is an
$(n-k)\times n$ matrix $H$ such that
$\mathcal{C}=\set{\mathbf{x}\in\mathbb{F}_2^n:H\mathbf{x}^T=\mathbf{0}}$.
\end{definition}

\begin{proposition}
\label{prop:GH-relation}
If $G=[I_k\mid A]$ is a generator matrix in systematic form, then
$H=[-A^T\mid I_{n-k}]$ is a parity-check matrix.  Over $\mathbb{F}_2$,
$-A^T=A^T$, so $H=[A^T\mid I_{n-k}]$.
\end{proposition}

\begin{proof}
We verify $GH^T=0$: $[I_k\mid A][A\mid I_{n-k}]^T=A+A=0$
over $\mathbb{F}_2$.  Since $H$ has rank $n-k$, the null space of $H$
is $k$-dimensional, hence equals~$\mathcal{C}$.
\end{proof}

\begin{theorem}[Minimum distance from parity-check matrix]
\label{thm:min-dist-H}
\index{minimum distance!from parity-check matrix}
The minimum distance of a linear code with parity-check matrix $H$
equals the minimum number of linearly dependent columns of~$H$.
\end{theorem}

\begin{proof}
A codeword $\mathbf{c}$ satisfies $H\mathbf{c}^T=\mathbf{0}$, meaning
the columns of $H$ indexed by the support of $\mathbf{c}$ sum to zero.
Thus $\operatorname{wt}(\mathbf{c})$ equals the number of columns in a
linear dependence.  The minimum weight (hence minimum distance) is the
smallest such dependence.
\end{proof}

\subsection{Syndrome Decoding}
\label{ssec:syndrome}
\index{syndrome decoding}

\begin{definition}[Syndrome]
\label{def:syndrome}
\index{syndrome}
The \emph{syndrome} of a received vector $\mathbf{r}\in\mathbb{F}_2^n$
is $\mathbf{s}=H\mathbf{r}^T\in\mathbb{F}_2^{n-k}$.
\end{definition}

\begin{proposition}
\label{prop:syndrome-coset}
Two vectors have the same syndrome if and only if they are in the same
coset of~$\mathcal{C}$ in $\mathbb{F}_2^n$.
\end{proposition}

\begin{proof}
$H\mathbf{r}_1^T=H\mathbf{r}_2^T$ iff
$H(\mathbf{r}_1-\mathbf{r}_2)^T=\mathbf{0}$ iff
$\mathbf{r}_1-\mathbf{r}_2\in\mathcal{C}$.
\end{proof}

The syndrome decoding procedure:
\begin{enumerate}
  \item Precompute a \emph{syndrome table}: for each syndrome
        $\mathbf{s}$, store the \emph{coset leader} (minimum-weight
        vector in the coset).
  \item Given received $\mathbf{r}$, compute
        $\mathbf{s}=H\mathbf{r}^T$.
  \item Look up the coset leader $\mathbf{e}$ for syndrome~$\mathbf{s}$.
  \item Decode as $\hat{\mathbf{c}}=\mathbf{r}-\mathbf{e}
        =\mathbf{r}+\mathbf{e}$ (over $\mathbb{F}_2$).
\end{enumerate}

\begin{example}
\label{ex:syndrome-decode}
Consider the $[6,3,3]$-code with generator matrix
\[
  G=\begin{pmatrix}
    1&0&0&1&1&0\\
    0&1&0&0&1&1\\
    0&0&1&1&0&1
  \end{pmatrix}.
\]
The parity-check matrix is
\[
  H=\begin{pmatrix}
    1&0&1&1&0&0\\
    1&1&0&0&1&0\\
    0&1&1&0&0&1
  \end{pmatrix}.
\]
If $\mathbf{r}=(1,1,0,1,1,1)$ is received, the syndrome is
$H\mathbf{r}^T=(1,0,0)^T$.  The coset leader with this syndrome is
$\mathbf{e}=(0,0,1,0,0,0)$, so we decode
$\hat{\mathbf{c}}=(1,1,1,1,1,1)$.
\end{example}

\index{linear code|)}

%------------------------------------------------------------------
\section{Hamming Codes}
\label{sec:hamming}
\index{Hamming code|(}

\begin{definition}[Hamming code]
\label{def:hamming-code}
For an integer $r\ge 2$, the \emph{Hamming code} $\operatorname{Ham}(r,2)$
is the $[2^r-1,\,2^r-1-r,\,3]$-code whose parity-check matrix $H$ has
as columns all nonzero vectors in $\mathbb{F}_2^r$ (in any order).
\end{definition}

\begin{theorem}
\label{thm:hamming-properties}
The Hamming code $\operatorname{Ham}(r,2)$ has the following properties:
\begin{enumerate}
  \item It is a $[2^r-1,\,2^r-1-r,\,3]$-code.
  \item It can correct any single error.
  \item It is a \emph{perfect code}: every vector in $\mathbb{F}_2^n$
        is within Hamming distance~$1$ of exactly one codeword.
\end{enumerate}
\end{theorem}

\begin{proof}
(1) The matrix $H$ is $r\times(2^r-1)$ with rank~$r$ (it contains
$I_r$ as a submatrix).  So the code has dimension $(2^r-1)-r$.  Any two
columns of $H$ are distinct nonzero vectors, hence linearly independent
(over $\mathbb{F}_2$, two distinct nonzero vectors are independent).
Thus no codeword has weight $1$ or~$2$, and $d\ge 3$.  Since $H$ has
$2^r-1\ge 3$ columns, some three are dependent, so $d=3$.

(2) With $d=3$, the code corrects $\lfloor(3-1)/2\rfloor=1$ error.

(3) The number of vectors within distance~$1$ of a codeword is
$1+n=1+(2^r-1)=2^r$.  The total number of such balls is
$\abs{\mathcal{C}}\cdot 2^r=2^{2^r-1-r}\cdot 2^r=2^{2^r-1}$, which
equals $\abs{\mathbb{F}_2^n}$.  Hence the balls partition
$\mathbb{F}_2^n$: the code is perfect.
\index{perfect code}
\end{proof}

\begin{example}[Hamming(7,4)]
\label{ex:hamming74}
\index{Hamming code!Ham(7,4)@$\operatorname{Ham}(7,4)$}
For $r=3$, we get the $[7,4,3]$-Hamming code.  A standard parity-check
matrix is
\[
  H=\begin{pmatrix}
    0&0&0&1&1&1&1\\
    0&1&1&0&0&1&1\\
    1&0&1&0&1&0&1
  \end{pmatrix},
\]
where column $j$ is the binary representation of~$j$ (for $j=1,\dots,7$).
The syndrome $H\mathbf{r}^T$ of a single-error vector directly gives
the binary index of the error position.
\end{example}

\begin{remark}[Decoding Hamming codes]
\label{rem:hamming-decode}
If the syndrome is $\mathbf{0}$, no error is detected.  Otherwise, the
syndrome (read as a binary number) gives the position of the single
error, which is then flipped.  This makes Hamming code decoding
exceptionally simple.
\end{remark}

\begin{example}[Hamming(7,4) decoding]
\label{ex:hamming-decode}
Suppose we send codeword $\mathbf{c}=(1,0,1,1,0,0,1)$ and receive
$\mathbf{r}=(1,0,1,1,0,1,1)$ (error in position~$6$).  The syndrome is
\[
  H\mathbf{r}^T=\begin{pmatrix}1\\1\\0\end{pmatrix}
  =(110)_2=6.
\]
We flip bit~$6$ to recover $\hat{\mathbf{c}}=(1,0,1,1,0,0,1)=\mathbf{c}$.
\end{example}

\index{Hamming code|)}

%------------------------------------------------------------------
\section{The Singleton and Hamming Bounds}
\label{sec:bounds}
\index{Singleton bound}
\index{Hamming bound}

\begin{theorem}[Singleton bound]
\label{thm:singleton}
An $(n,M,d)$-code satisfies $M\le 2^{n-d+1}$.  Equivalently, for an
$[n,k,d]$-linear code, $k\le n-d+1$.
\end{theorem}

\begin{proof}
Project each codeword onto the first $n-d+1$ coordinates.  Any two
distinct codewords that agree on these coordinates must differ only
in the last $d-1$ coordinates, giving distance at most $d-1<d$, a
contradiction.  Hence the projection is injective:
$M\le 2^{n-d+1}$.
\end{proof}

\begin{theorem}[Hamming bound (sphere-packing bound)]
\label{thm:hamming-bound}
An $(n,M,d)$-code with $d=2t+1$ satisfies
\[
  M\cdot\sum_{i=0}^{t}\binom{n}{i}\le 2^n.
\]
Codes achieving equality are \emph{perfect}.
\end{theorem}

\begin{proof}
The balls of radius $t$ centered at distinct codewords are disjoint
(since any two codewords are at distance $\ge 2t+1$).  Each ball
contains $\sum_{i=0}^t\binom{n}{i}$ vectors.  All balls lie within
$\mathbb{F}_2^n$, which has $2^n$ elements.
\end{proof}

%------------------------------------------------------------------
\section{Reed--Solomon Codes: A Brief Mention}
\label{sec:reed-solomon}
\index{Reed--Solomon code}

Reed--Solomon codes, introduced by Irving Reed and Gustave Solomon in
1960, are among the most widely deployed error-correcting codes.  They
operate over larger finite fields $\mathbb{F}_q$ (typically
$q=2^m$) and achieve the Singleton bound with equality, making them
\emph{maximum distance separable} (MDS) codes.
\index{MDS code}

A Reed--Solomon code of length $n=q-1$ and dimension $k$ encodes a
message polynomial $m(x)$ of degree $<k$ as the vector of evaluations
$(m(\alpha_1),\dots,m(\alpha_n))$, where $\alpha_1,\dots,\alpha_n$ are
the nonzero elements of $\mathbb{F}_q$.  The minimum distance is
$d=n-k+1$.

\begin{remark}
Reed--Solomon codes are used in compact discs, QR~codes, deep-space
communication (Voyager, Mars rovers), and digital television (DVB).
Their algebraic structure enables efficient decoding algorithms such as
the Berlekamp--Massey algorithm.
\index{Berlekamp--Massey algorithm}
\end{remark}

%------------------------------------------------------------------
\section{Exercises}
\label{sec:ex-ch11}

\begin{exercise}
\label{exo:hamming-verify}
Verify that the Hamming distance is a metric (prove the triangle
inequality).
\end{exercise}

\begin{exercise}
\label{exo:parity-check}
The $[n,n-1,2]$ single-parity-check code consists of all binary vectors
of even weight.  Find its generator and parity-check matrices.
\end{exercise}

\begin{exercise}
\label{exo:hamming-15}
Write down the parity-check matrix for $\operatorname{Ham}(4,2)$, a
$[15,11,3]$-code.  How many codewords does it have?
\end{exercise}

\begin{exercise}
\label{exo:syndrome-table}
For the $[7,4,3]$-Hamming code, construct the complete syndrome table
and decode the received vectors $\mathbf{r}_1=(1,1,0,0,1,0,1)$ and
$\mathbf{r}_2=(0,1,1,1,0,1,0)$.
\end{exercise}

\begin{exercise}
\label{exo:perfect-codes}
Show that the only binary perfect $e$-error-correcting codes (for
$e\ge 1$) are the Hamming codes ($e=1$), the binary Golay code
($e=3$, $n=23$), and the repetition codes of odd length.
(You may quote the classification theorem without proof.)
\end{exercise}

\begin{exercise}
\label{exo:singleton-mds}
Show that the $[n,1,n]$ repetition code and the $[n,n-1,2]$ parity-check
code both meet the Singleton bound.
\end{exercise}

\begin{exercise}
\label{exo:dual-code}
\index{dual code}
The \emph{dual code} $\mathcal{C}^\perp$ of an $[n,k]$-code
$\mathcal{C}$ is the set of all vectors orthogonal to every codeword.
Show that $\mathcal{C}^\perp$ is an $[n,n-k]$-code and that
$(\mathcal{C}^\perp)^\perp=\mathcal{C}$.
\end{exercise}

\begin{exercise}
\label{exo:weight-enumerator}
\index{weight enumerator}
The \emph{weight enumerator} of a code $\mathcal{C}$ is the polynomial
$W_{\mathcal{C}}(x,y)=\sum_{\mathbf{c}\in\mathcal{C}}
x^{n-\operatorname{wt}(\mathbf{c})}y^{\operatorname{wt}(\mathbf{c})}$.
Compute $W_{\mathcal{C}}(x,y)$ for the $[7,4,3]$-Hamming code.
\end{exercise}

\index{coding theory|)}


%%%%%%%%%%%%%%%%%%%%%%%%%%%%%%%%%%%%%%%%%%%%%%%%%%%%%%%%%%%%%
%  APPENDIX — Combinatorial Identities
%%%%%%%%%%%%%%%%%%%%%%%%%%%%%%%%%%%%%%%%%%%%%%%%%%%%%%%%%%%%%
\appendix
\chapter{Combinatorial Identities}
\label{app:comb-identities}
\index{combinatorial identities}

The following table collects important combinatorial identities that
appear throughout this course.  In each identity, $n$ and $k$ are
non-negative integers with $0\le k\le n$ unless otherwise stated, and
$x,y$ may be real or complex numbers as appropriate.

\bigskip

\renewcommand{\arraystretch}{1.6}
\begin{longtable}{p{5cm}p{8cm}}
\hline
\textbf{Identity} & \textbf{Name / Comment} \\
\hline\hline
\endfirsthead
\hline
\textbf{Identity} & \textbf{Name / Comment} \\
\hline\hline
\endhead
\hline
\endfoot

$\displaystyle\binom{n}{k}=\binom{n}{n-k}$
  & Symmetry of binomial coefficients
  \index{binomial coefficient!symmetry}\\

$\displaystyle\binom{n}{k}=\binom{n-1}{k-1}+\binom{n-1}{k}$
  & Pascal's identity
  \index{Pascal's identity}\\

$\displaystyle\sum_{k=0}^{n}\binom{n}{k}=2^n$
  & Sum of binomial coefficients\\

$\displaystyle\sum_{k=0}^{n}(-1)^k\binom{n}{k}=0$
  & Alternating sum ($n\ge 1$)\\

$\displaystyle(x+y)^n=\sum_{k=0}^{n}\binom{n}{k}x^{k}y^{n-k}$
  & Binomial theorem
  \index{binomial theorem}\\

$\displaystyle\sum_{k=0}^{r}\binom{m}{k}\binom{n}{r-k}
  =\binom{m+n}{r}$
  & Vandermonde's identity
  \index{Vandermonde's identity}\\

$\displaystyle\sum_{k=0}^{n}\binom{n}{k}^2=\binom{2n}{n}$
  & Special case of Vandermonde ($m=n=r$)\\

$\displaystyle\binom{n}{k}\binom{k}{j}
  =\binom{n}{j}\binom{n-j}{k-j}$
  & Absorption / triple binomial\\

$\displaystyle k\binom{n}{k}=n\binom{n-1}{k-1}$
  & Absorption identity
  \index{absorption identity}\\

$\displaystyle\sum_{k=0}^{n}k\binom{n}{k}=n\,2^{n-1}$
  & Weighted sum\\

$\displaystyle\sum_{j=k}^{n}\binom{j}{k}=\binom{n+1}{k+1}$
  & Hockey-stick / Christmas-stocking identity
  \index{hockey-stick identity}\\

$\displaystyle\binom{-n}{k}=(-1)^k\binom{n+k-1}{k}$
  & Upper negation\\

$\displaystyle\binom{n+k-1}{k}=\binom{n+k-1}{n-1}$
  & Stars and bars (multiset coefficient)
  \index{stars and bars}\\

$\displaystyle\sum_{k=0}^{n}k^2\binom{n}{k}
  =n(n+1)2^{n-2}$  \newline (for $n\ge 2$)
  & Second moment\\

$\displaystyle\sum_{k=0}^{n}\binom{n}{k}x^k=(1+x)^n$
  & Binomial series (finite)\\

$\displaystyle\sum_{k=0}^{\infty}\binom{n+k}{k}x^k
  =\frac{1}{(1-x)^{n+1}}$
  & Negative binomial series ($\abs{x}<1$)
  \index{negative binomial series}\\

$\displaystyle D_n=n!\sum_{k=0}^{n}\frac{(-1)^k}{k!}$
  & Number of derangements
  \index{derangement}\\

$\displaystyle\abs{A_1\cup\cdots\cup A_n}
  =\sum_{k=1}^{n}(-1)^{k+1}
  \!\!\sum_{I\subseteq[n],\,\abs{I}=k}\!\!\abs{\bigcap_{i\in I}A_i}$
  & Inclusion--exclusion principle
  \index{inclusion-exclusion}\\

$\displaystyle\sum_{d\mid n}\varphi(d)=n$
  & Euler totient sum
  \index{Euler totient function}\\

$\displaystyle S(n,k)=kS(n-1,k)+S(n-1,k-1)$
  & Stirling numbers of the second kind recurrence
  \index{Stirling number!second kind}\\

$\displaystyle p(n)=\sum_{k\ge 1}(-1)^{k+1}
  \bigl(p(n-\omega_k)+p(n-\bar\omega_k)\bigr)$
  & Euler's pentagonal recurrence for partitions
  \index{partition!pentagonal recurrence}\\

\end{longtable}

\noindent
Here $\omega_k=k(3k-1)/2$ and $\bar\omega_k=k(3k+1)/2$ are
generalized pentagonal numbers, $\varphi$ is Euler's totient function,
$S(n,k)$ denotes a Stirling number of the second kind, and $p(n)$ is
the number of integer partitions of~$n$.

%%%%%%%%%%%%%%%%%%%%%%%%%%%%%%%%%%%%%%%%%%%%%%%%%%%%%%%%%%%%%
%  BIBLIOGRAPHY
%%%%%%%%%%%%%%%%%%%%%%%%%%%%%%%%%%%%%%%%%%%%%%%%%%%%%%%%%%%%%
\begin{thebibliography}{99}

\bibitem{rosen}
\index{Rosen, K.}
K.\,H.\ Rosen,
\emph{Discrete Mathematics and Its Applications},
8th~ed., McGraw-Hill, 2019.

\bibitem{gkp}
\index{Graham, R.}
\index{Knuth, D.}
\index{Patashnik, O.}
R.\,L.\ Graham, D.\,E.\ Knuth, and O.\ Patashnik,
\emph{Concrete Mathematics: A Foundation for Computer Science},
2nd~ed., Addison-Wesley, 1994.

\bibitem{lovasz}
\index{Lovasz, L.@Lov\'asz, L.}
L.\ Lov\'asz,
\emph{Combinatorial Problems and Exercises},
2nd~ed., AMS Chelsea Publishing, 2007.

\bibitem{diestel}
\index{Diestel, R.}
R.\ Diestel,
\emph{Graph Theory},
5th~ed., Graduate Texts in Mathematics, vol.~173,
Springer, 2017.

\bibitem{cameron}
\index{Cameron, P.}
P.\,J.\ Cameron,
\emph{Combinatorics: Topics, Techniques, Algorithms},
Cambridge University Press, 1994.

\bibitem{vanlint-wilson}
\index{van Lint, J.}
\index{Wilson, R.}
J.\,H.\ van Lint and R.\,M.\ Wilson,
\emph{A Course in Combinatorics},
2nd~ed., Cambridge University Press, 2001.

\end{thebibliography}

%%%%%%%%%%%%%%%%%%%%%%%%%%%%%%%%%%%%%%%%%%%%%%%%%%%%%%%%%%%%%
%  INDEX & END
%%%%%%%%%%%%%%%%%%%%%%%%%%%%%%%%%%%%%%%%%%%%%%%%%%%%%%%%%%%%%
\printindex
\end{document}
